\documentclass[11pt,a4paper]{article}
\usepackage[a4paper,margin=20mm]{geometry}
\usepackage{fontspec}
\setmainfont{Noto Serif}
\setsansfont{Noto Sans}
\usepackage{amsmath,amssymb}
\usepackage{booktabs,tabularx,array}
\usepackage[table]{xcolor}
\usepackage{hyperref}
\definecolor{primary}{RGB}{27,79,114}
\definecolor{accent}{RGB}{40,116,166}
\definecolor{rowtint}{RGB}{235,245,251}
\hypersetup{colorlinks=true,linkcolor=primary,urlcolor=primary}
\setlength{\parskip}{5pt plus 2pt minus 1pt}
\setlength{\parindent}{0pt}
\begin{document}
\begin{center}
  {\LARGE\bfseries\color{primary} PHIẾU HỌC TẬP TRẮC NGHIỆM TOÁN 7}\\[4pt]
  {\normalsize\itshape\color{accent} Bộ sách: Kết nối tri thức với cuộc sống (Tập 1 \& Tập 2 – 10 Chương)}\\[4pt]
  {\small Cấu trúc mỗi chương: 20 câu Dễ (Câu 1--20) $\bullet$ 20 câu Trung bình (Câu 21--40) $\bullet$ 10 câu Khó (Câu 41--50) -- Tổng 50 câu/chương}
\end{center}
\vspace{6pt}\hrule\vspace{10pt}
\section*{\color{primary} CHƯƠNG I: SỐ HỮU TỈ}
\textit{Trọng tâm kiến thức SGK: Tập hợp các số hữu tỉ $\mathbb{Q}$, cộng, trừ, nhân, chia số hữu tỉ, lũy thừa với số mũ tự nhiên của một số hữu tỉ và thứ tự thực hiện các phép tính.}\par\vspace{4pt}
\subsection*{\color{accent} Phần I. Mức độ Dễ (Nhận biết -- Câu 1 đến Câu 20)}
\textbf{\color{primary} Câu 1.} Số đối của số hữu tỉ -3/7 là:\par
\hspace*{1.5em}\textbf{A.} 3/7\quad\quad \textbf{B.} 7/3\quad\quad \textbf{C.} -7/3\quad\quad \textbf{D.} -3/7\par\vspace{2pt}
\textbf{\color{primary} Câu 2.} Thực hiện yêu cầu nhận biết sau: Giá trị của lũy thừa (-1/3)$^2$ bằng.\par
\hspace*{1.5em}\textbf{A.} 1/9\quad\quad \textbf{B.} -1/9\quad\quad \textbf{C.} -2/3\quad\quad \textbf{D.} 2/3\par\vspace{2pt}
\textbf{\color{primary} Câu 3.} Trong bài học thuộc Chương I, xét nội dung: Kết quả của phép tính (2/3)$^{4}$ $\cdot$ (2/3)$^{3}$ bằng. Đáp án đúng là:\par
\hspace*{1.5em}\textbf{A.} (2/3)$^{12}$\quad\quad \textbf{B.} (2/3)$^{7}$\quad\quad \textbf{C.} (2/3)$^{6}$\quad\quad \textbf{D.} (4/9)$^{7}$\par\vspace{2pt}
\textbf{\color{primary} Câu 4.} Dựa vào kiến thức chuẩn SGK của Chương I (Số hữu tỉ), số đối của số hữu tỉ -6/13 là:\par
\hspace*{1.5em}\textbf{A.} 6/13\quad\quad \textbf{B.} -6/13\quad\quad \textbf{C.} 13/6\quad\quad \textbf{D.} -13/6\par\vspace{2pt}
\textbf{\color{primary} Câu 5.} Chọn kết quả đúng nhất khi xét: Giá trị của lũy thừa (-1/6)$^2$ bằng:\par
\hspace*{1.5em}\textbf{A.} -1/36\quad\quad \textbf{B.} -2/6\quad\quad \textbf{C.} 1/36\quad\quad \textbf{D.} 2/6\par\vspace{2pt}
\textbf{\color{primary} Câu 6.} Khi làm bài tập cơ bản về số hữu tỉ, một học sinh gặp câu hỏi: Kết quả của phép tính (2/3)$^{7}$ $\cdot$ (2/3)$^{3}$ bằng. Kết quả nào dưới đây đúng?\par
\hspace*{1.5em}\textbf{A.} (4/9)$^{10}$\quad\quad \textbf{B.} (2/3)$^{9}$\quad\quad \textbf{C.} (2/3)$^{21}$\quad\quad \textbf{D.} (2/3)$^{10}$\par\vspace{2pt}
\textbf{\color{primary} Câu 7.} Xác định giá trị hoặc khẳng định đúng cho bài toán sau: Số đối của số hữu tỉ -9/19 là:\par
\hspace*{1.5em}\textbf{A.} -19/9\quad\quad \textbf{B.} -9/19\quad\quad \textbf{C.} 19/9\quad\quad \textbf{D.} 9/19\par\vspace{2pt}
\textbf{\color{primary} Câu 8.} Áp dụng trực tiếp định nghĩa và quy tắc trong Chương I: Giá trị của lũy thừa (-1/9)$^2$ bằng:\par
\hspace*{1.5em}\textbf{A.} 1/81\quad\quad \textbf{B.} -1/81\quad\quad \textbf{C.} -2/9\quad\quad \textbf{D.} 2/9\par\vspace{2pt}
\textbf{\color{primary} Câu 9.} Kiểm tra nhanh kiến thức nền tảng: Kết quả của phép tính (2/3)$^{10}$ $\cdot$ (2/3)$^{3}$ bằng:\par
\hspace*{1.5em}\textbf{A.} (2/3)$^{30}$\quad\quad \textbf{B.} (2/3)$^{12}$\quad\quad \textbf{C.} (4/9)$^{13}$\quad\quad \textbf{D.} (2/3)$^{13}$\par\vspace{2pt}
\textbf{\color{primary} Câu 10.} Cho dữ kiện nhận biết sau đây: Số đối của số hữu tỉ -12/25 là. Phương án chính xác là:\par
\hspace*{1.5em}\textbf{A.} -12/25\quad\quad \textbf{B.} 12/25\quad\quad \textbf{C.} -25/12\quad\quad \textbf{D.} 25/12\par\vspace{2pt}
\textbf{\color{primary} Câu 11.} Hãy chỉ ra kết quả đúng trong bốn lựa chọn A, B, C, D đối với yêu cầu: Giá trị của lũy thừa (-1/12)$^2$ bằng:\par
\hspace*{1.5em}\textbf{A.} -1/144\quad\quad \textbf{B.} 1/144\quad\quad \textbf{C.} 2/12\quad\quad \textbf{D.} -2/12\par\vspace{2pt}
\textbf{\color{primary} Câu 12.} Câu hỏi ôn tập lý thuyết và kỹ năng cơ bản (Chương I): Kết quả của phép tính (2/3)$^{13}$ $\cdot$ (2/3)$^{3}$ bằng:\par
\hspace*{1.5em}\textbf{A.} (2/3)$^{39}$\quad\quad \textbf{B.} (2/3)$^{16}$\quad\quad \textbf{C.} (2/3)$^{15}$\quad\quad \textbf{D.} (4/9)$^{16}$\par\vspace{2pt}
\textbf{\color{primary} Câu 13.} Theo nội dung sách giáo khoa Kết nối tri thức với cuộc sống (Chương I), số đối của số hữu tỉ -15/31 là:\par
\hspace*{1.5em}\textbf{A.} -31/15\quad\quad \textbf{B.} -15/31\quad\quad \textbf{C.} 31/15\quad\quad \textbf{D.} 15/31\par\vspace{2pt}
\textbf{\color{primary} Câu 14.} Tìm đáp số đúng của bài toán cơ bản sau: Giá trị của lũy thừa (-1/15)$^2$ bằng:\par
\hspace*{1.5em}\textbf{A.} 2/15\quad\quad \textbf{B.} -1/225\quad\quad \textbf{C.} 1/225\quad\quad \textbf{D.} -2/15\par\vspace{2pt}
\textbf{\color{primary} Câu 15.} Xét đối tượng toán học trong Chương I: Kết quả của phép tính (2/3)$^{16}$ $\cdot$ (2/3)$^{3}$ bằng. Giá trị hoặc tính chất tương ứng là:\par
\hspace*{1.5em}\textbf{A.} (4/9)$^{19}$\quad\quad \textbf{B.} (2/3)$^{18}$\quad\quad \textbf{C.} (2/3)$^{19}$\quad\quad \textbf{D.} (2/3)$^{48}$\par\vspace{2pt}
\textbf{\color{primary} Câu 16.} Hoàn thành phép xác định cơ bản sau đây: Số đối của số hữu tỉ -18/37 là:\par
\hspace*{1.5em}\textbf{A.} -37/18\quad\quad \textbf{B.} -18/37\quad\quad \textbf{C.} 18/37\quad\quad \textbf{D.} 37/18\par\vspace{2pt}
\textbf{\color{primary} Câu 17.} Trong các khẳng định hoặc giá trị dưới đây, đâu là đáp án đúng cho câu hỏi: Giá trị của lũy thừa (-1/18)$^2$ bằng?\par
\hspace*{1.5em}\textbf{A.} 2/18\quad\quad \textbf{B.} -2/18\quad\quad \textbf{C.} -1/324\quad\quad \textbf{D.} 1/324\par\vspace{2pt}
\textbf{\color{primary} Câu 18.} Bài tập khởi động Chương I: Kết quả của phép tính (2/3)$^{19}$ $\cdot$ (2/3)$^{3}$ bằng:\par
\hspace*{1.5em}\textbf{A.} (2/3)$^{57}$\quad\quad \textbf{B.} (2/3)$^{22}$\quad\quad \textbf{C.} (4/9)$^{22}$\quad\quad \textbf{D.} (2/3)$^{21}$\par\vspace{2pt}
\textbf{\color{primary} Câu 19.} Quan sát dữ kiện đề bài: Số đối của số hữu tỉ -21/43 là. Hãy chọn phương án đúng:\par
\hspace*{1.5em}\textbf{A.} 21/43\quad\quad \textbf{B.} -21/43\quad\quad \textbf{C.} 43/21\quad\quad \textbf{D.} -43/21\par\vspace{2pt}
\textbf{\color{primary} Câu 20.} Sử dụng công thức hoặc tính chất đầu bài của Chương I để giải quyết yêu cầu: Giá trị của lũy thừa (-1/21)$^2$ bằng:\par
\hspace*{1.5em}\textbf{A.} -1/441\quad\quad \textbf{B.} -2/21\quad\quad \textbf{C.} 1/441\quad\quad \textbf{D.} 2/21\par\vspace{2pt}
\subsection*{\color{accent} Phần II. Mức độ Trung bình (Thông hiểu -- Câu 21 đến Câu 40)}
\textbf{\color{primary} Câu 21.} Tìm số hữu tỉ x, biết: x + 1/2 = 3/2.:\par
\hspace*{1.5em}\textbf{A.} x = 3/2\quad\quad \textbf{B.} x = -2/2\quad\quad \textbf{C.} x = 2/2\quad\quad \textbf{D.} x = 4/2\par\vspace{2pt}
\textbf{\color{primary} Câu 22.} Vận dụng quy tắc trong Chương I (Số hữu tỉ) để giải bài toán: Tính nhanh giá trị biểu thức A = (3/6) $\cdot$ (7/11) + (3/6) $\cdot$ (4/11).:\par
\hspace*{1.5em}\textbf{A.} 1/6\quad\quad \textbf{B.} 33/67\quad\quad \textbf{C.} 6/6\quad\quad \textbf{D.} 3/6\par\vspace{2pt}
\textbf{\color{primary} Câu 23.} Biến đổi và tính toán theo kiến thức Chương I: Tìm số tự nhiên n biết (1/2)$^{n}$ = 1/32.. Kết quả thu được là:\par
\hspace*{1.5em}\textbf{A.} n = 4\quad\quad \textbf{B.} n = 16\quad\quad \textbf{C.} n = 5\quad\quad \textbf{D.} n = 6\par\vspace{2pt}
\textbf{\color{primary} Câu 24.} Xét bài toán mức độ thông hiểu sau: Tìm số hữu tỉ x, biết: x + 1/5 = 3/5.. Đáp án đúng là:\par
\hspace*{1.5em}\textbf{A.} x = 3/5\quad\quad \textbf{B.} x = 4/5\quad\quad \textbf{C.} x = -2/5\quad\quad \textbf{D.} x = 2/5\par\vspace{2pt}
\textbf{\color{primary} Câu 25.} Để hoàn thành phiếu học tập Chương I, hãy xác định: Tính nhanh giá trị biểu thức A = (3/9) $\cdot$ (7/11) + (3/9) $\cdot$ (4/11).:\par
\hspace*{1.5em}\textbf{A.} 33/100\quad\quad \textbf{B.} 6/9\quad\quad \textbf{C.} 3/9\quad\quad \textbf{D.} 1/9\par\vspace{2pt}
\textbf{\color{primary} Câu 26.} Phân tích dữ kiện bài toán: Tìm số tự nhiên n biết (1/2)$^{n}$ = 1/8.. Giá trị cần tìm bằng:\par
\hspace*{1.5em}\textbf{A.} n = 3\quad\quad \textbf{B.} n = 4\quad\quad \textbf{C.} n = 2\quad\quad \textbf{D.} n = 3 (1)\par\vspace{2pt}
\textbf{\color{primary} Câu 27.} Áp dụng tính chất đã học ở Chương I vào bài tập sau: Tìm số hữu tỉ x, biết: x + 1/8 = 3/8.:\par
\hspace*{1.5em}\textbf{A.} x = 2/8\quad\quad \textbf{B.} x = -2/8\quad\quad \textbf{C.} x = 4/8\quad\quad \textbf{D.} x = 3/8\par\vspace{2pt}
\textbf{\color{primary} Câu 28.} Thực hiện các bước suy luận hợp lý để giải quyết yêu cầu: Tính nhanh giá trị biểu thức A = (3/12) $\cdot$ (7/11) + (3/12) $\cdot$ (4/11).:\par
\hspace*{1.5em}\textbf{A.} 6/12\quad\quad \textbf{B.} 1/12\quad\quad \textbf{C.} 33/133\quad\quad \textbf{D.} 3/12\par\vspace{2pt}
\textbf{\color{primary} Câu 29.} Cho bài toán tự luyện (Số hữu tỉ): Tìm số tự nhiên n biết (1/2)$^{n}$ = 1/64.. Phương án nào sau đây chính xác?\par
\hspace*{1.5em}\textbf{A.} n = 6\quad\quad \textbf{B.} n = 32\quad\quad \textbf{C.} n = 5\quad\quad \textbf{D.} n = 7\par\vspace{2pt}
\textbf{\color{primary} Câu 30.} Kết hợp các phép biến đổi trong Chương I, hãy tìm kết quả của bài toán: Tìm số hữu tỉ x, biết: x + 1/11 = 3/11.:\par
\hspace*{1.5em}\textbf{A.} x = 2/11\quad\quad \textbf{B.} x = 3/11\quad\quad \textbf{C.} x = -2/11\quad\quad \textbf{D.} x = 4/11\par\vspace{2pt}
\textbf{\color{primary} Câu 31.} Một bài kiểm tra định kỳ nội dung Chương I nêu câu hỏi: Tính nhanh giá trị biểu thức A = (3/15) $\cdot$ (7/11) + (3/15) $\cdot$ (4/11).. Hãy chọn đáp số đúng:\par
\hspace*{1.5em}\textbf{A.} 3/15\quad\quad \textbf{B.} 6/15\quad\quad \textbf{C.} 33/166\quad\quad \textbf{D.} 1/15\par\vspace{2pt}
\textbf{\color{primary} Câu 32.} Tìm lời giải đúng cho bài toán thông hiểu: Tìm số tự nhiên n biết (1/2)$^{n}$ = 1/16.:\par
\hspace*{1.5em}\textbf{A.} n = 8\quad\quad \textbf{B.} n = 5\quad\quad \textbf{C.} n = 3\quad\quad \textbf{D.} n = 4\par\vspace{2pt}
\textbf{\color{primary} Câu 33.} Sử dụng mối liên hệ giữa các đại lượng trong Chương I: Tìm số hữu tỉ x, biết: x + 1/14 = 3/14.:\par
\hspace*{1.5em}\textbf{A.} x = 2/14\quad\quad \textbf{B.} x = 4/14\quad\quad \textbf{C.} x = 3/14\quad\quad \textbf{D.} x = -2/14\par\vspace{2pt}
\textbf{\color{primary} Câu 34.} Hãy tính toán và chọn phương án đúng đối với bài tập: Tính nhanh giá trị biểu thức A = (3/18) $\cdot$ (7/11) + (3/18) $\cdot$ (4/11).:\par
\hspace*{1.5em}\textbf{A.} 33/199\quad\quad \textbf{B.} 1/18\quad\quad \textbf{C.} 3/18\quad\quad \textbf{D.} 6/18\par\vspace{2pt}
\textbf{\color{primary} Câu 35.} Dựa trên các tính chất trọng tâm của Chương I, xét bài toán: Tìm số tự nhiên n biết (1/2)$^{n}$ = 1/128.:\par
\hspace*{1.5em}\textbf{A.} n = 64\quad\quad \textbf{B.} n = 8\quad\quad \textbf{C.} n = 6\quad\quad \textbf{D.} n = 7\par\vspace{2pt}
\textbf{\color{primary} Câu 36.} Thực hành giải toán (Số hữu tỉ): Tìm số hữu tỉ x, biết: x + 1/17 = 3/17.. Đáp án nào dưới đây thỏa mãn?\par
\hspace*{1.5em}\textbf{A.} x = 3/17\quad\quad \textbf{B.} x = 2/17\quad\quad \textbf{C.} x = 4/17\quad\quad \textbf{D.} x = -2/17\par\vspace{2pt}
\textbf{\color{primary} Câu 37.} Kiểm tra kỹ năng biến đổi và suy luận qua bài toán: Tính nhanh giá trị biểu thức A = (3/21) $\cdot$ (7/11) + (3/21) $\cdot$ (4/11).:\par
\hspace*{1.5em}\textbf{A.} 6/21\quad\quad \textbf{B.} 3/21\quad\quad \textbf{C.} 33/232\quad\quad \textbf{D.} 1/21\par\vspace{2pt}
\textbf{\color{primary} Câu 38.} Cho giả thiết thuộc phạm vi Chương I: Tìm số tự nhiên n biết (1/2)$^{n}$ = 1/32.. Kết luận đúng là:\par
\hspace*{1.5em}\textbf{A.} n = 6\quad\quad \textbf{B.} n = 16\quad\quad \textbf{C.} n = 4\quad\quad \textbf{D.} n = 5\par\vspace{2pt}
\textbf{\color{primary} Câu 39.} Rút gọn hoặc tính giá trị theo yêu cầu sau của Chương I: Tìm số hữu tỉ x, biết: x + 1/20 = 3/20.:\par
\hspace*{1.5em}\textbf{A.} x = 3/20\quad\quad \textbf{B.} x = 4/20\quad\quad \textbf{C.} x = 2/20\quad\quad \textbf{D.} x = -2/20\par\vspace{2pt}
\textbf{\color{primary} Câu 40.} Xác định phương án đúng sau khi thực hiện phép tính / suy luận: Tính nhanh giá trị biểu thức A = (3/24) $\cdot$ (7/11) + (3/24) $\cdot$ (4/11).:\par
\hspace*{1.5em}\textbf{A.} 1/24\quad\quad \textbf{B.} 3/24\quad\quad \textbf{C.} 33/265\quad\quad \textbf{D.} 6/24\par\vspace{2pt}
\subsection*{\color{accent} Phần III. Mức độ Khó (Vận dụng \& Vận dụng cao -- Câu 41 đến Câu 50)}
\textbf{\color{primary} Câu 41.} Tìm giá trị nhỏ nhất của biểu thức P = (x - 1/2)$^2$ + 5.:\par
\hspace*{1.5em}\textbf{A.} 0\quad\quad \textbf{B.} 1/2\quad\quad \textbf{C.} -5\quad\quad \textbf{D.} 5\par\vspace{2pt}
\textbf{\color{primary} Câu 42.} Bài toán vận dụng (Chương I – Số hữu tỉ): So sánh hai lũy thừa A = 2$^{15}$ và B = 3$^{10}$.. Kết quả đúng là:\par
\hspace*{1.5em}\textbf{A.} A < B (vì 8 < 9)\quad\quad \textbf{B.} A > B\quad\quad \textbf{C.} A = B\quad\quad \textbf{D.} Không so sánh được\par\vspace{2pt}
\textbf{\color{primary} Câu 43.} Vận dụng kiến thức nâng cao trong phạm vi Chương I để giải bài toán: Tìm giá trị nhỏ nhất của biểu thức P = (x - 3/2)$^2$ + 9.:\par
\hspace*{1.5em}\textbf{A.} 3/2\quad\quad \textbf{B.} -9\quad\quad \textbf{C.} 9\quad\quad \textbf{D.} 0\par\vspace{2pt}
\textbf{\color{primary} Câu 44.} Phân tích và giải bài toán tổng hợp sau của Chương I: So sánh hai lũy thừa A = 2$^{21}$ và B = 3$^{14}$.:\par
\hspace*{1.5em}\textbf{A.} A = B\quad\quad \textbf{B.} A < B (vì 8 < 9)\quad\quad \textbf{C.} Không so sánh được\quad\quad \textbf{D.} A > B\par\vspace{2pt}
\textbf{\color{primary} Câu 45.} Thử thách tư duy trong phạm vi kiến thức Chương I: Tìm giá trị nhỏ nhất của biểu thức P = (x - 5/2)$^2$ + 13.. Đáp số chính xác là:\par
\hspace*{1.5em}\textbf{A.} 13\quad\quad \textbf{B.} -13\quad\quad \textbf{C.} 0\quad\quad \textbf{D.} 5/2\par\vspace{2pt}
\textbf{\color{primary} Câu 46.} Sử dụng phương pháp giải toán của Chương I để tìm đáp án cho bài toán: So sánh hai lũy thừa A = 2$^{27}$ và B = 3$^{18}$.:\par
\hspace*{1.5em}\textbf{A.} Không so sánh được\quad\quad \textbf{B.} A < B (vì 8 < 9)\quad\quad \textbf{C.} A > B\quad\quad \textbf{D.} A = B\par\vspace{2pt}
\textbf{\color{primary} Câu 47.} Xét bài toán vận dụng thực tiễn / nâng cao thuộc Chương I: Tìm giá trị nhỏ nhất của biểu thức P = (x - 7/2)$^2$ + 17.:\par
\hspace*{1.5em}\textbf{A.} 17\quad\quad \textbf{B.} -17\quad\quad \textbf{C.} 7/2\quad\quad \textbf{D.} 0\par\vspace{2pt}
\textbf{\color{primary} Câu 48.} Tìm kết quả cuối cùng của bài toán phân loại học sinh khá giỏi (Chương I): So sánh hai lũy thừa A = 2$^{33}$ và B = 3$^{22}$.:\par
\hspace*{1.5em}\textbf{A.} A = B\quad\quad \textbf{B.} Không so sánh được\quad\quad \textbf{C.} A > B\quad\quad \textbf{D.} A < B (vì 8 < 9)\par\vspace{2pt}
\textbf{\color{primary} Câu 49.} Kết hợp các tính chất đã học trong Chương I, hãy giải quyết bài toán: Tìm giá trị nhỏ nhất của biểu thức P = (x - 9/2)$^2$ + 21.:\par
\hspace*{1.5em}\textbf{A.} 9/2\quad\quad \textbf{B.} 21\quad\quad \textbf{C.} 0\quad\quad \textbf{D.} -21\par\vspace{2pt}
\textbf{\color{primary} Câu 50.} Câu hỏi vận dụng cao cuối Chương I: So sánh hai lũy thừa A = 2$^{39}$ và B = 3$^{26}$.. Phương án đúng là:\par
\hspace*{1.5em}\textbf{A.} Không so sánh được\quad\quad \textbf{B.} A < B (vì 8 < 9)\quad\quad \textbf{C.} A = B\quad\quad \textbf{D.} A > B\par\vspace{2pt}
\subsection*{\color{primary} Bảng đáp án Chương I: Số hữu tỉ}
\begin{center}
\renewcommand{\arraystretch}{1.25}
\begin{tabularx}{\textwidth}{*{10}{>{\centering\arraybackslash}X}}
\toprule
\textbf{\color{primary} Cột 1} & \textbf{\color{primary} Cột 2} & \textbf{\color{primary} Cột 3} & \textbf{\color{primary} Cột 4} & \textbf{\color{primary} Cột 5} & \textbf{\color{primary} Cột 6} & \textbf{\color{primary} Cột 7} & \textbf{\color{primary} Cột 8} & \textbf{\color{primary} Cột 9} & \textbf{\color{primary} Cột 10} \\
\midrule
\textbf{1.A} & \textbf{2.A} & \textbf{3.B} & \textbf{4.A} & \textbf{5.C} & \textbf{6.D} & \textbf{7.D} & \textbf{8.A} & \textbf{9.D} & \textbf{10.B} \\
\rowcolor{rowtint} \textbf{11.B} & \textbf{12.B} & \textbf{13.D} & \textbf{14.C} & \textbf{15.C} & \textbf{16.C} & \textbf{17.D} & \textbf{18.B} & \textbf{19.A} & \textbf{20.C} \\
\textbf{21.C} & \textbf{22.D} & \textbf{23.C} & \textbf{24.D} & \textbf{25.C} & \textbf{26.A} & \textbf{27.A} & \textbf{28.D} & \textbf{29.A} & \textbf{30.A} \\
\rowcolor{rowtint} \textbf{31.A} & \textbf{32.D} & \textbf{33.A} & \textbf{34.C} & \textbf{35.D} & \textbf{36.B} & \textbf{37.B} & \textbf{38.D} & \textbf{39.C} & \textbf{40.B} \\
\textbf{41.D} & \textbf{42.A} & \textbf{43.C} & \textbf{44.B} & \textbf{45.A} & \textbf{46.B} & \textbf{47.A} & \textbf{48.D} & \textbf{49.B} & \textbf{50.B} \\
\bottomrule
\end{tabularx}
\end{center}\vspace{10pt}
\section*{\color{primary} CHƯƠNG II: SỐ THỰC}
\textit{Trọng tâm kiến thức SGK: Số thập phân vô hạn tuần hoàn, số vô tỉ, căn bậc hai số học, tập hợp các số thực $\mathbb{R}$ và giá trị tuyệt đối của một số thực.}\par\vspace{4pt}
\subsection*{\color{accent} Phần I. Mức độ Dễ (Nhận biết -- Câu 1 đến Câu 20)}
\textbf{\color{primary} Câu 1.} Căn bậc hai số học của 9 là:\par
\hspace*{1.5em}\textbf{A.} -3\quad\quad \textbf{B.} 81\quad\quad \textbf{C.} $\pm$3\quad\quad \textbf{D.} 3\par\vspace{2pt}
\textbf{\color{primary} Câu 2.} Thực hiện yêu cầu nhận biết sau: Giá trị tuyệt đối của số thực -5,5 tức là |-5,5| bằng.\par
\hspace*{1.5em}\textbf{A.} 0\quad\quad \textbf{B.} 5,5\quad\quad \textbf{C.} -5,5\quad\quad \textbf{D.} 10,5\par\vspace{2pt}
\textbf{\color{primary} Câu 3.} Trong bài học thuộc Chương II, xét nội dung: Trong các số sau đây, số nào là số vô tỉ. Đáp án đúng là:\par
\hspace*{1.5em}\textbf{A.} $\sqrt{9}$\quad\quad \textbf{B.} $\sqrt{5}$\quad\quad \textbf{C.} 1,25\quad\quad \textbf{D.} -3/4\par\vspace{2pt}
\textbf{\color{primary} Câu 4.} Dựa vào kiến thức chuẩn SGK của Chương II (Số thực), căn bậc hai số học của 36 là:\par
\hspace*{1.5em}\textbf{A.} 1296\quad\quad \textbf{B.} -6\quad\quad \textbf{C.} 6\quad\quad \textbf{D.} $\pm$6\par\vspace{2pt}
\textbf{\color{primary} Câu 5.} Chọn kết quả đúng nhất khi xét: Giá trị tuyệt đối của số thực -8,5 tức là |-8,5| bằng:\par
\hspace*{1.5em}\textbf{A.} 8,5\quad\quad \textbf{B.} 16,5\quad\quad \textbf{C.} 0\quad\quad \textbf{D.} -8,5\par\vspace{2pt}
\textbf{\color{primary} Câu 6.} Khi làm bài tập cơ bản về số thực, một học sinh gặp câu hỏi: Trong các số sau đây, số nào là số vô tỉ. Kết quả nào dưới đây đúng?\par
\hspace*{1.5em}\textbf{A.} -3/4\quad\quad \textbf{B.} $\sqrt{9}$\quad\quad \textbf{C.} $\sqrt{2}$\quad\quad \textbf{D.} 1,25\par\vspace{2pt}
\textbf{\color{primary} Câu 7.} Xác định giá trị hoặc khẳng định đúng cho bài toán sau: Căn bậc hai số học của 81 là:\par
\hspace*{1.5em}\textbf{A.} 9\quad\quad \textbf{B.} -9\quad\quad \textbf{C.} $\pm$9\quad\quad \textbf{D.} 6561\par\vspace{2pt}
\textbf{\color{primary} Câu 8.} Áp dụng trực tiếp định nghĩa và quy tắc trong Chương II: Giá trị tuyệt đối của số thực -11,5 tức là |-11,5| bằng:\par
\hspace*{1.5em}\textbf{A.} 11,5\quad\quad \textbf{B.} 0\quad\quad \textbf{C.} -11,5\quad\quad \textbf{D.} 22,5\par\vspace{2pt}
\textbf{\color{primary} Câu 9.} Kiểm tra nhanh kiến thức nền tảng: Trong các số sau đây, số nào là số vô tỉ?\par
\hspace*{1.5em}\textbf{A.} 1,25\quad\quad \textbf{B.} -3/4\quad\quad \textbf{C.} $\sqrt{7}$\quad\quad \textbf{D.} $\sqrt{9}$\par\vspace{2pt}
\textbf{\color{primary} Câu 10.} Cho dữ kiện nhận biết sau đây: Căn bậc hai số học của 144 là. Phương án chính xác là:\par
\hspace*{1.5em}\textbf{A.} -12\quad\quad \textbf{B.} $\pm$12\quad\quad \textbf{C.} 20736\quad\quad \textbf{D.} 12\par\vspace{2pt}
\textbf{\color{primary} Câu 11.} Hãy chỉ ra kết quả đúng trong bốn lựa chọn A, B, C, D đối với yêu cầu: Giá trị tuyệt đối của số thực -14,5 tức là |-14,5| bằng:\par
\hspace*{1.5em}\textbf{A.} 14,5\quad\quad \textbf{B.} 0\quad\quad \textbf{C.} 28,5\quad\quad \textbf{D.} -14,5\par\vspace{2pt}
\textbf{\color{primary} Câu 12.} Câu hỏi ôn tập lý thuyết và kỹ năng cơ bản (Chương II): Trong các số sau đây, số nào là số vô tỉ:\par
\hspace*{1.5em}\textbf{A.} -3/4\quad\quad \textbf{B.} 1,25\quad\quad \textbf{C.} $\sqrt{3}$\quad\quad \textbf{D.} $\sqrt{9}$\par\vspace{2pt}
\textbf{\color{primary} Câu 13.} Theo nội dung sách giáo khoa Kết nối tri thức với cuộc sống (Chương II), căn bậc hai số học của 225 là:\par
\hspace*{1.5em}\textbf{A.} $\pm$15\quad\quad \textbf{B.} 15\quad\quad \textbf{C.} -15\quad\quad \textbf{D.} 50625\par\vspace{2pt}
\textbf{\color{primary} Câu 14.} Tìm đáp số đúng của bài toán cơ bản sau: Giá trị tuyệt đối của số thực -17,5 tức là |-17,5| bằng:\par
\hspace*{1.5em}\textbf{A.} 34,5\quad\quad \textbf{B.} -17,5\quad\quad \textbf{C.} 17,5\quad\quad \textbf{D.} 0\par\vspace{2pt}
\textbf{\color{primary} Câu 15.} Xét đối tượng toán học trong Chương II: Trong các số sau đây, số nào là số vô tỉ. Giá trị hoặc tính chất tương ứng là:\par
\hspace*{1.5em}\textbf{A.} $\pi$\quad\quad \textbf{B.} -3/4\quad\quad \textbf{C.} 1,25\quad\quad \textbf{D.} $\sqrt{9}$\par\vspace{2pt}
\textbf{\color{primary} Câu 16.} Hoàn thành phép xác định cơ bản sau đây: Căn bậc hai số học của 324 là:\par
\hspace*{1.5em}\textbf{A.} -18\quad\quad \textbf{B.} $\pm$18\quad\quad \textbf{C.} 18\quad\quad \textbf{D.} 104976\par\vspace{2pt}
\textbf{\color{primary} Câu 17.} Trong các khẳng định hoặc giá trị dưới đây, đâu là đáp án đúng cho câu hỏi: Giá trị tuyệt đối của số thực -20,5 tức là |-20,5| bằng?\par
\hspace*{1.5em}\textbf{A.} 20,5\quad\quad \textbf{B.} 40,5\quad\quad \textbf{C.} -20,5\quad\quad \textbf{D.} 0\par\vspace{2pt}
\textbf{\color{primary} Câu 18.} Bài tập khởi động Chương II: Trong các số sau đây, số nào là số vô tỉ:\par
\hspace*{1.5em}\textbf{A.} $\sqrt{5}$\quad\quad \textbf{B.} $\sqrt{9}$\quad\quad \textbf{C.} 1,25\quad\quad \textbf{D.} -3/4\par\vspace{2pt}
\textbf{\color{primary} Câu 19.} Quan sát dữ kiện đề bài: Căn bậc hai số học của 441 là. Hãy chọn phương án đúng:\par
\hspace*{1.5em}\textbf{A.} 194481\quad\quad \textbf{B.} $\pm$21\quad\quad \textbf{C.} 21\quad\quad \textbf{D.} -21\par\vspace{2pt}
\textbf{\color{primary} Câu 20.} Sử dụng công thức hoặc tính chất đầu bài của Chương II để giải quyết yêu cầu: Giá trị tuyệt đối của số thực -23,5 tức là |-23,5| bằng:\par
\hspace*{1.5em}\textbf{A.} 23,5\quad\quad \textbf{B.} -23,5\quad\quad \textbf{C.} 46,5\quad\quad \textbf{D.} 0\par\vspace{2pt}
\subsection*{\color{accent} Phần II. Mức độ Trung bình (Thông hiểu -- Câu 21 đến Câu 40)}
\textbf{\color{primary} Câu 21.} Tính giá trị của biểu thức M = $\sqrt{4}$ + |-5|.:\par
\hspace*{1.5em}\textbf{A.} 3\quad\quad \textbf{B.} 9\quad\quad \textbf{C.} 7\quad\quad \textbf{D.} -3\par\vspace{2pt}
\textbf{\color{primary} Câu 22.} Vận dụng quy tắc trong Chương II (Số thực) để giải bài toán: Tìm tất cả các số thực x thỏa mãn |x| = 7.:\par
\hspace*{1.5em}\textbf{A.} x = 7 hoặc x = -7\quad\quad \textbf{B.} Chỉ có x = -7\quad\quad \textbf{C.} Chỉ có x = 7\quad\quad \textbf{D.} x = 49\par\vspace{2pt}
\textbf{\color{primary} Câu 23.} Biến đổi và tính toán theo kiến thức Chương II: Biết $\sqrt{x}$ = 6 (x $\ge$ 0). Giá trị của x bằng. Kết quả thu được là:\par
\hspace*{1.5em}\textbf{A.} -36\quad\quad \textbf{B.} 12\quad\quad \textbf{C.} 36\quad\quad \textbf{D.} 6\par\vspace{2pt}
\textbf{\color{primary} Câu 24.} Xét bài toán mức độ thông hiểu sau: Tính giá trị của biểu thức M = $\sqrt{25}$ + |-8|.. Đáp án đúng là:\par
\hspace*{1.5em}\textbf{A.} 33\quad\quad \textbf{B.} 3\quad\quad \textbf{C.} -3\quad\quad \textbf{D.} 13\par\vspace{2pt}
\textbf{\color{primary} Câu 25.} Để hoàn thành phiếu học tập Chương II, hãy xác định: Tìm tất cả các số thực x thỏa mãn |x| = 10.:\par
\hspace*{1.5em}\textbf{A.} x = 10 hoặc x = -10\quad\quad \textbf{B.} Chỉ có x = -10\quad\quad \textbf{C.} x = 100\quad\quad \textbf{D.} Chỉ có x = 10\par\vspace{2pt}
\textbf{\color{primary} Câu 26.} Phân tích dữ kiện bài toán: Biết $\sqrt{x}$ = 9 (x $\ge$ 0). Giá trị của x bằng. Giá trị cần tìm bằng:\par
\hspace*{1.5em}\textbf{A.} 81\quad\quad \textbf{B.} 18\quad\quad \textbf{C.} 9\quad\quad \textbf{D.} -81\par\vspace{2pt}
\textbf{\color{primary} Câu 27.} Áp dụng tính chất đã học ở Chương II vào bài tập sau: Tính giá trị của biểu thức M = $\sqrt{64}$ + |-11|.:\par
\hspace*{1.5em}\textbf{A.} 75\quad\quad \textbf{B.} 3\quad\quad \textbf{C.} 19\quad\quad \textbf{D.} -3\par\vspace{2pt}
\textbf{\color{primary} Câu 28.} Thực hiện các bước suy luận hợp lý để giải quyết yêu cầu: Tìm tất cả các số thực x thỏa mãn |x| = 13.:\par
\hspace*{1.5em}\textbf{A.} Chỉ có x = 13\quad\quad \textbf{B.} x = 169\quad\quad \textbf{C.} Chỉ có x = -13\quad\quad \textbf{D.} x = 13 hoặc x = -13\par\vspace{2pt}
\textbf{\color{primary} Câu 29.} Cho bài toán tự luyện (Số thực): Biết $\sqrt{x}$ = 12 (x $\ge$ 0). Giá trị của x bằng. Phương án nào sau đây chính xác?\par
\hspace*{1.5em}\textbf{A.} 24\quad\quad \textbf{B.} 12\quad\quad \textbf{C.} 144\quad\quad \textbf{D.} -144\par\vspace{2pt}
\textbf{\color{primary} Câu 30.} Kết hợp các phép biến đổi trong Chương II, hãy tìm kết quả của bài toán: Tính giá trị của biểu thức M = $\sqrt{121}$ + |-14|.:\par
\hspace*{1.5em}\textbf{A.} -3\quad\quad \textbf{B.} 135\quad\quad \textbf{C.} 25\quad\quad \textbf{D.} 3\par\vspace{2pt}
\textbf{\color{primary} Câu 31.} Một bài kiểm tra định kỳ nội dung Chương II nêu câu hỏi: Tìm tất cả các số thực x thỏa mãn |x| = 16.. Hãy chọn đáp số đúng:\par
\hspace*{1.5em}\textbf{A.} x = 16 hoặc x = -16\quad\quad \textbf{B.} Chỉ có x = 16\quad\quad \textbf{C.} Chỉ có x = -16\quad\quad \textbf{D.} x = 256\par\vspace{2pt}
\textbf{\color{primary} Câu 32.} Tìm lời giải đúng cho bài toán thông hiểu: Biết $\sqrt{x}$ = 15 (x $\ge$ 0). Giá trị của x bằng:\par
\hspace*{1.5em}\textbf{A.} 15\quad\quad \textbf{B.} -225\quad\quad \textbf{C.} 30\quad\quad \textbf{D.} 225\par\vspace{2pt}
\textbf{\color{primary} Câu 33.} Sử dụng mối liên hệ giữa các đại lượng trong Chương II: Tính giá trị của biểu thức M = $\sqrt{196}$ + |-17|.:\par
\hspace*{1.5em}\textbf{A.} 3\quad\quad \textbf{B.} 31\quad\quad \textbf{C.} 213\quad\quad \textbf{D.} -3\par\vspace{2pt}
\textbf{\color{primary} Câu 34.} Hãy tính toán và chọn phương án đúng đối với bài tập: Tìm tất cả các số thực x thỏa mãn |x| = 19.:\par
\hspace*{1.5em}\textbf{A.} Chỉ có x = -19\quad\quad \textbf{B.} x = 361\quad\quad \textbf{C.} Chỉ có x = 19\quad\quad \textbf{D.} x = 19 hoặc x = -19\par\vspace{2pt}
\textbf{\color{primary} Câu 35.} Dựa trên các tính chất trọng tâm của Chương II, xét bài toán: Biết $\sqrt{x}$ = 18 (x $\ge$ 0). Giá trị của x bằng:\par
\hspace*{1.5em}\textbf{A.} -324\quad\quad \textbf{B.} 18\quad\quad \textbf{C.} 36\quad\quad \textbf{D.} 324\par\vspace{2pt}
\textbf{\color{primary} Câu 36.} Thực hành giải toán (Số thực): Tính giá trị của biểu thức M = $\sqrt{289}$ + |-20|.. Đáp án nào dưới đây thỏa mãn?\par
\hspace*{1.5em}\textbf{A.} -3\quad\quad \textbf{B.} 37\quad\quad \textbf{C.} 3\quad\quad \textbf{D.} 309\par\vspace{2pt}
\textbf{\color{primary} Câu 37.} Kiểm tra kỹ năng biến đổi và suy luận qua bài toán: Tìm tất cả các số thực x thỏa mãn |x| = 22.:\par
\hspace*{1.5em}\textbf{A.} x = 484\quad\quad \textbf{B.} x = 22 hoặc x = -22\quad\quad \textbf{C.} Chỉ có x = 22\quad\quad \textbf{D.} Chỉ có x = -22\par\vspace{2pt}
\textbf{\color{primary} Câu 38.} Cho giả thiết thuộc phạm vi Chương II: Biết $\sqrt{x}$ = 21 (x $\ge$ 0). Giá trị của x bằng. Kết luận đúng là:\par
\hspace*{1.5em}\textbf{A.} 42\quad\quad \textbf{B.} -441\quad\quad \textbf{C.} 21\quad\quad \textbf{D.} 441\par\vspace{2pt}
\textbf{\color{primary} Câu 39.} Rút gọn hoặc tính giá trị theo yêu cầu sau của Chương II: Tính giá trị của biểu thức M = $\sqrt{400}$ + |-23|.:\par
\hspace*{1.5em}\textbf{A.} 43\quad\quad \textbf{B.} 423\quad\quad \textbf{C.} 3\quad\quad \textbf{D.} -3\par\vspace{2pt}
\textbf{\color{primary} Câu 40.} Xác định phương án đúng sau khi thực hiện phép tính / suy luận: Tìm tất cả các số thực x thỏa mãn |x| = 25.:\par
\hspace*{1.5em}\textbf{A.} Chỉ có x = -25\quad\quad \textbf{B.} x = 25 hoặc x = -25\quad\quad \textbf{C.} x = 625\quad\quad \textbf{D.} Chỉ có x = 25\par\vspace{2pt}
\subsection*{\color{accent} Phần III. Mức độ Khó (Vận dụng \& Vận dụng cao -- Câu 41 đến Câu 50)}
\textbf{\color{primary} Câu 41.} Tìm giá trị lớn nhất của biểu thức Q = 15 - |2x - 3|.:\par
\hspace*{1.5em}\textbf{A.} 15\quad\quad \textbf{B.} -15\quad\quad \textbf{C.} 3\quad\quad \textbf{D.} 0\par\vspace{2pt}
\textbf{\color{primary} Câu 42.} Bài toán vận dụng (Chương II – Số thực): Tìm số thực x không âm thỏa mãn 3$\cdot$$\sqrt{x}$ - 8 = 4.. Kết quả đúng là:\par
\hspace*{1.5em}\textbf{A.} x = 36\quad\quad \textbf{B.} x = 16\quad\quad \textbf{C.} x = 12\quad\quad \textbf{D.} x = 4\par\vspace{2pt}
\textbf{\color{primary} Câu 43.} Vận dụng kiến thức nâng cao trong phạm vi Chương II để giải bài toán: Tìm giá trị lớn nhất của biểu thức Q = 19 - |2x - 5|.:\par
\hspace*{1.5em}\textbf{A.} 0\quad\quad \textbf{B.} 19\quad\quad \textbf{C.} 5\quad\quad \textbf{D.} -19\par\vspace{2pt}
\textbf{\color{primary} Câu 44.} Phân tích và giải bài toán tổng hợp sau của Chương II: Tìm số thực x không âm thỏa mãn 3$\cdot$$\sqrt{x}$ - 12 = 6.:\par
\hspace*{1.5em}\textbf{A.} x = 6\quad\quad \textbf{B.} x = 54\quad\quad \textbf{C.} x = 18\quad\quad \textbf{D.} x = 36\par\vspace{2pt}
\textbf{\color{primary} Câu 45.} Thử thách tư duy trong phạm vi kiến thức Chương II: Tìm giá trị lớn nhất của biểu thức Q = 23 - |2x - 7|.. Đáp số chính xác là:\par
\hspace*{1.5em}\textbf{A.} 7\quad\quad \textbf{B.} 23\quad\quad \textbf{C.} -23\quad\quad \textbf{D.} 0\par\vspace{2pt}
\textbf{\color{primary} Câu 46.} Sử dụng phương pháp giải toán của Chương II để tìm đáp án cho bài toán: Tìm số thực x không âm thỏa mãn 3$\cdot$$\sqrt{x}$ - 16 = 8.:\par
\hspace*{1.5em}\textbf{A.} x = 72\quad\quad \textbf{B.} x = 64\quad\quad \textbf{C.} x = 24\quad\quad \textbf{D.} x = 8\par\vspace{2pt}
\textbf{\color{primary} Câu 47.} Xét bài toán vận dụng thực tiễn / nâng cao thuộc Chương II: Tìm giá trị lớn nhất của biểu thức Q = 27 - |2x - 9|.:\par
\hspace*{1.5em}\textbf{A.} 9\quad\quad \textbf{B.} 27\quad\quad \textbf{C.} -27\quad\quad \textbf{D.} 0\par\vspace{2pt}
\textbf{\color{primary} Câu 48.} Tìm kết quả cuối cùng của bài toán phân loại học sinh khá giỏi (Chương II): Tìm số thực x không âm thỏa mãn 3$\cdot$$\sqrt{x}$ - 20 = 10.:\par
\hspace*{1.5em}\textbf{A.} x = 10\quad\quad \textbf{B.} x = 100\quad\quad \textbf{C.} x = 90\quad\quad \textbf{D.} x = 30\par\vspace{2pt}
\textbf{\color{primary} Câu 49.} Kết hợp các tính chất đã học trong Chương II, hãy giải quyết bài toán: Tìm giá trị lớn nhất của biểu thức Q = 31 - |2x - 11|.:\par
\hspace*{1.5em}\textbf{A.} 0\quad\quad \textbf{B.} 31\quad\quad \textbf{C.} -31\quad\quad \textbf{D.} 11\par\vspace{2pt}
\textbf{\color{primary} Câu 50.} Câu hỏi vận dụng cao cuối Chương II: Tìm số thực x không âm thỏa mãn 3$\cdot$$\sqrt{x}$ - 24 = 12.. Phương án đúng là:\par
\hspace*{1.5em}\textbf{A.} x = 36\quad\quad \textbf{B.} x = 144\quad\quad \textbf{C.} x = 108\quad\quad \textbf{D.} x = 12\par\vspace{2pt}
\subsection*{\color{primary} Bảng đáp án Chương II: Số thực}
\begin{center}
\renewcommand{\arraystretch}{1.25}
\begin{tabularx}{\textwidth}{*{10}{>{\centering\arraybackslash}X}}
\toprule
\textbf{\color{primary} Cột 1} & \textbf{\color{primary} Cột 2} & \textbf{\color{primary} Cột 3} & \textbf{\color{primary} Cột 4} & \textbf{\color{primary} Cột 5} & \textbf{\color{primary} Cột 6} & \textbf{\color{primary} Cột 7} & \textbf{\color{primary} Cột 8} & \textbf{\color{primary} Cột 9} & \textbf{\color{primary} Cột 10} \\
\midrule
\textbf{1.D} & \textbf{2.B} & \textbf{3.B} & \textbf{4.C} & \textbf{5.A} & \textbf{6.C} & \textbf{7.A} & \textbf{8.A} & \textbf{9.C} & \textbf{10.D} \\
\rowcolor{rowtint} \textbf{11.A} & \textbf{12.C} & \textbf{13.B} & \textbf{14.C} & \textbf{15.A} & \textbf{16.C} & \textbf{17.A} & \textbf{18.A} & \textbf{19.C} & \textbf{20.A} \\
\textbf{21.C} & \textbf{22.A} & \textbf{23.C} & \textbf{24.D} & \textbf{25.A} & \textbf{26.A} & \textbf{27.C} & \textbf{28.D} & \textbf{29.C} & \textbf{30.C} \\
\rowcolor{rowtint} \textbf{31.A} & \textbf{32.D} & \textbf{33.B} & \textbf{34.D} & \textbf{35.D} & \textbf{36.B} & \textbf{37.B} & \textbf{38.D} & \textbf{39.A} & \textbf{40.B} \\
\textbf{41.A} & \textbf{42.B} & \textbf{43.B} & \textbf{44.D} & \textbf{45.B} & \textbf{46.B} & \textbf{47.B} & \textbf{48.B} & \textbf{49.B} & \textbf{50.B} \\
\bottomrule
\end{tabularx}
\end{center}\vspace{10pt}
\section*{\color{primary} CHƯƠNG III: GÓC VÀ ĐƯỜNG THẲNG SONG SONG}
\textit{Trọng tâm kiến thức SGK: Góc ở vị trí đặc biệt (hai góc kề bù, hai góc đối đỉnh), tia phân giác của một góc, dấu hiệu nhận biết và tính chất hai đường thẳng song song, tiên đề Euclid, định lí.}\par\vspace{4pt}
\subsection*{\color{accent} Phần I. Mức độ Dễ (Nhận biết -- Câu 1 đến Câu 20)}
\textbf{\color{primary} Câu 1.} Cho hai góc xOy và yOz là hai góc kề bù. Biết số đo góc xOy = 40°. Số đo góc yOz bằng:\par
\hspace*{1.5em}\textbf{A.} 50°\quad\quad \textbf{B.} 40°\quad\quad \textbf{C.} 140°\quad\quad \textbf{D.} 150°\par\vspace{2pt}
\textbf{\color{primary} Câu 2.} Thực hiện yêu cầu nhận biết sau: Hai đường thẳng xx' và yy' cắt nhau tại O tạo thành góc xOy = 39°. Số đo góc đối đỉnh x'Oy' bằng.\par
\hspace*{1.5em}\textbf{A.} 51°\quad\quad \textbf{B.} 78°\quad\quad \textbf{C.} 39°\quad\quad \textbf{D.} 141°\par\vspace{2pt}
\textbf{\color{primary} Câu 3.} Trong bài học thuộc Chương III, xét nội dung: Cho góc mOn có số đo bằng 62°. Tia Oz là tia phân giác của góc mOn thì số đo góc mOz bằng. Đáp án đúng là:\par
\hspace*{1.5em}\textbf{A.} 124°\quad\quad \textbf{B.} 118°\quad\quad \textbf{C.} 31°\quad\quad \textbf{D.} 62°\par\vspace{2pt}
\textbf{\color{primary} Câu 4.} Dựa vào kiến thức chuẩn SGK của Chương III (Góc và đường thẳng song song), cho hai góc xOy và yOz là hai góc kề bù. Biết số đo góc xOy = 55°. Số đo góc yOz bằng:\par
\hspace*{1.5em}\textbf{A.} 35°\quad\quad \textbf{B.} 55°\quad\quad \textbf{C.} 135°\quad\quad \textbf{D.} 125°\par\vspace{2pt}
\textbf{\color{primary} Câu 5.} Chọn kết quả đúng nhất khi xét: Hai đường thẳng xx' và yy' cắt nhau tại O tạo thành góc xOy = 51°. Số đo góc đối đỉnh x'Oy' bằng:\par
\hspace*{1.5em}\textbf{A.} 51°\quad\quad \textbf{B.} 129°\quad\quad \textbf{C.} 39°\quad\quad \textbf{D.} 102°\par\vspace{2pt}
\textbf{\color{primary} Câu 6.} Khi làm bài tập cơ bản về góc và đường thẳng song song, một học sinh gặp câu hỏi: Cho góc mOn có số đo bằng 80°. Tia Oz là tia phân giác của góc mOn thì số đo góc mOz bằng. Kết quả nào dưới đây đúng?\par
\hspace*{1.5em}\textbf{A.} 160°\quad\quad \textbf{B.} 40°\quad\quad \textbf{C.} 80°\quad\quad \textbf{D.} 100°\par\vspace{2pt}
\textbf{\color{primary} Câu 7.} Xác định giá trị hoặc khẳng định đúng cho bài toán sau: Cho hai góc xOy và yOz là hai góc kề bù. Biết số đo góc xOy = 70°. Số đo góc yOz bằng:\par
\hspace*{1.5em}\textbf{A.} 20°\quad\quad \textbf{B.} 70°\quad\quad \textbf{C.} 120°\quad\quad \textbf{D.} 110°\par\vspace{2pt}
\textbf{\color{primary} Câu 8.} Áp dụng trực tiếp định nghĩa và quy tắc trong Chương III: Hai đường thẳng xx' và yy' cắt nhau tại O tạo thành góc xOy = 63°. Số đo góc đối đỉnh x'Oy' bằng:\par
\hspace*{1.5em}\textbf{A.} 27°\quad\quad \textbf{B.} 63°\quad\quad \textbf{C.} 126°\quad\quad \textbf{D.} 117°\par\vspace{2pt}
\textbf{\color{primary} Câu 9.} Kiểm tra nhanh kiến thức nền tảng: Cho góc mOn có số đo bằng 98°. Tia Oz là tia phân giác của góc mOn thì số đo góc mOz bằng:\par
\hspace*{1.5em}\textbf{A.} 49°\quad\quad \textbf{B.} 196°\quad\quad \textbf{C.} 98°\quad\quad \textbf{D.} 82°\par\vspace{2pt}
\textbf{\color{primary} Câu 10.} Cho dữ kiện nhận biết sau đây: Cho hai góc xOy và yOz là hai góc kề bù. Biết số đo góc xOy = 85°. Số đo góc yOz bằng. Phương án chính xác là:\par
\hspace*{1.5em}\textbf{A.} 85°\quad\quad \textbf{B.} 105°\quad\quad \textbf{C.} 95°\quad\quad \textbf{D.} 5°\par\vspace{2pt}
\textbf{\color{primary} Câu 11.} Hãy chỉ ra kết quả đúng trong bốn lựa chọn A, B, C, D đối với yêu cầu: Hai đường thẳng xx' và yy' cắt nhau tại O tạo thành góc xOy = 75°. Số đo góc đối đỉnh x'Oy' bằng:\par
\hspace*{1.5em}\textbf{A.} 75°\quad\quad \textbf{B.} 150°\quad\quad \textbf{C.} 15°\quad\quad \textbf{D.} 105°\par\vspace{2pt}
\textbf{\color{primary} Câu 12.} Câu hỏi ôn tập lý thuyết và kỹ năng cơ bản (Chương III): Cho góc mOn có số đo bằng 116°. Tia Oz là tia phân giác của góc mOn thì số đo góc mOz bằng:\par
\hspace*{1.5em}\textbf{A.} 58°\quad\quad \textbf{B.} 232°\quad\quad \textbf{C.} 116°\quad\quad \textbf{D.} 64°\par\vspace{2pt}
\textbf{\color{primary} Câu 13.} Theo nội dung sách giáo khoa Kết nối tri thức với cuộc sống (Chương III), cho hai góc xOy và yOz là hai góc kề bù. Biết số đo góc xOy = 100°. Số đo góc yOz bằng:\par
\hspace*{1.5em}\textbf{A.} 80°\quad\quad \textbf{B.} 90°\quad\quad \textbf{C.} 100°\quad\quad \textbf{D.} -10°\par\vspace{2pt}
\textbf{\color{primary} Câu 14.} Tìm đáp số đúng của bài toán cơ bản sau: Hai đường thẳng xx' và yy' cắt nhau tại O tạo thành góc xOy = 87°. Số đo góc đối đỉnh x'Oy' bằng:\par
\hspace*{1.5em}\textbf{A.} 87°\quad\quad \textbf{B.} 174°\quad\quad \textbf{C.} 93°\quad\quad \textbf{D.} 3°\par\vspace{2pt}
\textbf{\color{primary} Câu 15.} Xét đối tượng toán học trong Chương III: Cho góc mOn có số đo bằng 134°. Tia Oz là tia phân giác của góc mOn thì số đo góc mOz bằng. Giá trị hoặc tính chất tương ứng là:\par
\hspace*{1.5em}\textbf{A.} 67°\quad\quad \textbf{B.} 46°\quad\quad \textbf{C.} 134°\quad\quad \textbf{D.} 268°\par\vspace{2pt}
\textbf{\color{primary} Câu 16.} Hoàn thành phép xác định cơ bản sau đây: Cho hai góc xOy và yOz là hai góc kề bù. Biết số đo góc xOy = 115°. Số đo góc yOz bằng:\par
\hspace*{1.5em}\textbf{A.} -25°\quad\quad \textbf{B.} 65°\quad\quad \textbf{C.} 115°\quad\quad \textbf{D.} 75°\par\vspace{2pt}
\textbf{\color{primary} Câu 17.} Trong các khẳng định hoặc giá trị dưới đây, đâu là đáp án đúng cho câu hỏi: Hai đường thẳng xx' và yy' cắt nhau tại O tạo thành góc xOy = 99°. Số đo góc đối đỉnh x'Oy' bằng?\par
\hspace*{1.5em}\textbf{A.} 198°\quad\quad \textbf{B.} -9°\quad\quad \textbf{C.} 99°\quad\quad \textbf{D.} 81°\par\vspace{2pt}
\textbf{\color{primary} Câu 18.} Bài tập khởi động Chương III: Cho góc mOn có số đo bằng 152°. Tia Oz là tia phân giác của góc mOn thì số đo góc mOz bằng:\par
\hspace*{1.5em}\textbf{A.} 76°\quad\quad \textbf{B.} 28°\quad\quad \textbf{C.} 152°\quad\quad \textbf{D.} 304°\par\vspace{2pt}
\textbf{\color{primary} Câu 19.} Quan sát dữ kiện đề bài: Cho hai góc xOy và yOz là hai góc kề bù. Biết số đo góc xOy = 130°. Số đo góc yOz bằng. Hãy chọn phương án đúng:\par
\hspace*{1.5em}\textbf{A.} 60°\quad\quad \textbf{B.} 50°\quad\quad \textbf{C.} -40°\quad\quad \textbf{D.} 130°\par\vspace{2pt}
\textbf{\color{primary} Câu 20.} Sử dụng công thức hoặc tính chất đầu bài của Chương III để giải quyết yêu cầu: Hai đường thẳng xx' và yy' cắt nhau tại O tạo thành góc xOy = 111°. Số đo góc đối đỉnh x'Oy' bằng:\par
\hspace*{1.5em}\textbf{A.} -21°\quad\quad \textbf{B.} 222°\quad\quad \textbf{C.} 69°\quad\quad \textbf{D.} 111°\par\vspace{2pt}
\subsection*{\color{accent} Phần II. Mức độ Trung bình (Thông hiểu -- Câu 21 đến Câu 40)}
\textbf{\color{primary} Câu 21.} Đường thẳng c cắt hai đường thẳng song song a và b. Biết một góc so le trong bằng 50°, góc so le trong còn lại ứng với nó có số đo là:\par
\hspace*{1.5em}\textbf{A.} 40°\quad\quad \textbf{B.} 50°\quad\quad \textbf{C.} 90°\quad\quad \textbf{D.} 130°\par\vspace{2pt}
\textbf{\color{primary} Câu 22.} Vận dụng quy tắc trong Chương III (Góc và đường thẳng song song) để giải bài toán: Đường thẳng c cắt hai đường thẳng song song a và b tạo thành một cặp góc trong cùng phía A$_1$ và B$_2$. Biết góc A$_1$ = 63°, số đo góc B$_2$ bằng:\par
\hspace*{1.5em}\textbf{A.} 27°\quad\quad \textbf{B.} 117°\quad\quad \textbf{C.} 63°\quad\quad \textbf{D.} 127°\par\vspace{2pt}
\textbf{\color{primary} Câu 23.} Biến đổi và tính toán theo kiến thức Chương III: Theo tiên đề Euclid, qua một điểm M nằm ngoài đường thẳng a cho trước, có bao nhiêu đường thẳng đi qua M và song song với a. Kết quả thu được là:\par
\begin{itemize}\setlength{\itemsep}{1pt}\setlength{\parskip}{0pt}
  \item[] \textbf{A.} Có vô số đường thẳng
  \item[] \textbf{B.} Không có đường thẳng nào
  \item[] \textbf{C.} Chỉ có một đường thẳng duy nhất
  \item[] \textbf{D.} Có hai đường thẳng
\end{itemize}\vspace{2pt}
\textbf{\color{primary} Câu 24.} Xét bài toán mức độ thông hiểu sau: Đường thẳng c cắt hai đường thẳng song song a và b. Biết một góc so le trong bằng 62°, góc so le trong còn lại ứng với nó có số đo là. Đáp án đúng là:\par
\hspace*{1.5em}\textbf{A.} 90°\quad\quad \textbf{B.} 118°\quad\quad \textbf{C.} 28°\quad\quad \textbf{D.} 62°\par\vspace{2pt}
\textbf{\color{primary} Câu 25.} Để hoàn thành phiếu học tập Chương III, hãy xác định: Đường thẳng c cắt hai đường thẳng song song a và b tạo thành một cặp góc trong cùng phía A$_1$ và B$_2$. Biết góc A$_1$ = 72°, số đo góc B$_2$ bằng:\par
\hspace*{1.5em}\textbf{A.} 18°\quad\quad \textbf{B.} 118°\quad\quad \textbf{C.} 72°\quad\quad \textbf{D.} 108°\par\vspace{2pt}
\textbf{\color{primary} Câu 26.} Phân tích dữ kiện bài toán: Theo tiên đề Euclid, qua một điểm M nằm ngoài đường thẳng a cho trước, có bao nhiêu đường thẳng đi qua M và song song với a. Giá trị cần tìm bằng:\par
\begin{itemize}\setlength{\itemsep}{1pt}\setlength{\parskip}{0pt}
  \item[] \textbf{A.} Có vô số đường thẳng
  \item[] \textbf{B.} Không có đường thẳng nào
  \item[] \textbf{C.} Chỉ có một đường thẳng duy nhất
  \item[] \textbf{D.} Có hai đường thẳng
\end{itemize}\vspace{2pt}
\textbf{\color{primary} Câu 27.} Áp dụng tính chất đã học ở Chương III vào bài tập sau: Đường thẳng c cắt hai đường thẳng song song a và b. Biết một góc so le trong bằng 74°, góc so le trong còn lại ứng với nó có số đo là:\par
\hspace*{1.5em}\textbf{A.} 16°\quad\quad \textbf{B.} 106°\quad\quad \textbf{C.} 90°\quad\quad \textbf{D.} 74°\par\vspace{2pt}
\textbf{\color{primary} Câu 28.} Thực hiện các bước suy luận hợp lý để giải quyết yêu cầu: Đường thẳng c cắt hai đường thẳng song song a và b tạo thành một cặp góc trong cùng phía A$_1$ và B$_2$. Biết góc A$_1$ = 81°, số đo góc B$_2$ bằng:\par
\hspace*{1.5em}\textbf{A.} 99°\quad\quad \textbf{B.} 109°\quad\quad \textbf{C.} 81°\quad\quad \textbf{D.} 9°\par\vspace{2pt}
\textbf{\color{primary} Câu 29.} Cho bài toán tự luyện (Góc và đường thẳng song song): Theo tiên đề Euclid, qua một điểm M nằm ngoài đường thẳng a cho trước, có bao nhiêu đường thẳng đi qua M và song song với a. Phương án nào sau đây chính xác?\par
\begin{itemize}\setlength{\itemsep}{1pt}\setlength{\parskip}{0pt}
  \item[] \textbf{A.} Có vô số đường thẳng
  \item[] \textbf{B.} Chỉ có một đường thẳng duy nhất
  \item[] \textbf{C.} Không có đường thẳng nào
  \item[] \textbf{D.} Có hai đường thẳng
\end{itemize}\vspace{2pt}
\textbf{\color{primary} Câu 30.} Kết hợp các phép biến đổi trong Chương III, hãy tìm kết quả của bài toán: Đường thẳng c cắt hai đường thẳng song song a và b. Biết một góc so le trong bằng 86°, góc so le trong còn lại ứng với nó có số đo là:\par
\hspace*{1.5em}\textbf{A.} 90°\quad\quad \textbf{B.} 4°\quad\quad \textbf{C.} 94°\quad\quad \textbf{D.} 86°\par\vspace{2pt}
\textbf{\color{primary} Câu 31.} Một bài kiểm tra định kỳ nội dung Chương III nêu câu hỏi: Đường thẳng c cắt hai đường thẳng song song a và b tạo thành một cặp góc trong cùng phía A$_1$ và B$_2$. Biết góc A$_1$ = 90°, số đo góc B$_2$ bằng. Hãy chọn đáp số đúng:\par
\hspace*{1.5em}\textbf{A.} 90° (1)\quad\quad \textbf{B.} 100°\quad\quad \textbf{C.} 90°\quad\quad \textbf{D.} 0°\par\vspace{2pt}
\textbf{\color{primary} Câu 32.} Tìm lời giải đúng cho bài toán thông hiểu: Theo tiên đề Euclid, qua một điểm M nằm ngoài đường thẳng a cho trước, có bao nhiêu đường thẳng đi qua M và song song với a:\par
\begin{itemize}\setlength{\itemsep}{1pt}\setlength{\parskip}{0pt}
  \item[] \textbf{A.} Có vô số đường thẳng
  \item[] \textbf{B.} Không có đường thẳng nào
  \item[] \textbf{C.} Chỉ có một đường thẳng duy nhất
  \item[] \textbf{D.} Có hai đường thẳng
\end{itemize}\vspace{2pt}
\textbf{\color{primary} Câu 33.} Sử dụng mối liên hệ giữa các đại lượng trong Chương III: Đường thẳng c cắt hai đường thẳng song song a và b. Biết một góc so le trong bằng 98°, góc so le trong còn lại ứng với nó có số đo là:\par
\hspace*{1.5em}\textbf{A.} 90°\quad\quad \textbf{B.} 82°\quad\quad \textbf{C.} 98°\quad\quad \textbf{D.} -8°\par\vspace{2pt}
\textbf{\color{primary} Câu 34.} Hãy tính toán và chọn phương án đúng đối với bài tập: Đường thẳng c cắt hai đường thẳng song song a và b tạo thành một cặp góc trong cùng phía A$_1$ và B$_2$. Biết góc A$_1$ = 99°, số đo góc B$_2$ bằng:\par
\hspace*{1.5em}\textbf{A.} 91°\quad\quad \textbf{B.} 99°\quad\quad \textbf{C.} 81°\quad\quad \textbf{D.} -9°\par\vspace{2pt}
\textbf{\color{primary} Câu 35.} Dựa trên các tính chất trọng tâm của Chương III, xét bài toán: Theo tiên đề Euclid, qua một điểm M nằm ngoài đường thẳng a cho trước, có bao nhiêu đường thẳng đi qua M và song song với a:\par
\begin{itemize}\setlength{\itemsep}{1pt}\setlength{\parskip}{0pt}
  \item[] \textbf{A.} Chỉ có một đường thẳng duy nhất
  \item[] \textbf{B.} Không có đường thẳng nào
  \item[] \textbf{C.} Có vô số đường thẳng
  \item[] \textbf{D.} Có hai đường thẳng
\end{itemize}\vspace{2pt}
\textbf{\color{primary} Câu 36.} Thực hành giải toán (Góc và đường thẳng song song): Đường thẳng c cắt hai đường thẳng song song a và b. Biết một góc so le trong bằng 110°, góc so le trong còn lại ứng với nó có số đo là. Đáp án nào dưới đây thỏa mãn?\par
\hspace*{1.5em}\textbf{A.} 70°\quad\quad \textbf{B.} 90°\quad\quad \textbf{C.} -20°\quad\quad \textbf{D.} 110°\par\vspace{2pt}
\textbf{\color{primary} Câu 37.} Kiểm tra kỹ năng biến đổi và suy luận qua bài toán: Đường thẳng c cắt hai đường thẳng song song a và b tạo thành một cặp góc trong cùng phía A$_1$ và B$_2$. Biết góc A$_1$ = 108°, số đo góc B$_2$ bằng:\par
\hspace*{1.5em}\textbf{A.} -18°\quad\quad \textbf{B.} 82°\quad\quad \textbf{C.} 108°\quad\quad \textbf{D.} 72°\par\vspace{2pt}
\textbf{\color{primary} Câu 38.} Cho giả thiết thuộc phạm vi Chương III: Theo tiên đề Euclid, qua một điểm M nằm ngoài đường thẳng a cho trước, có bao nhiêu đường thẳng đi qua M và song song với a. Kết luận đúng là:\par
\begin{itemize}\setlength{\itemsep}{1pt}\setlength{\parskip}{0pt}
  \item[] \textbf{A.} Có vô số đường thẳng
  \item[] \textbf{B.} Chỉ có một đường thẳng duy nhất
  \item[] \textbf{C.} Không có đường thẳng nào
  \item[] \textbf{D.} Có hai đường thẳng
\end{itemize}\vspace{2pt}
\textbf{\color{primary} Câu 39.} Rút gọn hoặc tính giá trị theo yêu cầu sau của Chương III: Đường thẳng c cắt hai đường thẳng song song a và b. Biết một góc so le trong bằng 122°, góc so le trong còn lại ứng với nó có số đo là:\par
\hspace*{1.5em}\textbf{A.} -32°\quad\quad \textbf{B.} 122°\quad\quad \textbf{C.} 58°\quad\quad \textbf{D.} 90°\par\vspace{2pt}
\textbf{\color{primary} Câu 40.} Xác định phương án đúng sau khi thực hiện phép tính / suy luận: Đường thẳng c cắt hai đường thẳng song song a và b tạo thành một cặp góc trong cùng phía A$_1$ và B$_2$. Biết góc A$_1$ = 117°, số đo góc B$_2$ bằng:\par
\hspace*{1.5em}\textbf{A.} 63°\quad\quad \textbf{B.} 117°\quad\quad \textbf{C.} -27°\quad\quad \textbf{D.} 73°\par\vspace{2pt}
\subsection*{\color{accent} Phần III. Mức độ Khó (Vận dụng \& Vận dụng cao -- Câu 41 đến Câu 50)}
\textbf{\color{primary} Câu 41.} Cho hai góc kề bù AOB và BOC, trong đó góc AOB = 60°. Gọi OM là tia phân giác của góc AOB và ON là tia phân giác của góc BOC. Số đo góc MON bằng:\par
\hspace*{1.5em}\textbf{A.} 60°\quad\quad \textbf{B.} 90°\quad\quad \textbf{C.} 180°\quad\quad \textbf{D.} 90° (1)\par\vspace{2pt}
\textbf{\color{primary} Câu 42.} Bài toán vận dụng (Chương III – Góc và đường thẳng song song): Cho hai đường thẳng song song a và b bị cắt bởi đường thẳng c tại A và B. Hai góc trong cùng phía có tỉ số số đo là 2 : 3. Hiệu số đo của góc lớn và góc nhỏ trong cặp góc đó bằng. Kết quả đúng là:\par
\hspace*{1.5em}\textbf{A.} 18°\quad\quad \textbf{B.} 36°\quad\quad \textbf{C.} 72°\quad\quad \textbf{D.} 108°\par\vspace{2pt}
\textbf{\color{primary} Câu 43.} Vận dụng kiến thức nâng cao trong phạm vi Chương III để giải bài toán: Cho hai góc kề bù AOB và BOC, trong đó góc AOB = 72°. Gọi OM là tia phân giác của góc AOB và ON là tia phân giác của góc BOC. Số đo góc MON bằng:\par
\hspace*{1.5em}\textbf{A.} 60°\quad\quad \textbf{B.} 180°\quad\quad \textbf{C.} 72°\quad\quad \textbf{D.} 90°\par\vspace{2pt}
\textbf{\color{primary} Câu 44.} Phân tích và giải bài toán tổng hợp sau của Chương III: Cho hai đường thẳng song song a và b bị cắt bởi đường thẳng c tại A và B. Hai góc trong cùng phía có tỉ số số đo là 2 : 3. Hiệu số đo của góc lớn và góc nhỏ trong cặp góc đó bằng:\par
\hspace*{1.5em}\textbf{A.} 108°\quad\quad \textbf{B.} 36°\quad\quad \textbf{C.} 18°\quad\quad \textbf{D.} 72°\par\vspace{2pt}
\textbf{\color{primary} Câu 45.} Thử thách tư duy trong phạm vi kiến thức Chương III: Cho hai góc kề bù AOB và BOC, trong đó góc AOB = 84°. Gọi OM là tia phân giác của góc AOB và ON là tia phân giác của góc BOC. Số đo góc MON bằng. Đáp số chính xác là:\par
\hspace*{1.5em}\textbf{A.} 84°\quad\quad \textbf{B.} 180°\quad\quad \textbf{C.} 90°\quad\quad \textbf{D.} 60°\par\vspace{2pt}
\textbf{\color{primary} Câu 46.} Sử dụng phương pháp giải toán của Chương III để tìm đáp án cho bài toán: Cho hai đường thẳng song song a và b bị cắt bởi đường thẳng c tại A và B. Hai góc trong cùng phía có tỉ số số đo là 2 : 3. Hiệu số đo của góc lớn và góc nhỏ trong cặp góc đó bằng:\par
\hspace*{1.5em}\textbf{A.} 36°\quad\quad \textbf{B.} 108°\quad\quad \textbf{C.} 72°\quad\quad \textbf{D.} 18°\par\vspace{2pt}
\textbf{\color{primary} Câu 47.} Xét bài toán vận dụng thực tiễn / nâng cao thuộc Chương III: Cho hai góc kề bù AOB và BOC, trong đó góc AOB = 96°. Gọi OM là tia phân giác của góc AOB và ON là tia phân giác của góc BOC. Số đo góc MON bằng:\par
\hspace*{1.5em}\textbf{A.} 60°\quad\quad \textbf{B.} 96°\quad\quad \textbf{C.} 180°\quad\quad \textbf{D.} 90°\par\vspace{2pt}
\textbf{\color{primary} Câu 48.} Tìm kết quả cuối cùng của bài toán phân loại học sinh khá giỏi (Chương III): Cho hai đường thẳng song song a và b bị cắt bởi đường thẳng c tại A và B. Hai góc trong cùng phía có tỉ số số đo là 2 : 3. Hiệu số đo của góc lớn và góc nhỏ trong cặp góc đó bằng:\par
\hspace*{1.5em}\textbf{A.} 72°\quad\quad \textbf{B.} 108°\quad\quad \textbf{C.} 18°\quad\quad \textbf{D.} 36°\par\vspace{2pt}
\textbf{\color{primary} Câu 49.} Kết hợp các tính chất đã học trong Chương III, hãy giải quyết bài toán: Cho hai góc kề bù AOB và BOC, trong đó góc AOB = 108°. Gọi OM là tia phân giác của góc AOB và ON là tia phân giác của góc BOC. Số đo góc MON bằng:\par
\hspace*{1.5em}\textbf{A.} 60°\quad\quad \textbf{B.} 180°\quad\quad \textbf{C.} 90°\quad\quad \textbf{D.} 108°\par\vspace{2pt}
\textbf{\color{primary} Câu 50.} Câu hỏi vận dụng cao cuối Chương III: Cho hai đường thẳng song song a và b bị cắt bởi đường thẳng c tại A và B. Hai góc trong cùng phía có tỉ số số đo là 2 : 3. Hiệu số đo của góc lớn và góc nhỏ trong cặp góc đó bằng. Phương án đúng là:\par
\hspace*{1.5em}\textbf{A.} 36°\quad\quad \textbf{B.} 72°\quad\quad \textbf{C.} 108°\quad\quad \textbf{D.} 18°\par\vspace{2pt}
\subsection*{\color{primary} Bảng đáp án Chương III: Góc và đường thẳng song song}
\begin{center}
\renewcommand{\arraystretch}{1.25}
\begin{tabularx}{\textwidth}{*{10}{>{\centering\arraybackslash}X}}
\toprule
\textbf{\color{primary} Cột 1} & \textbf{\color{primary} Cột 2} & \textbf{\color{primary} Cột 3} & \textbf{\color{primary} Cột 4} & \textbf{\color{primary} Cột 5} & \textbf{\color{primary} Cột 6} & \textbf{\color{primary} Cột 7} & \textbf{\color{primary} Cột 8} & \textbf{\color{primary} Cột 9} & \textbf{\color{primary} Cột 10} \\
\midrule
\textbf{1.C} & \textbf{2.C} & \textbf{3.C} & \textbf{4.D} & \textbf{5.A} & \textbf{6.B} & \textbf{7.D} & \textbf{8.B} & \textbf{9.A} & \textbf{10.C} \\
\rowcolor{rowtint} \textbf{11.A} & \textbf{12.A} & \textbf{13.A} & \textbf{14.A} & \textbf{15.A} & \textbf{16.B} & \textbf{17.C} & \textbf{18.A} & \textbf{19.B} & \textbf{20.D} \\
\textbf{21.B} & \textbf{22.B} & \textbf{23.C} & \textbf{24.D} & \textbf{25.D} & \textbf{26.C} & \textbf{27.D} & \textbf{28.A} & \textbf{29.B} & \textbf{30.D} \\
\rowcolor{rowtint} \textbf{31.C} & \textbf{32.C} & \textbf{33.C} & \textbf{34.C} & \textbf{35.A} & \textbf{36.D} & \textbf{37.D} & \textbf{38.B} & \textbf{39.B} & \textbf{40.A} \\
\textbf{41.B} & \textbf{42.B} & \textbf{43.D} & \textbf{44.B} & \textbf{45.C} & \textbf{46.A} & \textbf{47.D} & \textbf{48.D} & \textbf{49.C} & \textbf{50.A} \\
\bottomrule
\end{tabularx}
\end{center}\vspace{10pt}
\section*{\color{primary} CHƯƠNG IV: TAM GIÁC BẰNG NHAU}
\textit{Trọng tâm kiến thức SGK: Tổng các góc trong một tam giác, hai tam giác bằng nhau, các trường hợp bằng nhau của tam giác thường và tam giác vuông, tam giác cân, đường trung trực của đoạn thẳng.}\par\vspace{4pt}
\subsection*{\color{accent} Phần I. Mức độ Dễ (Nhận biết -- Câu 1 đến Câu 20)}
\textbf{\color{primary} Câu 1.} Cho tam giác ABC có góc A = 50° và góc B = 60°. Số đo góc C của tam giác ABC bằng:\par
\hspace*{1.5em}\textbf{A.} 60°\quad\quad \textbf{B.} 70°\quad\quad \textbf{C.} 110°\quad\quad \textbf{D.} 80°\par\vspace{2pt}
\textbf{\color{primary} Câu 2.} Thực hiện yêu cầu nhận biết sau: Cho $\Delta$ABC = $\Delta$DEF. Biết cạnh AB = 6 cm, BC = 8 cm, AC = 9 cm. Độ dài cạnh DE bằng.\par
\hspace*{1.5em}\textbf{A.} 6 cm\quad\quad \textbf{B.} 23 cm\quad\quad \textbf{C.} 9 cm\quad\quad \textbf{D.} 8 cm\par\vspace{2pt}
\textbf{\color{primary} Câu 3.} Trong bài học thuộc Chương IV, xét nội dung: Cho tam giác ABC cân tại A có góc ở đáy B = 49°. Số đo góc ở đáy C bằng. Đáp án đúng là:\par
\hspace*{1.5em}\textbf{A.} 82°\quad\quad \textbf{B.} 131°\quad\quad \textbf{C.} 90°\quad\quad \textbf{D.} 49°\par\vspace{2pt}
\textbf{\color{primary} Câu 4.} Dựa vào kiến thức chuẩn SGK của Chương IV (Tam giác bằng nhau), cho tam giác ABC có góc A = 56° và góc B = 63°. Số đo góc C của tam giác ABC bằng:\par
\hspace*{1.5em}\textbf{A.} 61°\quad\quad \textbf{B.} 119°\quad\quad \textbf{C.} 71°\quad\quad \textbf{D.} 51°\par\vspace{2pt}
\textbf{\color{primary} Câu 5.} Chọn kết quả đúng nhất khi xét: Cho $\Delta$ABC = $\Delta$DEF. Biết cạnh AB = 9 cm, BC = 11 cm, AC = 12 cm. Độ dài cạnh DE bằng:\par
\hspace*{1.5em}\textbf{A.} 9 cm\quad\quad \textbf{B.} 12 cm\quad\quad \textbf{C.} 32 cm\quad\quad \textbf{D.} 11 cm\par\vspace{2pt}
\textbf{\color{primary} Câu 6.} Khi làm bài tập cơ bản về tam giác bằng nhau, một học sinh gặp câu hỏi: Cho tam giác ABC cân tại A có góc ở đáy B = 55°. Số đo góc ở đáy C bằng. Kết quả nào dưới đây đúng?\par
\hspace*{1.5em}\textbf{A.} 90°\quad\quad \textbf{B.} 55°\quad\quad \textbf{C.} 70°\quad\quad \textbf{D.} 125°\par\vspace{2pt}
\textbf{\color{primary} Câu 7.} Xác định giá trị hoặc khẳng định đúng cho bài toán sau: Cho tam giác ABC có góc A = 62° và góc B = 66°. Số đo góc C của tam giác ABC bằng:\par
\hspace*{1.5em}\textbf{A.} 62°\quad\quad \textbf{B.} 128°\quad\quad \textbf{C.} 42°\quad\quad \textbf{D.} 52°\par\vspace{2pt}
\textbf{\color{primary} Câu 8.} Áp dụng trực tiếp định nghĩa và quy tắc trong Chương IV: Cho $\Delta$ABC = $\Delta$DEF. Biết cạnh AB = 12 cm, BC = 14 cm, AC = 15 cm. Độ dài cạnh DE bằng:\par
\hspace*{1.5em}\textbf{A.} 41 cm\quad\quad \textbf{B.} 12 cm\quad\quad \textbf{C.} 14 cm\quad\quad \textbf{D.} 15 cm\par\vspace{2pt}
\textbf{\color{primary} Câu 9.} Kiểm tra nhanh kiến thức nền tảng: Cho tam giác ABC cân tại A có góc ở đáy B = 61°. Số đo góc ở đáy C bằng:\par
\hspace*{1.5em}\textbf{A.} 119°\quad\quad \textbf{B.} 90°\quad\quad \textbf{C.} 61°\quad\quad \textbf{D.} 58°\par\vspace{2pt}
\textbf{\color{primary} Câu 10.} Cho dữ kiện nhận biết sau đây: Cho tam giác ABC có góc A = 68° và góc B = 69°. Số đo góc C của tam giác ABC bằng. Phương án chính xác là:\par
\hspace*{1.5em}\textbf{A.} 33°\quad\quad \textbf{B.} 43°\quad\quad \textbf{C.} 53°\quad\quad \textbf{D.} 137°\par\vspace{2pt}
\textbf{\color{primary} Câu 11.} Hãy chỉ ra kết quả đúng trong bốn lựa chọn A, B, C, D đối với yêu cầu: Cho $\Delta$ABC = $\Delta$DEF. Biết cạnh AB = 15 cm, BC = 17 cm, AC = 18 cm. Độ dài cạnh DE bằng:\par
\hspace*{1.5em}\textbf{A.} 17 cm\quad\quad \textbf{B.} 15 cm\quad\quad \textbf{C.} 18 cm\quad\quad \textbf{D.} 50 cm\par\vspace{2pt}
\textbf{\color{primary} Câu 12.} Câu hỏi ôn tập lý thuyết và kỹ năng cơ bản (Chương IV): Cho tam giác ABC cân tại A có góc ở đáy B = 67°. Số đo góc ở đáy C bằng:\par
\hspace*{1.5em}\textbf{A.} 113°\quad\quad \textbf{B.} 67°\quad\quad \textbf{C.} 46°\quad\quad \textbf{D.} 90°\par\vspace{2pt}
\textbf{\color{primary} Câu 13.} Theo nội dung sách giáo khoa Kết nối tri thức với cuộc sống (Chương IV), cho tam giác ABC có góc A = 74° và góc B = 72°. Số đo góc C của tam giác ABC bằng:\par
\hspace*{1.5em}\textbf{A.} 24°\quad\quad \textbf{B.} 44°\quad\quad \textbf{C.} 34°\quad\quad \textbf{D.} 146°\par\vspace{2pt}
\textbf{\color{primary} Câu 14.} Tìm đáp số đúng của bài toán cơ bản sau: Cho $\Delta$ABC = $\Delta$DEF. Biết cạnh AB = 18 cm, BC = 20 cm, AC = 21 cm. Độ dài cạnh DE bằng:\par
\hspace*{1.5em}\textbf{A.} 18 cm\quad\quad \textbf{B.} 59 cm\quad\quad \textbf{C.} 20 cm\quad\quad \textbf{D.} 21 cm\par\vspace{2pt}
\textbf{\color{primary} Câu 15.} Xét đối tượng toán học trong Chương IV: Cho tam giác ABC cân tại A có góc ở đáy B = 73°. Số đo góc ở đáy C bằng. Giá trị hoặc tính chất tương ứng là:\par
\hspace*{1.5em}\textbf{A.} 73°\quad\quad \textbf{B.} 34°\quad\quad \textbf{C.} 107°\quad\quad \textbf{D.} 90°\par\vspace{2pt}
\textbf{\color{primary} Câu 16.} Hoàn thành phép xác định cơ bản sau đây: Cho tam giác ABC có góc A = 80° và góc B = 75°. Số đo góc C của tam giác ABC bằng:\par
\hspace*{1.5em}\textbf{A.} 25°\quad\quad \textbf{B.} 155°\quad\quad \textbf{C.} 35°\quad\quad \textbf{D.} 15°\par\vspace{2pt}
\textbf{\color{primary} Câu 17.} Trong các khẳng định hoặc giá trị dưới đây, đâu là đáp án đúng cho câu hỏi: Cho $\Delta$ABC = $\Delta$DEF. Biết cạnh AB = 21 cm, BC = 23 cm, AC = 24 cm. Độ dài cạnh DE bằng?\par
\hspace*{1.5em}\textbf{A.} 23 cm\quad\quad \textbf{B.} 24 cm\quad\quad \textbf{C.} 68 cm\quad\quad \textbf{D.} 21 cm\par\vspace{2pt}
\textbf{\color{primary} Câu 18.} Bài tập khởi động Chương IV: Cho tam giác ABC cân tại A có góc ở đáy B = 79°. Số đo góc ở đáy C bằng:\par
\hspace*{1.5em}\textbf{A.} 22°\quad\quad \textbf{B.} 101°\quad\quad \textbf{C.} 79°\quad\quad \textbf{D.} 90°\par\vspace{2pt}
\textbf{\color{primary} Câu 19.} Quan sát dữ kiện đề bài: Cho tam giác ABC có góc A = 86° và góc B = 78°. Số đo góc C của tam giác ABC bằng. Hãy chọn phương án đúng:\par
\hspace*{1.5em}\textbf{A.} 26°\quad\quad \textbf{B.} 16°\quad\quad \textbf{C.} 164°\quad\quad \textbf{D.} 6°\par\vspace{2pt}
\textbf{\color{primary} Câu 20.} Sử dụng công thức hoặc tính chất đầu bài của Chương IV để giải quyết yêu cầu: Cho $\Delta$ABC = $\Delta$DEF. Biết cạnh AB = 24 cm, BC = 26 cm, AC = 27 cm. Độ dài cạnh DE bằng:\par
\hspace*{1.5em}\textbf{A.} 24 cm\quad\quad \textbf{B.} 27 cm\quad\quad \textbf{C.} 77 cm\quad\quad \textbf{D.} 26 cm\par\vspace{2pt}
\subsection*{\color{accent} Phần II. Mức độ Trung bình (Thông hiểu -- Câu 21 đến Câu 40)}
\textbf{\color{primary} Câu 21.} Tam giác ABC cân tại A có góc ở đỉnh A = 40°. Số đo mỗi góc ở đáy B và C bằng:\par
\hspace*{1.5em}\textbf{A.} 50°\quad\quad \textbf{B.} 70°\quad\quad \textbf{C.} 140°\quad\quad \textbf{D.} 40°\par\vspace{2pt}
\textbf{\color{primary} Câu 22.} Vận dụng quy tắc trong Chương IV (Tam giác bằng nhau) để giải bài toán: Cho tam giác ABC vuông tại A, biết góc B = 28°. Số đo góc C bằng:\par
\hspace*{1.5em}\textbf{A.} 152°\quad\quad \textbf{B.} 62°\quad\quad \textbf{C.} 28°\quad\quad \textbf{D.} 72°\par\vspace{2pt}
\textbf{\color{primary} Câu 23.} Biến đổi và tính toán theo kiến thức Chương IV: Điểm M nằm trên đường trung trực của đoạn thẳng AB. Biết MA = 11 cm, độ dài đoạn thẳng MB bằng. Kết quả thu được là:\par
\hspace*{1.5em}\textbf{A.} 5.5 cm\quad\quad \textbf{B.} 11 cm\quad\quad \textbf{C.} 22 cm\quad\quad \textbf{D.} 9 cm\par\vspace{2pt}
\textbf{\color{primary} Câu 24.} Xét bài toán mức độ thông hiểu sau: Tam giác ABC cân tại A có góc ở đỉnh A = 52°. Số đo mỗi góc ở đáy B và C bằng. Đáp án đúng là:\par
\hspace*{1.5em}\textbf{A.} 64°\quad\quad \textbf{B.} 52°\quad\quad \textbf{C.} 44°\quad\quad \textbf{D.} 128°\par\vspace{2pt}
\textbf{\color{primary} Câu 25.} Để hoàn thành phiếu học tập Chương IV, hãy xác định: Cho tam giác ABC vuông tại A, biết góc B = 37°. Số đo góc C bằng:\par
\hspace*{1.5em}\textbf{A.} 143°\quad\quad \textbf{B.} 53°\quad\quad \textbf{C.} 63°\quad\quad \textbf{D.} 37°\par\vspace{2pt}
\textbf{\color{primary} Câu 26.} Phân tích dữ kiện bài toán: Điểm M nằm trên đường trung trực của đoạn thẳng AB. Biết MA = 17 cm, độ dài đoạn thẳng MB bằng. Giá trị cần tìm bằng:\par
\hspace*{1.5em}\textbf{A.} 17 cm\quad\quad \textbf{B.} 34 cm\quad\quad \textbf{C.} 15 cm\quad\quad \textbf{D.} 8.5 cm\par\vspace{2pt}
\textbf{\color{primary} Câu 27.} Áp dụng tính chất đã học ở Chương IV vào bài tập sau: Tam giác ABC cân tại A có góc ở đỉnh A = 64°. Số đo mỗi góc ở đáy B và C bằng:\par
\hspace*{1.5em}\textbf{A.} 116°\quad\quad \textbf{B.} 38°\quad\quad \textbf{C.} 58°\quad\quad \textbf{D.} 64°\par\vspace{2pt}
\textbf{\color{primary} Câu 28.} Thực hiện các bước suy luận hợp lý để giải quyết yêu cầu: Cho tam giác ABC vuông tại A, biết góc B = 46°. Số đo góc C bằng:\par
\hspace*{1.5em}\textbf{A.} 44°\quad\quad \textbf{B.} 134°\quad\quad \textbf{C.} 46°\quad\quad \textbf{D.} 54°\par\vspace{2pt}
\textbf{\color{primary} Câu 29.} Cho bài toán tự luyện (Tam giác bằng nhau): Điểm M nằm trên đường trung trực của đoạn thẳng AB. Biết MA = 23 cm, độ dài đoạn thẳng MB bằng. Phương án nào sau đây chính xác?\par
\hspace*{1.5em}\textbf{A.} 11.5 cm\quad\quad \textbf{B.} 21 cm\quad\quad \textbf{C.} 46 cm\quad\quad \textbf{D.} 23 cm\par\vspace{2pt}
\textbf{\color{primary} Câu 30.} Kết hợp các phép biến đổi trong Chương IV, hãy tìm kết quả của bài toán: Tam giác ABC cân tại A có góc ở đỉnh A = 76°. Số đo mỗi góc ở đáy B và C bằng:\par
\hspace*{1.5em}\textbf{A.} 104°\quad\quad \textbf{B.} 52°\quad\quad \textbf{C.} 76°\quad\quad \textbf{D.} 32°\par\vspace{2pt}
\textbf{\color{primary} Câu 31.} Một bài kiểm tra định kỳ nội dung Chương IV nêu câu hỏi: Cho tam giác ABC vuông tại A, biết góc B = 55°. Số đo góc C bằng. Hãy chọn đáp số đúng:\par
\hspace*{1.5em}\textbf{A.} 35°\quad\quad \textbf{B.} 125°\quad\quad \textbf{C.} 45°\quad\quad \textbf{D.} 55°\par\vspace{2pt}
\textbf{\color{primary} Câu 32.} Tìm lời giải đúng cho bài toán thông hiểu: Điểm M nằm trên đường trung trực của đoạn thẳng AB. Biết MA = 29 cm, độ dài đoạn thẳng MB bằng:\par
\hspace*{1.5em}\textbf{A.} 29 cm\quad\quad \textbf{B.} 58 cm\quad\quad \textbf{C.} 14.5 cm\quad\quad \textbf{D.} 27 cm\par\vspace{2pt}
\textbf{\color{primary} Câu 33.} Sử dụng mối liên hệ giữa các đại lượng trong Chương IV: Tam giác ABC cân tại A có góc ở đỉnh A = 88°. Số đo mỗi góc ở đáy B và C bằng:\par
\hspace*{1.5em}\textbf{A.} 92°\quad\quad \textbf{B.} 26°\quad\quad \textbf{C.} 46°\quad\quad \textbf{D.} 88°\par\vspace{2pt}
\textbf{\color{primary} Câu 34.} Hãy tính toán và chọn phương án đúng đối với bài tập: Cho tam giác ABC vuông tại A, biết góc B = 64°. Số đo góc C bằng:\par
\hspace*{1.5em}\textbf{A.} 64°\quad\quad \textbf{B.} 116°\quad\quad \textbf{C.} 36°\quad\quad \textbf{D.} 26°\par\vspace{2pt}
\textbf{\color{primary} Câu 35.} Dựa trên các tính chất trọng tâm của Chương IV, xét bài toán: Điểm M nằm trên đường trung trực của đoạn thẳng AB. Biết MA = 35 cm, độ dài đoạn thẳng MB bằng:\par
\hspace*{1.5em}\textbf{A.} 33 cm\quad\quad \textbf{B.} 70 cm\quad\quad \textbf{C.} 17.5 cm\quad\quad \textbf{D.} 35 cm\par\vspace{2pt}
\textbf{\color{primary} Câu 36.} Thực hành giải toán (Tam giác bằng nhau): Tam giác ABC cân tại A có góc ở đỉnh A = 100°. Số đo mỗi góc ở đáy B và C bằng. Đáp án nào dưới đây thỏa mãn?\par
\hspace*{1.5em}\textbf{A.} 80°\quad\quad \textbf{B.} 100°\quad\quad \textbf{C.} 20°\quad\quad \textbf{D.} 40°\par\vspace{2pt}
\textbf{\color{primary} Câu 37.} Kiểm tra kỹ năng biến đổi và suy luận qua bài toán: Cho tam giác ABC vuông tại A, biết góc B = 73°. Số đo góc C bằng:\par
\hspace*{1.5em}\textbf{A.} 107°\quad\quad \textbf{B.} 27°\quad\quad \textbf{C.} 17°\quad\quad \textbf{D.} 73°\par\vspace{2pt}
\textbf{\color{primary} Câu 38.} Cho giả thiết thuộc phạm vi Chương IV: Điểm M nằm trên đường trung trực của đoạn thẳng AB. Biết MA = 41 cm, độ dài đoạn thẳng MB bằng. Kết luận đúng là:\par
\hspace*{1.5em}\textbf{A.} 39 cm\quad\quad \textbf{B.} 41 cm\quad\quad \textbf{C.} 82 cm\quad\quad \textbf{D.} 20.5 cm\par\vspace{2pt}
\textbf{\color{primary} Câu 39.} Rút gọn hoặc tính giá trị theo yêu cầu sau của Chương IV: Tam giác ABC cân tại A có góc ở đỉnh A = 112°. Số đo mỗi góc ở đáy B và C bằng:\par
\hspace*{1.5em}\textbf{A.} 34°\quad\quad \textbf{B.} 14°\quad\quad \textbf{C.} 112°\quad\quad \textbf{D.} 68°\par\vspace{2pt}
\textbf{\color{primary} Câu 40.} Xác định phương án đúng sau khi thực hiện phép tính / suy luận: Cho tam giác ABC vuông tại A, biết góc B = 82°. Số đo góc C bằng:\par
\hspace*{1.5em}\textbf{A.} 98°\quad\quad \textbf{B.} 8°\quad\quad \textbf{C.} 82°\quad\quad \textbf{D.} 18°\par\vspace{2pt}
\subsection*{\color{accent} Phần III. Mức độ Khó (Vận dụng \& Vận dụng cao -- Câu 41 đến Câu 50)}
\textbf{\color{primary} Câu 41.} Trong tam giác ABC, số đo ba góc A, B, C tỉ lệ với 2 : 3 : 4. Số đo góc lớn nhất của tam giác ABC là:\par
\hspace*{1.5em}\textbf{A.} 60°\quad\quad \textbf{B.} 80°\quad\quad \textbf{C.} 40°\quad\quad \textbf{D.} 100°\par\vspace{2pt}
\textbf{\color{primary} Câu 42.} Bài toán vận dụng (Chương IV – Tam giác bằng nhau): Cho tam giác ABC có góc A = 54°. Các tia phân giác trong của góc B và góc C cắt nhau tại I. Số đo góc BIC bằng. Kết quả đúng là:\par
\hspace*{1.5em}\textbf{A.} 126°\quad\quad \textbf{B.} 102°\quad\quad \textbf{C.} 117°\quad\quad \textbf{D.} 63°\par\vspace{2pt}
\textbf{\color{primary} Câu 43.} Vận dụng kiến thức nâng cao trong phạm vi Chương IV để giải bài toán: Trong tam giác ABC, số đo ba góc A, B, C tỉ lệ với 2 : 3 : 4. Số đo góc lớn nhất của tam giác ABC là:\par
\hspace*{1.5em}\textbf{A.} 60°\quad\quad \textbf{B.} 100°\quad\quad \textbf{C.} 40°\quad\quad \textbf{D.} 80°\par\vspace{2pt}
\textbf{\color{primary} Câu 44.} Phân tích và giải bài toán tổng hợp sau của Chương IV: Cho tam giác ABC có góc A = 62°. Các tia phân giác trong của góc B và góc C cắt nhau tại I. Số đo góc BIC bằng:\par
\hspace*{1.5em}\textbf{A.} 118°\quad\quad \textbf{B.} 106°\quad\quad \textbf{C.} 59°\quad\quad \textbf{D.} 121°\par\vspace{2pt}
\textbf{\color{primary} Câu 45.} Thử thách tư duy trong phạm vi kiến thức Chương IV: Trong tam giác ABC, số đo ba góc A, B, C tỉ lệ với 2 : 3 : 4. Số đo góc lớn nhất của tam giác ABC là. Đáp số chính xác là:\par
\hspace*{1.5em}\textbf{A.} 60°\quad\quad \textbf{B.} 100°\quad\quad \textbf{C.} 80°\quad\quad \textbf{D.} 40°\par\vspace{2pt}
\textbf{\color{primary} Câu 46.} Sử dụng phương pháp giải toán của Chương IV để tìm đáp án cho bài toán: Cho tam giác ABC có góc A = 70°. Các tia phân giác trong của góc B và góc C cắt nhau tại I. Số đo góc BIC bằng:\par
\hspace*{1.5em}\textbf{A.} 55°\quad\quad \textbf{B.} 125°\quad\quad \textbf{C.} 110°\quad\quad \textbf{D.} 125° (1)\par\vspace{2pt}
\textbf{\color{primary} Câu 47.} Xét bài toán vận dụng thực tiễn / nâng cao thuộc Chương IV: Trong tam giác ABC, số đo ba góc A, B, C tỉ lệ với 2 : 3 : 4. Số đo góc lớn nhất của tam giác ABC là:\par
\hspace*{1.5em}\textbf{A.} 100°\quad\quad \textbf{B.} 60°\quad\quad \textbf{C.} 40°\quad\quad \textbf{D.} 80°\par\vspace{2pt}
\textbf{\color{primary} Câu 48.} Tìm kết quả cuối cùng của bài toán phân loại học sinh khá giỏi (Chương IV): Cho tam giác ABC có góc A = 78°. Các tia phân giác trong của góc B và góc C cắt nhau tại I. Số đo góc BIC bằng:\par
\hspace*{1.5em}\textbf{A.} 114°\quad\quad \textbf{B.} 129°\quad\quad \textbf{C.} 51°\quad\quad \textbf{D.} 102°\par\vspace{2pt}
\textbf{\color{primary} Câu 49.} Kết hợp các tính chất đã học trong Chương IV, hãy giải quyết bài toán: Trong tam giác ABC, số đo ba góc A, B, C tỉ lệ với 2 : 3 : 4. Số đo góc lớn nhất của tam giác ABC là:\par
\hspace*{1.5em}\textbf{A.} 100°\quad\quad \textbf{B.} 60°\quad\quad \textbf{C.} 80°\quad\quad \textbf{D.} 40°\par\vspace{2pt}
\textbf{\color{primary} Câu 50.} Câu hỏi vận dụng cao cuối Chương IV: Cho tam giác ABC có góc A = 86°. Các tia phân giác trong của góc B và góc C cắt nhau tại I. Số đo góc BIC bằng. Phương án đúng là:\par
\hspace*{1.5em}\textbf{A.} 118°\quad\quad \textbf{B.} 47°\quad\quad \textbf{C.} 94°\quad\quad \textbf{D.} 133°\par\vspace{2pt}
\subsection*{\color{primary} Bảng đáp án Chương IV: Tam giác bằng nhau}
\begin{center}
\renewcommand{\arraystretch}{1.25}
\begin{tabularx}{\textwidth}{*{10}{>{\centering\arraybackslash}X}}
\toprule
\textbf{\color{primary} Cột 1} & \textbf{\color{primary} Cột 2} & \textbf{\color{primary} Cột 3} & \textbf{\color{primary} Cột 4} & \textbf{\color{primary} Cột 5} & \textbf{\color{primary} Cột 6} & \textbf{\color{primary} Cột 7} & \textbf{\color{primary} Cột 8} & \textbf{\color{primary} Cột 9} & \textbf{\color{primary} Cột 10} \\
\midrule
\textbf{1.B} & \textbf{2.A} & \textbf{3.D} & \textbf{4.A} & \textbf{5.A} & \textbf{6.B} & \textbf{7.D} & \textbf{8.B} & \textbf{9.C} & \textbf{10.B} \\
\rowcolor{rowtint} \textbf{11.B} & \textbf{12.B} & \textbf{13.C} & \textbf{14.A} & \textbf{15.A} & \textbf{16.A} & \textbf{17.D} & \textbf{18.C} & \textbf{19.B} & \textbf{20.A} \\
\textbf{21.B} & \textbf{22.B} & \textbf{23.B} & \textbf{24.A} & \textbf{25.B} & \textbf{26.A} & \textbf{27.C} & \textbf{28.A} & \textbf{29.D} & \textbf{30.B} \\
\rowcolor{rowtint} \textbf{31.A} & \textbf{32.A} & \textbf{33.C} & \textbf{34.D} & \textbf{35.D} & \textbf{36.D} & \textbf{37.C} & \textbf{38.B} & \textbf{39.A} & \textbf{40.B} \\
\textbf{41.B} & \textbf{42.C} & \textbf{43.D} & \textbf{44.D} & \textbf{45.C} & \textbf{46.B} & \textbf{47.D} & \textbf{48.B} & \textbf{49.C} & \textbf{50.D} \\
\bottomrule
\end{tabularx}
\end{center}\vspace{10pt}
\section*{\color{primary} CHƯƠNG V: THU THẬP VÀ BIỂU DIỄN DỮ LIỆU}
\textit{Trọng tâm kiến thức SGK: Thu thập và phân loại dữ liệu (dữ liệu định tính, định lượng, tính đại diện), biểu đồ hình quạt tròn và biểu đồ đoạn thẳng.}\par\vspace{4pt}
\subsection*{\color{accent} Phần I. Mức độ Dễ (Nhận biết -- Câu 1 đến Câu 20)}
\textbf{\color{primary} Câu 1.} Trong biểu đồ hình quạt tròn, tổng các tỉ lệ phần trăm của tất cả các thành phần (cả hình tròn) luôn bằng:\par
\hspace*{1.5em}\textbf{A.} 90\%\quad\quad \textbf{B.} 360\%\quad\quad \textbf{C.} 50\%\quad\quad \textbf{D.} 100\%\par\vspace{2pt}
\textbf{\color{primary} Câu 2.} Hãy trả lời câu hỏi cơ bản sau: Biểu đồ nào thích hợp nhất để biểu diễn sự thay đổi của một đại lượng (như nhiệt độ, doanh thu) theo thời gian?\par
\hspace*{1.5em}\textbf{A.} Bảng kiểm đếm\quad\quad \textbf{B.} Biểu đồ hình quạt tròn\quad\quad \textbf{C.} Biểu đồ tranh\quad\quad \textbf{D.} Biểu đồ đoạn thẳng\par\vspace{2pt}
\textbf{\color{primary} Câu 3.} Trong bài học thuộc Chương V, xét nội dung: Trong biểu đồ hình quạt tròn thống kê xếp loại học lực của lớp 7A: Tốt chiếm 22\%, Khá chiếm 45\%, Đạt chiếm 28\%. Tỉ lệ phần trăm học sinh Chưa đạt là. Đáp án đúng là:\par
\hspace*{1.5em}\textbf{A.} 0\%\quad\quad \textbf{B.} 5\%\quad\quad \textbf{C.} 10\%\quad\quad \textbf{D.} 15\%\par\vspace{2pt}
\textbf{\color{primary} Câu 4.} Dựa vào kiến thức chuẩn SGK của Chương V (Thu thập và biểu diễn dữ liệu), trong biểu đồ hình quạt tròn, tổng các tỉ lệ phần trăm của tất cả các thành phần (cả hình tròn) luôn bằng:\par
\hspace*{1.5em}\textbf{A.} 90\%\quad\quad \textbf{B.} 100\%\quad\quad \textbf{C.} 50\%\quad\quad \textbf{D.} 360\%\par\vspace{2pt}
\textbf{\color{primary} Câu 5.} Chọn phương án trả lời đúng nhất: Biểu đồ nào thích hợp nhất để biểu diễn sự thay đổi của một đại lượng (như nhiệt độ, doanh thu) theo thời gian?\par
\hspace*{1.5em}\textbf{A.} Biểu đồ đoạn thẳng\quad\quad \textbf{B.} Biểu đồ hình quạt tròn\quad\quad \textbf{C.} Biểu đồ tranh\quad\quad \textbf{D.} Bảng kiểm đếm\par\vspace{2pt}
\textbf{\color{primary} Câu 6.} Khi làm bài tập cơ bản về thu thập và biểu diễn dữ liệu, một học sinh gặp câu hỏi: Trong biểu đồ hình quạt tròn thống kê xếp loại học lực của lớp 7A: Tốt chiếm 25\%, Khá chiếm 45\%, Đạt chiếm 25\%. Tỉ lệ phần trăm học sinh Chưa đạt là. Kết quả nào dưới đây đúng?\par
\hspace*{1.5em}\textbf{A.} 10\%\quad\quad \textbf{B.} 15\%\quad\quad \textbf{C.} 5\%\quad\quad \textbf{D.} 0\%\par\vspace{2pt}
\textbf{\color{primary} Câu 7.} Xác định giá trị hoặc khẳng định đúng cho bài toán sau: Trong biểu đồ hình quạt tròn, tổng các tỉ lệ phần trăm của tất cả các thành phần (cả hình tròn) luôn bằng:\par
\hspace*{1.5em}\textbf{A.} 90\%\quad\quad \textbf{B.} 360\%\quad\quad \textbf{C.} 100\%\quad\quad \textbf{D.} 50\%\par\vspace{2pt}
\textbf{\color{primary} Câu 8.} Áp dụng trực tiếp định nghĩa và quy tắc trong Chương V: Biểu đồ nào thích hợp nhất để biểu diễn sự thay đổi của một đại lượng (như nhiệt độ, doanh thu) theo thời gian:\par
\hspace*{1.5em}\textbf{A.} Biểu đồ đoạn thẳng\quad\quad \textbf{B.} Biểu đồ hình quạt tròn\quad\quad \textbf{C.} Biểu đồ tranh\quad\quad \textbf{D.} Bảng kiểm đếm\par\vspace{2pt}
\textbf{\color{primary} Câu 9.} Kiểm tra nhanh kiến thức nền tảng: Trong biểu đồ hình quạt tròn thống kê xếp loại học lực của lớp 7A: Tốt chiếm 28\%, Khá chiếm 45\%, Đạt chiếm 22\%. Tỉ lệ phần trăm học sinh Chưa đạt là:\par
\hspace*{1.5em}\textbf{A.} 0\%\quad\quad \textbf{B.} 5\%\quad\quad \textbf{C.} 15\%\quad\quad \textbf{D.} 10\%\par\vspace{2pt}
\textbf{\color{primary} Câu 10.} Cho dữ kiện nhận biết sau đây: Trong biểu đồ hình quạt tròn, tổng các tỉ lệ phần trăm của tất cả các thành phần (cả hình tròn) luôn bằng. Phương án chính xác là:\par
\hspace*{1.5em}\textbf{A.} 50\%\quad\quad \textbf{B.} 100\%\quad\quad \textbf{C.} 360\%\quad\quad \textbf{D.} 90\%\par\vspace{2pt}
\textbf{\color{primary} Câu 11.} Hãy chỉ ra kết quả đúng trong bốn lựa chọn A, B, C, D đối với yêu cầu: Biểu đồ nào thích hợp nhất để biểu diễn sự thay đổi của một đại lượng (như nhiệt độ, doanh thu) theo thời gian:\par
\hspace*{1.5em}\textbf{A.} Biểu đồ đoạn thẳng\quad\quad \textbf{B.} Bảng kiểm đếm\quad\quad \textbf{C.} Biểu đồ hình quạt tròn\quad\quad \textbf{D.} Biểu đồ tranh\par\vspace{2pt}
\textbf{\color{primary} Câu 12.} Câu hỏi ôn tập lý thuyết và kỹ năng cơ bản (Chương V): Trong biểu đồ hình quạt tròn thống kê xếp loại học lực của lớp 7A: Tốt chiếm 31\%, Khá chiếm 45\%, Đạt chiếm 19\%. Tỉ lệ phần trăm học sinh Chưa đạt là:\par
\hspace*{1.5em}\textbf{A.} 15\%\quad\quad \textbf{B.} 0\%\quad\quad \textbf{C.} 5\%\quad\quad \textbf{D.} 10\%\par\vspace{2pt}
\textbf{\color{primary} Câu 13.} Theo nội dung sách giáo khoa Kết nối tri thức với cuộc sống (Chương V), trong biểu đồ hình quạt tròn, tổng các tỉ lệ phần trăm của tất cả các thành phần (cả hình tròn) luôn bằng:\par
\hspace*{1.5em}\textbf{A.} 360\%\quad\quad \textbf{B.} 90\%\quad\quad \textbf{C.} 100\%\quad\quad \textbf{D.} 50\%\par\vspace{2pt}
\textbf{\color{primary} Câu 14.} Tìm đáp số đúng của bài toán cơ bản sau: Biểu đồ nào thích hợp nhất để biểu diễn sự thay đổi của một đại lượng (như nhiệt độ, doanh thu) theo thời gian:\par
\hspace*{1.5em}\textbf{A.} Bảng kiểm đếm\quad\quad \textbf{B.} Biểu đồ tranh\quad\quad \textbf{C.} Biểu đồ đoạn thẳng\quad\quad \textbf{D.} Biểu đồ hình quạt tròn\par\vspace{2pt}
\textbf{\color{primary} Câu 15.} Xét đối tượng toán học trong Chương V: Trong biểu đồ hình quạt tròn thống kê xếp loại học lực của lớp 7A: Tốt chiếm 34\%, Khá chiếm 45\%, Đạt chiếm 16\%. Tỉ lệ phần trăm học sinh Chưa đạt là. Giá trị hoặc tính chất tương ứng là:\par
\hspace*{1.5em}\textbf{A.} 0\%\quad\quad \textbf{B.} 15\%\quad\quad \textbf{C.} 5\%\quad\quad \textbf{D.} 10\%\par\vspace{2pt}
\textbf{\color{primary} Câu 16.} Hoàn thành phép xác định cơ bản sau đây: Trong biểu đồ hình quạt tròn, tổng các tỉ lệ phần trăm của tất cả các thành phần (cả hình tròn) luôn bằng:\par
\hspace*{1.5em}\textbf{A.} 100\%\quad\quad \textbf{B.} 50\%\quad\quad \textbf{C.} 90\%\quad\quad \textbf{D.} 360\%\par\vspace{2pt}
\textbf{\color{primary} Câu 17.} Trong các khẳng định hoặc giá trị dưới đây, đâu là đáp án đúng cho câu hỏi: Biểu đồ nào thích hợp nhất để biểu diễn sự thay đổi của một đại lượng (như nhiệt độ, doanh thu) theo thời gian?\par
\hspace*{1.5em}\textbf{A.} Biểu đồ hình quạt tròn\quad\quad \textbf{B.} Biểu đồ tranh\quad\quad \textbf{C.} Biểu đồ đoạn thẳng\quad\quad \textbf{D.} Bảng kiểm đếm\par\vspace{2pt}
\textbf{\color{primary} Câu 18.} Bài tập khởi động Chương V: Trong biểu đồ hình quạt tròn thống kê xếp loại học lực của lớp 7A: Tốt chiếm 37\%, Khá chiếm 45\%, Đạt chiếm 13\%. Tỉ lệ phần trăm học sinh Chưa đạt là:\par
\hspace*{1.5em}\textbf{A.} 15\%\quad\quad \textbf{B.} 10\%\quad\quad \textbf{C.} 5\%\quad\quad \textbf{D.} 0\%\par\vspace{2pt}
\textbf{\color{primary} Câu 19.} Quan sát dữ kiện đề bài: Trong biểu đồ hình quạt tròn, tổng các tỉ lệ phần trăm của tất cả các thành phần (cả hình tròn) luôn bằng. Hãy chọn phương án đúng:\par
\hspace*{1.5em}\textbf{A.} 100\%\quad\quad \textbf{B.} 90\%\quad\quad \textbf{C.} 50\%\quad\quad \textbf{D.} 360\%\par\vspace{2pt}
\textbf{\color{primary} Câu 20.} Sử dụng công thức hoặc tính chất đầu bài của Chương V để giải quyết yêu cầu: Biểu đồ nào thích hợp nhất để biểu diễn sự thay đổi của một đại lượng (như nhiệt độ, doanh thu) theo thời gian:\par
\hspace*{1.5em}\textbf{A.} Bảng kiểm đếm\quad\quad \textbf{B.} Biểu đồ đoạn thẳng\quad\quad \textbf{C.} Biểu đồ tranh\quad\quad \textbf{D.} Biểu đồ hình quạt tròn\par\vspace{2pt}
\subsection*{\color{accent} Phần II. Mức độ Trung bình (Thông hiểu -- Câu 21 đến Câu 40)}
\textbf{\color{primary} Câu 21.} Một trường THCS có 400 học sinh khối 7. Biết trên biểu đồ hình quạt tròn, số học sinh đăng ký câu lạc bộ Toán chiếm 25\%. Số học sinh đăng ký câu lạc bộ Toán là:\par
\hspace*{1.5em}\textbf{A.} 80 học sinh\quad\quad \textbf{B.} 300 học sinh\quad\quad \textbf{C.} 100 học sinh\quad\quad \textbf{D.} 200 học sinh\par\vspace{2pt}
\textbf{\color{primary} Câu 22.} Vận dụng quy tắc trong Chương V (Thu thập và biểu diễn dữ liệu) để giải bài toán: Doanh thu của một cửa hàng sách trong 3 tháng đầu năm lần lượt là 45, 55 và 65 triệu đồng. Doanh thu trung bình mỗi tháng trong quý I là:\par
\hspace*{1.5em}\textbf{A.} 55 triệu đồng\quad\quad \textbf{B.} 60 triệu đồng\quad\quad \textbf{C.} 165 triệu đồng\quad\quad \textbf{D.} 50 triệu đồng\par\vspace{2pt}
\textbf{\color{primary} Câu 23.} Biến đổi và tính toán theo kiến thức Chương V: Muốn khảo sát 'Thời gian tự học ở nhà mỗi ngày của học sinh khối 7 trong trường', cách lấy mẫu nào sau đây đảm bảo tính đại diện tốt nhất. Kết quả thu được là:\par
\begin{itemize}\setlength{\itemsep}{1pt}\setlength{\parskip}{0pt}
  \item[] \textbf{A.} Chỉ hỏi các bạn nam trong câu lạc bộ bóng đá
  \item[] \textbf{B.} Chỉ hỏi các bạn ngồi bàn đầu của lớp 7A1
  \item[] \textbf{C.} Chỉ hỏi các bạn trong đội tuyển học sinh giỏi Toán khối 7
  \item[] \textbf{D.} Chọn ngẫu nhiên 10 học sinh từ mỗi lớp 7 trong toàn khối
\end{itemize}\vspace{2pt}
\textbf{\color{primary} Câu 24.} Xét bài toán mức độ thông hiểu sau: Một trường THCS có 700 học sinh khối 7. Biết trên biểu đồ hình quạt tròn, số học sinh đăng ký câu lạc bộ Toán chiếm 25\%. Số học sinh đăng ký câu lạc bộ Toán là. Đáp án đúng là:\par
\hspace*{1.5em}\textbf{A.} 175 học sinh\quad\quad \textbf{B.} 350 học sinh\quad\quad \textbf{C.} 140 học sinh\quad\quad \textbf{D.} 525 học sinh\par\vspace{2pt}
\textbf{\color{primary} Câu 25.} Để hoàn thành phiếu học tập Chương V, hãy xác định: Doanh thu của một cửa hàng sách trong 3 tháng đầu năm lần lượt là 60, 70 và 80 triệu đồng. Doanh thu trung bình mỗi tháng trong quý I là:\par
\hspace*{1.5em}\textbf{A.} 75 triệu đồng\quad\quad \textbf{B.} 65 triệu đồng\quad\quad \textbf{C.} 210 triệu đồng\quad\quad \textbf{D.} 70 triệu đồng\par\vspace{2pt}
\textbf{\color{primary} Câu 26.} Phân tích dữ kiện bài toán: Muốn khảo sát 'Thời gian tự học ở nhà mỗi ngày của học sinh khối 7 trong trường', cách lấy mẫu nào sau đây đảm bảo tính đại diện tốt nhất. Giá trị cần tìm bằng:\par
\begin{itemize}\setlength{\itemsep}{1pt}\setlength{\parskip}{0pt}
  \item[] \textbf{A.} Chỉ hỏi các bạn nam trong câu lạc bộ bóng đá
  \item[] \textbf{B.} Chỉ hỏi các bạn ngồi bàn đầu của lớp 7A1
  \item[] \textbf{C.} Chỉ hỏi các bạn trong đội tuyển học sinh giỏi Toán khối 7
  \item[] \textbf{D.} Chọn ngẫu nhiên 10 học sinh từ mỗi lớp 7 trong toàn khối
\end{itemize}\vspace{2pt}
\textbf{\color{primary} Câu 27.} Áp dụng tính chất đã học ở Chương V vào bài tập sau: Một trường THCS có 500 học sinh khối 7. Biết trên biểu đồ hình quạt tròn, số học sinh đăng ký câu lạc bộ Toán chiếm 25\%. Số học sinh đăng ký câu lạc bộ Toán là:\par
\hspace*{1.5em}\textbf{A.} 250 học sinh\quad\quad \textbf{B.} 100 học sinh\quad\quad \textbf{C.} 125 học sinh\quad\quad \textbf{D.} 375 học sinh\par\vspace{2pt}
\textbf{\color{primary} Câu 28.} Thực hiện các bước suy luận hợp lý để giải quyết yêu cầu: Doanh thu của một cửa hàng sách trong 3 tháng đầu năm lần lượt là 75, 85 và 95 triệu đồng. Doanh thu trung bình mỗi tháng trong quý I là:\par
\hspace*{1.5em}\textbf{A.} 80 triệu đồng\quad\quad \textbf{B.} 255 triệu đồng\quad\quad \textbf{C.} 90 triệu đồng\quad\quad \textbf{D.} 85 triệu đồng\par\vspace{2pt}
\textbf{\color{primary} Câu 29.} Cho bài toán tự luyện (Thu thập và biểu diễn dữ liệu): Muốn khảo sát 'Thời gian tự học ở nhà mỗi ngày của học sinh khối 7 trong trường', cách lấy mẫu nào sau đây đảm bảo tính đại diện tốt nhất. Phương án nào sau đây chính xác?\par
\begin{itemize}\setlength{\itemsep}{1pt}\setlength{\parskip}{0pt}
  \item[] \textbf{A.} Chỉ hỏi các bạn ngồi bàn đầu của lớp 7A1
  \item[] \textbf{B.} Chỉ hỏi các bạn trong đội tuyển học sinh giỏi Toán khối 7
  \item[] \textbf{C.} Chọn ngẫu nhiên 10 học sinh từ mỗi lớp 7 trong toàn khối
  \item[] \textbf{D.} Chỉ hỏi các bạn nam trong câu lạc bộ bóng đá
\end{itemize}\vspace{2pt}
\textbf{\color{primary} Câu 30.} Kết hợp các phép biến đổi trong Chương V, hãy tìm kết quả của bài toán: Một trường THCS có 800 học sinh khối 7. Biết trên biểu đồ hình quạt tròn, số học sinh đăng ký câu lạc bộ Toán chiếm 25\%. Số học sinh đăng ký câu lạc bộ Toán là:\par
\hspace*{1.5em}\textbf{A.} 600 học sinh\quad\quad \textbf{B.} 200 học sinh\quad\quad \textbf{C.} 160 học sinh\quad\quad \textbf{D.} 400 học sinh\par\vspace{2pt}
\textbf{\color{primary} Câu 31.} Một bài kiểm tra định kỳ nội dung Chương V nêu câu hỏi: Doanh thu của một cửa hàng sách trong 3 tháng đầu năm lần lượt là 90, 100 và 110 triệu đồng. Doanh thu trung bình mỗi tháng trong quý I là. Hãy chọn đáp số đúng:\par
\hspace*{1.5em}\textbf{A.} 105 triệu đồng\quad\quad \textbf{B.} 300 triệu đồng\quad\quad \textbf{C.} 95 triệu đồng\quad\quad \textbf{D.} 100 triệu đồng\par\vspace{2pt}
\textbf{\color{primary} Câu 32.} Tìm lời giải đúng cho bài toán thông hiểu: Muốn khảo sát 'Thời gian tự học ở nhà mỗi ngày của học sinh khối 7 trong trường', cách lấy mẫu nào sau đây đảm bảo tính đại diện tốt nhất:\par
\begin{itemize}\setlength{\itemsep}{1pt}\setlength{\parskip}{0pt}
  \item[] \textbf{A.} Chỉ hỏi các bạn nam trong câu lạc bộ bóng đá
  \item[] \textbf{B.} Chọn ngẫu nhiên 10 học sinh từ mỗi lớp 7 trong toàn khối
  \item[] \textbf{C.} Chỉ hỏi các bạn trong đội tuyển học sinh giỏi Toán khối 7
  \item[] \textbf{D.} Chỉ hỏi các bạn ngồi bàn đầu của lớp 7A1
\end{itemize}\vspace{2pt}
\textbf{\color{primary} Câu 33.} Sử dụng mối liên hệ giữa các đại lượng trong Chương V: Một trường THCS có 600 học sinh khối 7. Biết trên biểu đồ hình quạt tròn, số học sinh đăng ký câu lạc bộ Toán chiếm 25\%. Số học sinh đăng ký câu lạc bộ Toán là:\par
\hspace*{1.5em}\textbf{A.} 300 học sinh\quad\quad \textbf{B.} 120 học sinh\quad\quad \textbf{C.} 450 học sinh\quad\quad \textbf{D.} 150 học sinh\par\vspace{2pt}
\textbf{\color{primary} Câu 34.} Hãy tính toán và chọn phương án đúng đối với bài tập: Doanh thu của một cửa hàng sách trong 3 tháng đầu năm lần lượt là 105, 115 và 125 triệu đồng. Doanh thu trung bình mỗi tháng trong quý I là:\par
\hspace*{1.5em}\textbf{A.} 120 triệu đồng\quad\quad \textbf{B.} 110 triệu đồng\quad\quad \textbf{C.} 345 triệu đồng\quad\quad \textbf{D.} 115 triệu đồng\par\vspace{2pt}
\textbf{\color{primary} Câu 35.} Dựa trên các tính chất trọng tâm của Chương V, xét bài toán: Muốn khảo sát 'Thời gian tự học ở nhà mỗi ngày của học sinh khối 7 trong trường', cách lấy mẫu nào sau đây đảm bảo tính đại diện tốt nhất:\par
\begin{itemize}\setlength{\itemsep}{1pt}\setlength{\parskip}{0pt}
  \item[] \textbf{A.} Chọn ngẫu nhiên 10 học sinh từ mỗi lớp 7 trong toàn khối
  \item[] \textbf{B.} Chỉ hỏi các bạn trong đội tuyển học sinh giỏi Toán khối 7
  \item[] \textbf{C.} Chỉ hỏi các bạn ngồi bàn đầu của lớp 7A1
  \item[] \textbf{D.} Chỉ hỏi các bạn nam trong câu lạc bộ bóng đá
\end{itemize}\vspace{2pt}
\textbf{\color{primary} Câu 36.} Thực hành giải toán (Thu thập và biểu diễn dữ liệu): Một trường THCS có 400 học sinh khối 7. Biết trên biểu đồ hình quạt tròn, số học sinh đăng ký câu lạc bộ Toán chiếm 25\%. Số học sinh đăng ký câu lạc bộ Toán là. Đáp án nào dưới đây thỏa mãn?\par
\hspace*{1.5em}\textbf{A.} 200 học sinh\quad\quad \textbf{B.} 100 học sinh\quad\quad \textbf{C.} 300 học sinh\quad\quad \textbf{D.} 80 học sinh\par\vspace{2pt}
\textbf{\color{primary} Câu 37.} Kiểm tra kỹ năng biến đổi và suy luận qua bài toán: Doanh thu của một cửa hàng sách trong 3 tháng đầu năm lần lượt là 120, 130 và 140 triệu đồng. Doanh thu trung bình mỗi tháng trong quý I là:\par
\hspace*{1.5em}\textbf{A.} 135 triệu đồng\quad\quad \textbf{B.} 130 triệu đồng\quad\quad \textbf{C.} 125 triệu đồng\quad\quad \textbf{D.} 390 triệu đồng\par\vspace{2pt}
\textbf{\color{primary} Câu 38.} Cho giả thiết thuộc phạm vi Chương V: Muốn khảo sát 'Thời gian tự học ở nhà mỗi ngày của học sinh khối 7 trong trường', cách lấy mẫu nào sau đây đảm bảo tính đại diện tốt nhất. Kết luận đúng là:\par
\begin{itemize}\setlength{\itemsep}{1pt}\setlength{\parskip}{0pt}
  \item[] \textbf{A.} Chỉ hỏi các bạn ngồi bàn đầu của lớp 7A1
  \item[] \textbf{B.} Chỉ hỏi các bạn nam trong câu lạc bộ bóng đá
  \item[] \textbf{C.} Chọn ngẫu nhiên 10 học sinh từ mỗi lớp 7 trong toàn khối
  \item[] \textbf{D.} Chỉ hỏi các bạn trong đội tuyển học sinh giỏi Toán khối 7
\end{itemize}\vspace{2pt}
\textbf{\color{primary} Câu 39.} Rút gọn hoặc tính giá trị theo yêu cầu sau của Chương V: Một trường THCS có 700 học sinh khối 7. Biết trên biểu đồ hình quạt tròn, số học sinh đăng ký câu lạc bộ Toán chiếm 25\%. Số học sinh đăng ký câu lạc bộ Toán là:\par
\hspace*{1.5em}\textbf{A.} 140 học sinh\quad\quad \textbf{B.} 175 học sinh\quad\quad \textbf{C.} 350 học sinh\quad\quad \textbf{D.} 525 học sinh\par\vspace{2pt}
\textbf{\color{primary} Câu 40.} Xác định phương án đúng sau khi thực hiện phép tính / suy luận: Doanh thu của một cửa hàng sách trong 3 tháng đầu năm lần lượt là 135, 145 và 155 triệu đồng. Doanh thu trung bình mỗi tháng trong quý I là:\par
\hspace*{1.5em}\textbf{A.} 435 triệu đồng\quad\quad \textbf{B.} 140 triệu đồng\quad\quad \textbf{C.} 145 triệu đồng\quad\quad \textbf{D.} 150 triệu đồng\par\vspace{2pt}
\subsection*{\color{accent} Phần III. Mức độ Khó (Vận dụng \& Vận dụng cao -- Câu 41 đến Câu 50)}
\textbf{\color{primary} Câu 41.} Trên biểu đồ hình quạt tròn, hình quạt biểu diễn tỉ lệ 15\% tương ứng với góc ở tâm có số đo bằng bao nhiêu độ (trên tổng 360°)?\par
\hspace*{1.5em}\textbf{A.} 27°\quad\quad \textbf{B.} 54°\quad\quad \textbf{C.} 72°\quad\quad \textbf{D.} 30°\par\vspace{2pt}
\textbf{\color{primary} Câu 42.} Bài toán vận dụng (Chương V – Thu thập và biểu diễn dữ liệu): Số lượng sản phẩm bán ra của một cửa hàng tăng từ 220 sản phẩm trong tháng 1 lên 275 sản phẩm trong tháng 2. Tỉ lệ phần trăm tăng trưởng của tháng 2 so với tháng 1 là. Kết quả đúng là:\par
\hspace*{1.5em}\textbf{A.} 20\%\quad\quad \textbf{B.} 50\%\quad\quad \textbf{C.} 25\%\quad\quad \textbf{D.} 15\%\par\vspace{2pt}
\textbf{\color{primary} Câu 43.} Vận dụng kiến thức nâng cao trong phạm vi Chương V để giải bài toán: Trên biểu đồ hình quạt tròn, hình quạt biểu diễn tỉ lệ 25\% tương ứng với góc ở tâm có số đo bằng bao nhiêu độ (trên tổng 360°):\par
\hspace*{1.5em}\textbf{A.} 50°\quad\quad \textbf{B.} 90°\quad\quad \textbf{C.} 108°\quad\quad \textbf{D.} 45°\par\vspace{2pt}
\textbf{\color{primary} Câu 44.} Phân tích và giải bài toán tổng hợp sau của Chương V: Số lượng sản phẩm bán ra của một cửa hàng tăng từ 260 sản phẩm trong tháng 1 lên 325 sản phẩm trong tháng 2. Tỉ lệ phần trăm tăng trưởng của tháng 2 so với tháng 1 là:\par
\hspace*{1.5em}\textbf{A.} 50\%\quad\quad \textbf{B.} 15\%\quad\quad \textbf{C.} 25\%\quad\quad \textbf{D.} 20\%\par\vspace{2pt}
\textbf{\color{primary} Câu 45.} Thử thách tư duy trong phạm vi kiến thức Chương V: Trên biểu đồ hình quạt tròn, hình quạt biểu diễn tỉ lệ 15\% tương ứng với góc ở tâm có số đo bằng bao nhiêu độ (trên tổng 360°). Đáp số chính xác là:\par
\hspace*{1.5em}\textbf{A.} 54°\quad\quad \textbf{B.} 27°\quad\quad \textbf{C.} 72°\quad\quad \textbf{D.} 30°\par\vspace{2pt}
\textbf{\color{primary} Câu 46.} Sử dụng phương pháp giải toán của Chương V để tìm đáp án cho bài toán: Số lượng sản phẩm bán ra của một cửa hàng tăng từ 300 sản phẩm trong tháng 1 lên 375 sản phẩm trong tháng 2. Tỉ lệ phần trăm tăng trưởng của tháng 2 so với tháng 1 là:\par
\hspace*{1.5em}\textbf{A.} 15\%\quad\quad \textbf{B.} 50\%\quad\quad \textbf{C.} 25\%\quad\quad \textbf{D.} 20\%\par\vspace{2pt}
\textbf{\color{primary} Câu 47.} Xét bài toán vận dụng thực tiễn / nâng cao thuộc Chương V: Trên biểu đồ hình quạt tròn, hình quạt biểu diễn tỉ lệ 25\% tương ứng với góc ở tâm có số đo bằng bao nhiêu độ (trên tổng 360°):\par
\hspace*{1.5em}\textbf{A.} 108°\quad\quad \textbf{B.} 50°\quad\quad \textbf{C.} 90°\quad\quad \textbf{D.} 45°\par\vspace{2pt}
\textbf{\color{primary} Câu 48.} Tìm kết quả cuối cùng của bài toán phân loại học sinh khá giỏi (Chương V): Số lượng sản phẩm bán ra của một cửa hàng tăng từ 340 sản phẩm trong tháng 1 lên 425 sản phẩm trong tháng 2. Tỉ lệ phần trăm tăng trưởng của tháng 2 so với tháng 1 là:\par
\hspace*{1.5em}\textbf{A.} 15\%\quad\quad \textbf{B.} 25\%\quad\quad \textbf{C.} 20\%\quad\quad \textbf{D.} 50\%\par\vspace{2pt}
\textbf{\color{primary} Câu 49.} Kết hợp các tính chất đã học trong Chương V, hãy giải quyết bài toán: Trên biểu đồ hình quạt tròn, hình quạt biểu diễn tỉ lệ 15\% tương ứng với góc ở tâm có số đo bằng bao nhiêu độ (trên tổng 360°):\par
\hspace*{1.5em}\textbf{A.} 54°\quad\quad \textbf{B.} 27°\quad\quad \textbf{C.} 72°\quad\quad \textbf{D.} 30°\par\vspace{2pt}
\textbf{\color{primary} Câu 50.} Câu hỏi vận dụng cao cuối Chương V: Số lượng sản phẩm bán ra của một cửa hàng tăng từ 380 sản phẩm trong tháng 1 lên 475 sản phẩm trong tháng 2. Tỉ lệ phần trăm tăng trưởng của tháng 2 so với tháng 1 là. Phương án đúng là:\par
\hspace*{1.5em}\textbf{A.} 25\%\quad\quad \textbf{B.} 20\%\quad\quad \textbf{C.} 50\%\quad\quad \textbf{D.} 15\%\par\vspace{2pt}
\subsection*{\color{primary} Bảng đáp án Chương V: Thu thập và biểu diễn dữ liệu}
\begin{center}
\renewcommand{\arraystretch}{1.25}
\begin{tabularx}{\textwidth}{*{10}{>{\centering\arraybackslash}X}}
\toprule
\textbf{\color{primary} Cột 1} & \textbf{\color{primary} Cột 2} & \textbf{\color{primary} Cột 3} & \textbf{\color{primary} Cột 4} & \textbf{\color{primary} Cột 5} & \textbf{\color{primary} Cột 6} & \textbf{\color{primary} Cột 7} & \textbf{\color{primary} Cột 8} & \textbf{\color{primary} Cột 9} & \textbf{\color{primary} Cột 10} \\
\midrule
\textbf{1.D} & \textbf{2.D} & \textbf{3.B} & \textbf{4.B} & \textbf{5.A} & \textbf{6.C} & \textbf{7.C} & \textbf{8.A} & \textbf{9.B} & \textbf{10.B} \\
\rowcolor{rowtint} \textbf{11.A} & \textbf{12.C} & \textbf{13.C} & \textbf{14.C} & \textbf{15.C} & \textbf{16.A} & \textbf{17.C} & \textbf{18.C} & \textbf{19.A} & \textbf{20.B} \\
\textbf{21.C} & \textbf{22.A} & \textbf{23.D} & \textbf{24.A} & \textbf{25.D} & \textbf{26.D} & \textbf{27.C} & \textbf{28.D} & \textbf{29.C} & \textbf{30.B} \\
\rowcolor{rowtint} \textbf{31.D} & \textbf{32.B} & \textbf{33.D} & \textbf{34.D} & \textbf{35.A} & \textbf{36.B} & \textbf{37.B} & \textbf{38.C} & \textbf{39.B} & \textbf{40.C} \\
\textbf{41.B} & \textbf{42.C} & \textbf{43.B} & \textbf{44.C} & \textbf{45.A} & \textbf{46.C} & \textbf{47.C} & \textbf{48.B} & \textbf{49.A} & \textbf{50.A} \\
\bottomrule
\end{tabularx}
\end{center}\vspace{10pt}
\section*{\color{primary} CHƯƠNG VI: TỈ LỆ THỨC VÀ ĐẠI LƯỢNG TỈ LỆ}
\textit{Trọng tâm kiến thức SGK: Tỉ lệ thức, tính chất của dãy tỉ số bằng nhau, đại lượng tỉ lệ thuận và đại lượng tỉ lệ nghịch.}\par\vspace{4pt}
\subsection*{\color{accent} Phần I. Mức độ Dễ (Nhận biết -- Câu 1 đến Câu 20)}
\textbf{\color{primary} Câu 1.} Tìm x trong tỉ lệ thức x/3 = 4/3.:\par
\hspace*{1.5em}\textbf{A.} x = 2\quad\quad \textbf{B.} x = 3\quad\quad \textbf{C.} x = 4\quad\quad \textbf{D.} x = 12\par\vspace{2pt}
\textbf{\color{primary} Câu 2.} Thực hiện yêu cầu nhận biết sau: Cho biết đại lượng y tỉ lệ thuận với đại lượng x theo hệ số tỉ lệ k = 3. Khi x = 5 thì giá trị của y bằng.\par
\hspace*{1.5em}\textbf{A.} -15\quad\quad \textbf{B.} 3/5\quad\quad \textbf{C.} 8\quad\quad \textbf{D.} 15\par\vspace{2pt}
\textbf{\color{primary} Câu 3.} Trong bài học thuộc Chương VI, xét nội dung: Cho hai đại lượng x và y tỉ lệ nghịch với nhau và khi x = 4 thì y = 6. Hệ số tỉ lệ a = x$\cdot$y bằng. Đáp án đúng là:\par
\hspace*{1.5em}\textbf{A.} 12\quad\quad \textbf{B.} 1.50\quad\quad \textbf{C.} 24\quad\quad \textbf{D.} 10\par\vspace{2pt}
\textbf{\color{primary} Câu 4.} Dựa vào kiến thức chuẩn SGK của Chương VI (Tỉ lệ thức và đại lượng tỉ lệ), tìm x trong tỉ lệ thức x/6 = 10/6.:\par
\hspace*{1.5em}\textbf{A.} x = 10\quad\quad \textbf{B.} x = 60\quad\quad \textbf{C.} x = 5\quad\quad \textbf{D.} x = 6\par\vspace{2pt}
\textbf{\color{primary} Câu 5.} Chọn kết quả đúng nhất khi xét: Cho biết đại lượng y tỉ lệ thuận với đại lượng x theo hệ số tỉ lệ k = 6. Khi x = 5 thì giá trị của y bằng:\par
\hspace*{1.5em}\textbf{A.} -30\quad\quad \textbf{B.} 6/5\quad\quad \textbf{C.} 11\quad\quad \textbf{D.} 30\par\vspace{2pt}
\textbf{\color{primary} Câu 6.} Khi làm bài tập cơ bản về tỉ lệ thức và đại lượng tỉ lệ, một học sinh gặp câu hỏi: Cho hai đại lượng x và y tỉ lệ nghịch với nhau và khi x = 7 thì y = 6. Hệ số tỉ lệ a = x$\cdot$y bằng. Kết quả nào dưới đây đúng?\par
\hspace*{1.5em}\textbf{A.} 0.86\quad\quad \textbf{B.} 42\quad\quad \textbf{C.} 21\quad\quad \textbf{D.} 13\par\vspace{2pt}
\textbf{\color{primary} Câu 7.} Xác định giá trị hoặc khẳng định đúng cho bài toán sau: Tìm x trong tỉ lệ thức x/9 = 16/9.:\par
\hspace*{1.5em}\textbf{A.} x = 8\quad\quad \textbf{B.} x = 16\quad\quad \textbf{C.} x = 144\quad\quad \textbf{D.} x = 9\par\vspace{2pt}
\textbf{\color{primary} Câu 8.} Áp dụng trực tiếp định nghĩa và quy tắc trong Chương VI: Cho biết đại lượng y tỉ lệ thuận với đại lượng x theo hệ số tỉ lệ k = 9. Khi x = 5 thì giá trị của y bằng:\par
\hspace*{1.5em}\textbf{A.} -45\quad\quad \textbf{B.} 45\quad\quad \textbf{C.} 9/5\quad\quad \textbf{D.} 14\par\vspace{2pt}
\textbf{\color{primary} Câu 9.} Kiểm tra nhanh kiến thức nền tảng: Cho hai đại lượng x và y tỉ lệ nghịch với nhau và khi x = 10 thì y = 6. Hệ số tỉ lệ a = x$\cdot$y bằng:\par
\hspace*{1.5em}\textbf{A.} 30\quad\quad \textbf{B.} 60\quad\quad \textbf{C.} 0.60\quad\quad \textbf{D.} 16\par\vspace{2pt}
\textbf{\color{primary} Câu 10.} Cho dữ kiện nhận biết sau đây: Tìm x trong tỉ lệ thức x/12 = 22/12.. Phương án chính xác là:\par
\hspace*{1.5em}\textbf{A.} x = 12\quad\quad \textbf{B.} x = 11\quad\quad \textbf{C.} x = 264\quad\quad \textbf{D.} x = 22\par\vspace{2pt}
\textbf{\color{primary} Câu 11.} Hãy chỉ ra kết quả đúng trong bốn lựa chọn A, B, C, D đối với yêu cầu: Cho biết đại lượng y tỉ lệ thuận với đại lượng x theo hệ số tỉ lệ k = 12. Khi x = 5 thì giá trị của y bằng:\par
\hspace*{1.5em}\textbf{A.} 12/5\quad\quad \textbf{B.} 17\quad\quad \textbf{C.} 60\quad\quad \textbf{D.} -60\par\vspace{2pt}
\textbf{\color{primary} Câu 12.} Câu hỏi ôn tập lý thuyết và kỹ năng cơ bản (Chương VI): Cho hai đại lượng x và y tỉ lệ nghịch với nhau và khi x = 13 thì y = 6. Hệ số tỉ lệ a = x$\cdot$y bằng:\par
\hspace*{1.5em}\textbf{A.} 19\quad\quad \textbf{B.} 0.46\quad\quad \textbf{C.} 39\quad\quad \textbf{D.} 78\par\vspace{2pt}
\textbf{\color{primary} Câu 13.} Theo nội dung sách giáo khoa Kết nối tri thức với cuộc sống (Chương VI), tìm x trong tỉ lệ thức x/15 = 28/15.:\par
\hspace*{1.5em}\textbf{A.} x = 15\quad\quad \textbf{B.} x = 420\quad\quad \textbf{C.} x = 14\quad\quad \textbf{D.} x = 28\par\vspace{2pt}
\textbf{\color{primary} Câu 14.} Tìm đáp số đúng của bài toán cơ bản sau: Cho biết đại lượng y tỉ lệ thuận với đại lượng x theo hệ số tỉ lệ k = 15. Khi x = 5 thì giá trị của y bằng:\par
\hspace*{1.5em}\textbf{A.} 15/5\quad\quad \textbf{B.} 75\quad\quad \textbf{C.} 20\quad\quad \textbf{D.} -75\par\vspace{2pt}
\textbf{\color{primary} Câu 15.} Xét đối tượng toán học trong Chương VI: Cho hai đại lượng x và y tỉ lệ nghịch với nhau và khi x = 16 thì y = 6. Hệ số tỉ lệ a = x$\cdot$y bằng. Giá trị hoặc tính chất tương ứng là:\par
\hspace*{1.5em}\textbf{A.} 48\quad\quad \textbf{B.} 96\quad\quad \textbf{C.} 0.38\quad\quad \textbf{D.} 22\par\vspace{2pt}
\textbf{\color{primary} Câu 16.} Hoàn thành phép xác định cơ bản sau đây: Tìm x trong tỉ lệ thức x/18 = 34/18.:\par
\hspace*{1.5em}\textbf{A.} x = 18\quad\quad \textbf{B.} x = 34\quad\quad \textbf{C.} x = 17\quad\quad \textbf{D.} x = 612\par\vspace{2pt}
\textbf{\color{primary} Câu 17.} Trong các khẳng định hoặc giá trị dưới đây, đâu là đáp án đúng cho câu hỏi: Cho biết đại lượng y tỉ lệ thuận với đại lượng x theo hệ số tỉ lệ k = 18. Khi x = 5 thì giá trị của y bằng?\par
\hspace*{1.5em}\textbf{A.} -90\quad\quad \textbf{B.} 90\quad\quad \textbf{C.} 23\quad\quad \textbf{D.} 18/5\par\vspace{2pt}
\textbf{\color{primary} Câu 18.} Bài tập khởi động Chương VI: Cho hai đại lượng x và y tỉ lệ nghịch với nhau và khi x = 19 thì y = 6. Hệ số tỉ lệ a = x$\cdot$y bằng:\par
\hspace*{1.5em}\textbf{A.} 114\quad\quad \textbf{B.} 57\quad\quad \textbf{C.} 25\quad\quad \textbf{D.} 0.32\par\vspace{2pt}
\textbf{\color{primary} Câu 19.} Quan sát dữ kiện đề bài: Tìm x trong tỉ lệ thức x/21 = 40/21.. Hãy chọn phương án đúng:\par
\hspace*{1.5em}\textbf{A.} x = 40\quad\quad \textbf{B.} x = 20\quad\quad \textbf{C.} x = 21\quad\quad \textbf{D.} x = 840\par\vspace{2pt}
\textbf{\color{primary} Câu 20.} Sử dụng công thức hoặc tính chất đầu bài của Chương VI để giải quyết yêu cầu: Cho biết đại lượng y tỉ lệ thuận với đại lượng x theo hệ số tỉ lệ k = 21. Khi x = 5 thì giá trị của y bằng:\par
\hspace*{1.5em}\textbf{A.} 105\quad\quad \textbf{B.} 21/5\quad\quad \textbf{C.} -105\quad\quad \textbf{D.} 26\par\vspace{2pt}
\subsection*{\color{accent} Phần II. Mức độ Trung bình (Thông hiểu -- Câu 21 đến Câu 40)}
\textbf{\color{primary} Câu 21.} Tìm hai số x và y biết x/3 = y/5 và x + y = 16.:\par
\hspace*{1.5em}\textbf{A.} x = 4; y = 12\quad\quad \textbf{B.} x = 8; y = 8\quad\quad \textbf{C.} x = 6; y = 10\quad\quad \textbf{D.} x = 10; y = 6\par\vspace{2pt}
\textbf{\color{primary} Câu 22.} Vận dụng quy tắc trong Chương VI (Tỉ lệ thức và đại lượng tỉ lệ) để giải bài toán: Tìm hai số x, y biết x/7 = y/4 và x - y = 12.:\par
\hspace*{1.5em}\textbf{A.} x = 21; y = 12\quad\quad \textbf{B.} x = 16; y = 28\quad\quad \textbf{C.} x = 56; y = 32\quad\quad \textbf{D.} x = 28; y = 16\par\vspace{2pt}
\textbf{\color{primary} Câu 23.} Biến đổi và tính toán theo kiến thức Chương VI: Với cùng năng suất, 6 công nhân hoàn thành một công việc trong 12 giờ. Hỏi 12 công nhân hoàn thành công việc đó trong bao nhiêu giờ. Kết quả thu được là:\par
\hspace*{1.5em}\textbf{A.} 24 giờ\quad\quad \textbf{B.} 8 giờ\quad\quad \textbf{C.} 6 giờ\quad\quad \textbf{D.} 4 giờ\par\vspace{2pt}
\textbf{\color{primary} Câu 24.} Xét bài toán mức độ thông hiểu sau: Tìm hai số x và y biết x/3 = y/5 và x + y = 40.. Đáp án đúng là:\par
\hspace*{1.5em}\textbf{A.} x = 15; y = 25\quad\quad \textbf{B.} x = 10; y = 30\quad\quad \textbf{C.} x = 25; y = 15\quad\quad \textbf{D.} x = 20; y = 20\par\vspace{2pt}
\textbf{\color{primary} Câu 25.} Để hoàn thành phiếu học tập Chương VI, hãy xác định: Tìm hai số x, y biết x/7 = y/4 và x - y = 21.:\par
\hspace*{1.5em}\textbf{A.} x = 42; y = 24\quad\quad \textbf{B.} x = 98; y = 56\quad\quad \textbf{C.} x = 28; y = 49\quad\quad \textbf{D.} x = 49; y = 28\par\vspace{2pt}
\textbf{\color{primary} Câu 26.} Phân tích dữ kiện bài toán: Với cùng năng suất, 6 công nhân hoàn thành một công việc trong 12 giờ. Hỏi 12 công nhân hoàn thành công việc đó trong bao nhiêu giờ. Giá trị cần tìm bằng:\par
\hspace*{1.5em}\textbf{A.} 8 giờ\quad\quad \textbf{B.} 6 giờ\quad\quad \textbf{C.} 24 giờ\quad\quad \textbf{D.} 4 giờ\par\vspace{2pt}
\textbf{\color{primary} Câu 27.} Áp dụng tính chất đã học ở Chương VI vào bài tập sau: Tìm hai số x và y biết x/3 = y/5 và x + y = 64.:\par
\hspace*{1.5em}\textbf{A.} x = 32; y = 32\quad\quad \textbf{B.} x = 40; y = 24\quad\quad \textbf{C.} x = 16; y = 48\quad\quad \textbf{D.} x = 24; y = 40\par\vspace{2pt}
\textbf{\color{primary} Câu 28.} Thực hiện các bước suy luận hợp lý để giải quyết yêu cầu: Tìm hai số x, y biết x/7 = y/4 và x - y = 30.:\par
\hspace*{1.5em}\textbf{A.} x = 140; y = 80\quad\quad \textbf{B.} x = 40; y = 70\quad\quad \textbf{C.} x = 63; y = 36\quad\quad \textbf{D.} x = 70; y = 40\par\vspace{2pt}
\textbf{\color{primary} Câu 29.} Cho bài toán tự luyện (Tỉ lệ thức và đại lượng tỉ lệ): Với cùng năng suất, 6 công nhân hoàn thành một công việc trong 12 giờ. Hỏi 12 công nhân hoàn thành công việc đó trong bao nhiêu giờ. Phương án nào sau đây chính xác?\par
\hspace*{1.5em}\textbf{A.} 6 giờ\quad\quad \textbf{B.} 24 giờ\quad\quad \textbf{C.} 4 giờ\quad\quad \textbf{D.} 8 giờ\par\vspace{2pt}
\textbf{\color{primary} Câu 30.} Kết hợp các phép biến đổi trong Chương VI, hãy tìm kết quả của bài toán: Tìm hai số x và y biết x/3 = y/5 và x + y = 88.:\par
\hspace*{1.5em}\textbf{A.} x = 44; y = 44\quad\quad \textbf{B.} x = 55; y = 33\quad\quad \textbf{C.} x = 33; y = 55\quad\quad \textbf{D.} x = 22; y = 66\par\vspace{2pt}
\textbf{\color{primary} Câu 31.} Một bài kiểm tra định kỳ nội dung Chương VI nêu câu hỏi: Tìm hai số x, y biết x/7 = y/4 và x - y = 39.. Hãy chọn đáp số đúng:\par
\hspace*{1.5em}\textbf{A.} x = 52; y = 91\quad\quad \textbf{B.} x = 182; y = 104\quad\quad \textbf{C.} x = 91; y = 52\quad\quad \textbf{D.} x = 84; y = 48\par\vspace{2pt}
\textbf{\color{primary} Câu 32.} Tìm lời giải đúng cho bài toán thông hiểu: Với cùng năng suất, 6 công nhân hoàn thành một công việc trong 12 giờ. Hỏi 12 công nhân hoàn thành công việc đó trong bao nhiêu giờ:\par
\hspace*{1.5em}\textbf{A.} 4 giờ\quad\quad \textbf{B.} 24 giờ\quad\quad \textbf{C.} 6 giờ\quad\quad \textbf{D.} 8 giờ\par\vspace{2pt}
\textbf{\color{primary} Câu 33.} Sử dụng mối liên hệ giữa các đại lượng trong Chương VI: Tìm hai số x và y biết x/3 = y/5 và x + y = 112.:\par
\hspace*{1.5em}\textbf{A.} x = 70; y = 42\quad\quad \textbf{B.} x = 28; y = 84\quad\quad \textbf{C.} x = 56; y = 56\quad\quad \textbf{D.} x = 42; y = 70\par\vspace{2pt}
\textbf{\color{primary} Câu 34.} Hãy tính toán và chọn phương án đúng đối với bài tập: Tìm hai số x, y biết x/7 = y/4 và x - y = 48.:\par
\hspace*{1.5em}\textbf{A.} x = 112; y = 64\quad\quad \textbf{B.} x = 64; y = 112\quad\quad \textbf{C.} x = 224; y = 128\quad\quad \textbf{D.} x = 105; y = 60\par\vspace{2pt}
\textbf{\color{primary} Câu 35.} Dựa trên các tính chất trọng tâm của Chương VI, xét bài toán: Với cùng năng suất, 6 công nhân hoàn thành một công việc trong 12 giờ. Hỏi 12 công nhân hoàn thành công việc đó trong bao nhiêu giờ:\par
\hspace*{1.5em}\textbf{A.} 8 giờ\quad\quad \textbf{B.} 24 giờ\quad\quad \textbf{C.} 6 giờ\quad\quad \textbf{D.} 4 giờ\par\vspace{2pt}
\textbf{\color{primary} Câu 36.} Thực hành giải toán (Tỉ lệ thức và đại lượng tỉ lệ): Tìm hai số x và y biết x/3 = y/5 và x + y = 136.. Đáp án nào dưới đây thỏa mãn?\par
\hspace*{1.5em}\textbf{A.} x = 51; y = 85\quad\quad \textbf{B.} x = 85; y = 51\quad\quad \textbf{C.} x = 68; y = 68\quad\quad \textbf{D.} x = 34; y = 102\par\vspace{2pt}
\textbf{\color{primary} Câu 37.} Kiểm tra kỹ năng biến đổi và suy luận qua bài toán: Tìm hai số x, y biết x/7 = y/4 và x - y = 57.:\par
\hspace*{1.5em}\textbf{A.} x = 266; y = 152\quad\quad \textbf{B.} x = 76; y = 133\quad\quad \textbf{C.} x = 133; y = 76\quad\quad \textbf{D.} x = 126; y = 72\par\vspace{2pt}
\textbf{\color{primary} Câu 38.} Cho giả thiết thuộc phạm vi Chương VI: Với cùng năng suất, 6 công nhân hoàn thành một công việc trong 12 giờ. Hỏi 12 công nhân hoàn thành công việc đó trong bao nhiêu giờ. Kết luận đúng là:\par
\hspace*{1.5em}\textbf{A.} 6 giờ\quad\quad \textbf{B.} 8 giờ\quad\quad \textbf{C.} 4 giờ\quad\quad \textbf{D.} 24 giờ\par\vspace{2pt}
\textbf{\color{primary} Câu 39.} Rút gọn hoặc tính giá trị theo yêu cầu sau của Chương VI: Tìm hai số x và y biết x/3 = y/5 và x + y = 160.:\par
\hspace*{1.5em}\textbf{A.} x = 40; y = 120\quad\quad \textbf{B.} x = 60; y = 100\quad\quad \textbf{C.} x = 80; y = 80\quad\quad \textbf{D.} x = 100; y = 60\par\vspace{2pt}
\textbf{\color{primary} Câu 40.} Xác định phương án đúng sau khi thực hiện phép tính / suy luận: Tìm hai số x, y biết x/7 = y/4 và x - y = 66.:\par
\hspace*{1.5em}\textbf{A.} x = 88; y = 154\quad\quad \textbf{B.} x = 147; y = 84\quad\quad \textbf{C.} x = 154; y = 88\quad\quad \textbf{D.} x = 308; y = 176\par\vspace{2pt}
\subsection*{\color{accent} Phần III. Mức độ Khó (Vận dụng \& Vận dụng cao -- Câu 41 đến Câu 50)}
\textbf{\color{primary} Câu 41.} Ba lớp 7A, 7B, 7C trồng được tất cả 72 cây xanh. Biết số cây trồng được của ba lớp tỉ lệ với 3 : 4 : 5. Số cây lớp 7C trồng được là:\par
\hspace*{1.5em}\textbf{A.} 30 cây\quad\quad \textbf{B.} 18 cây\quad\quad \textbf{C.} 36 cây\quad\quad \textbf{D.} 24 cây\par\vspace{2pt}
\textbf{\color{primary} Câu 42.} Bài toán vận dụng (Chương VI – Tỉ lệ thức và đại lượng tỉ lệ): Tìm hai số dương x, y biết x/2 = y/3 và x$\cdot$y = 54.. Kết quả đúng là:\par
\hspace*{1.5em}\textbf{A.} x = 9; y = 6\quad\quad \textbf{B.} x = 6; y = 9\quad\quad \textbf{C.} x = 18; y = 27\quad\quad \textbf{D.} x = 3; y = 18\par\vspace{2pt}
\textbf{\color{primary} Câu 43.} Vận dụng kiến thức nâng cao trong phạm vi Chương VI để giải bài toán: Ba lớp 7A, 7B, 7C trồng được tất cả 144 cây xanh. Biết số cây trồng được của ba lớp tỉ lệ với 3 : 4 : 5. Số cây lớp 7C trồng được là:\par
\hspace*{1.5em}\textbf{A.} 60 cây\quad\quad \textbf{B.} 36 cây\quad\quad \textbf{C.} 48 cây\quad\quad \textbf{D.} 72 cây\par\vspace{2pt}
\textbf{\color{primary} Câu 44.} Phân tích và giải bài toán tổng hợp sau của Chương VI: Tìm hai số dương x, y biết x/2 = y/3 và x$\cdot$y = 150.:\par
\hspace*{1.5em}\textbf{A.} x = 5; y = 30\quad\quad \textbf{B.} x = 50; y = 75\quad\quad \textbf{C.} x = 15; y = 10\quad\quad \textbf{D.} x = 10; y = 15\par\vspace{2pt}
\textbf{\color{primary} Câu 45.} Thử thách tư duy trong phạm vi kiến thức Chương VI: Ba lớp 7A, 7B, 7C trồng được tất cả 216 cây xanh. Biết số cây trồng được của ba lớp tỉ lệ với 3 : 4 : 5. Số cây lớp 7C trồng được là. Đáp số chính xác là:\par
\hspace*{1.5em}\textbf{A.} 90 cây\quad\quad \textbf{B.} 72 cây\quad\quad \textbf{C.} 54 cây\quad\quad \textbf{D.} 108 cây\par\vspace{2pt}
\textbf{\color{primary} Câu 46.} Sử dụng phương pháp giải toán của Chương VI để tìm đáp án cho bài toán: Tìm hai số dương x, y biết x/2 = y/3 và x$\cdot$y = 294.:\par
\hspace*{1.5em}\textbf{A.} x = 21; y = 14\quad\quad \textbf{B.} x = 14; y = 21\quad\quad \textbf{C.} x = 7; y = 42\quad\quad \textbf{D.} x = 98; y = 147\par\vspace{2pt}
\textbf{\color{primary} Câu 47.} Xét bài toán vận dụng thực tiễn / nâng cao thuộc Chương VI: Ba lớp 7A, 7B, 7C trồng được tất cả 288 cây xanh. Biết số cây trồng được của ba lớp tỉ lệ với 3 : 4 : 5. Số cây lớp 7C trồng được là:\par
\hspace*{1.5em}\textbf{A.} 72 cây\quad\quad \textbf{B.} 144 cây\quad\quad \textbf{C.} 96 cây\quad\quad \textbf{D.} 120 cây\par\vspace{2pt}
\textbf{\color{primary} Câu 48.} Tìm kết quả cuối cùng của bài toán phân loại học sinh khá giỏi (Chương VI): Tìm hai số dương x, y biết x/2 = y/3 và x$\cdot$y = 486.:\par
\hspace*{1.5em}\textbf{A.} x = 27; y = 18\quad\quad \textbf{B.} x = 18; y = 27\quad\quad \textbf{C.} x = 162; y = 243\quad\quad \textbf{D.} x = 9; y = 54\par\vspace{2pt}
\textbf{\color{primary} Câu 49.} Kết hợp các tính chất đã học trong Chương VI, hãy giải quyết bài toán: Ba lớp 7A, 7B, 7C trồng được tất cả 360 cây xanh. Biết số cây trồng được của ba lớp tỉ lệ với 3 : 4 : 5. Số cây lớp 7C trồng được là:\par
\hspace*{1.5em}\textbf{A.} 120 cây\quad\quad \textbf{B.} 90 cây\quad\quad \textbf{C.} 150 cây\quad\quad \textbf{D.} 180 cây\par\vspace{2pt}
\textbf{\color{primary} Câu 50.} Câu hỏi vận dụng cao cuối Chương VI: Tìm hai số dương x, y biết x/2 = y/3 và x$\cdot$y = 726.. Phương án đúng là:\par
\hspace*{1.5em}\textbf{A.} x = 11; y = 66\quad\quad \textbf{B.} x = 33; y = 22\quad\quad \textbf{C.} x = 22; y = 33\quad\quad \textbf{D.} x = 242; y = 363\par\vspace{2pt}
\subsection*{\color{primary} Bảng đáp án Chương VI: Tỉ lệ thức và đại lượng tỉ lệ}
\begin{center}
\renewcommand{\arraystretch}{1.25}
\begin{tabularx}{\textwidth}{*{10}{>{\centering\arraybackslash}X}}
\toprule
\textbf{\color{primary} Cột 1} & \textbf{\color{primary} Cột 2} & \textbf{\color{primary} Cột 3} & \textbf{\color{primary} Cột 4} & \textbf{\color{primary} Cột 5} & \textbf{\color{primary} Cột 6} & \textbf{\color{primary} Cột 7} & \textbf{\color{primary} Cột 8} & \textbf{\color{primary} Cột 9} & \textbf{\color{primary} Cột 10} \\
\midrule
\textbf{1.C} & \textbf{2.D} & \textbf{3.C} & \textbf{4.A} & \textbf{5.D} & \textbf{6.B} & \textbf{7.B} & \textbf{8.B} & \textbf{9.B} & \textbf{10.D} \\
\rowcolor{rowtint} \textbf{11.C} & \textbf{12.D} & \textbf{13.D} & \textbf{14.B} & \textbf{15.B} & \textbf{16.B} & \textbf{17.B} & \textbf{18.A} & \textbf{19.A} & \textbf{20.A} \\
\textbf{21.C} & \textbf{22.D} & \textbf{23.C} & \textbf{24.A} & \textbf{25.D} & \textbf{26.B} & \textbf{27.D} & \textbf{28.D} & \textbf{29.A} & \textbf{30.C} \\
\rowcolor{rowtint} \textbf{31.C} & \textbf{32.C} & \textbf{33.D} & \textbf{34.A} & \textbf{35.C} & \textbf{36.A} & \textbf{37.C} & \textbf{38.A} & \textbf{39.B} & \textbf{40.C} \\
\textbf{41.A} & \textbf{42.B} & \textbf{43.A} & \textbf{44.D} & \textbf{45.A} & \textbf{46.B} & \textbf{47.D} & \textbf{48.B} & \textbf{49.C} & \textbf{50.C} \\
\bottomrule
\end{tabularx}
\end{center}\vspace{10pt}
\section*{\color{primary} CHƯƠNG VII: BIỂU THỨC ĐẠI SỐ VÀ ĐA THỨC MỘT BIẾN}
\textit{Trọng tâm kiến thức SGK: Biểu thức đại số, đa thức một biến, bậc và hệ số của đa thức, nghiệm của đa thức một biến và các phép cộng, trừ, nhân, chia đa thức một biến.}\par\vspace{4pt}
\subsection*{\color{accent} Phần I. Mức độ Dễ (Nhận biết -- Câu 1 đến Câu 20)}
\textbf{\color{primary} Câu 1.} Bậc của đa thức một biến P(x) = 2x$^4$ - 3x$^2$ + 5x - 1 là:\par
\hspace*{1.5em}\textbf{A.} Bậc 4\quad\quad \textbf{B.} Bậc 7\quad\quad \textbf{C.} Bậc 4 (1)\quad\quad \textbf{D.} Bậc 2\par\vspace{2pt}
\textbf{\color{primary} Câu 2.} Thực hiện yêu cầu nhận biết sau: Hệ số cao nhất và hệ số tự do của đa thức Q(x) = 4x$^3$ - 5x + 3 lần lượt là.\par
\hspace*{1.5em}\textbf{A.} -5 và 3\quad\quad \textbf{B.} 3 và 3\quad\quad \textbf{C.} 4 và -5\quad\quad \textbf{D.} 4 và 3\par\vspace{2pt}
\textbf{\color{primary} Câu 3.} Trong bài học thuộc Chương VII, xét nội dung: Nghiệm của đa thức bậc nhất F(x) = 2x - 10 là. Đáp án đúng là:\par
\hspace*{1.5em}\textbf{A.} x = -5\quad\quad \textbf{B.} x = 10\quad\quad \textbf{C.} x = 5\quad\quad \textbf{D.} x = 0\par\vspace{2pt}
\textbf{\color{primary} Câu 4.} Dựa vào kiến thức chuẩn SGK của Chương VII (Biểu thức đại số và đa thức một biến), bậc của đa thức một biến P(x) = 5x$^4$ - 3x$^2$ + 8x - 1 là:\par
\hspace*{1.5em}\textbf{A.} Bậc 5\quad\quad \textbf{B.} Bậc 7\quad\quad \textbf{C.} Bậc 2\quad\quad \textbf{D.} Bậc 4\par\vspace{2pt}
\textbf{\color{primary} Câu 5.} Chọn kết quả đúng nhất khi xét: Hệ số cao nhất và hệ số tự do của đa thức Q(x) = 7x$^3$ - 5x + 9 lần lượt là:\par
\hspace*{1.5em}\textbf{A.} 3 và 9\quad\quad \textbf{B.} 7 và 9\quad\quad \textbf{C.} -5 và 9\quad\quad \textbf{D.} 7 và -5\par\vspace{2pt}
\textbf{\color{primary} Câu 6.} Khi làm bài tập cơ bản về biểu thức đại số và đa thức một biến, một học sinh gặp câu hỏi: Nghiệm của đa thức bậc nhất F(x) = 2x - 16 là. Kết quả nào dưới đây đúng?\par
\hspace*{1.5em}\textbf{A.} x = -8\quad\quad \textbf{B.} x = 8\quad\quad \textbf{C.} x = 16\quad\quad \textbf{D.} x = 0\par\vspace{2pt}
\textbf{\color{primary} Câu 7.} Xác định giá trị hoặc khẳng định đúng cho bài toán sau: Bậc của đa thức một biến P(x) = 8x$^4$ - 3x$^2$ + 11x - 1 là:\par
\hspace*{1.5em}\textbf{A.} Bậc 4\quad\quad \textbf{B.} Bậc 8\quad\quad \textbf{C.} Bậc 7\quad\quad \textbf{D.} Bậc 2\par\vspace{2pt}
\textbf{\color{primary} Câu 8.} Áp dụng trực tiếp định nghĩa và quy tắc trong Chương VII: Hệ số cao nhất và hệ số tự do của đa thức Q(x) = 10x$^3$ - 5x + 15 lần lượt là:\par
\hspace*{1.5em}\textbf{A.} 10 và 15\quad\quad \textbf{B.} 10 và -5\quad\quad \textbf{C.} 3 và 15\quad\quad \textbf{D.} -5 và 15\par\vspace{2pt}
\textbf{\color{primary} Câu 9.} Kiểm tra nhanh kiến thức nền tảng: Nghiệm của đa thức bậc nhất F(x) = 2x - 22 là:\par
\hspace*{1.5em}\textbf{A.} x = -11\quad\quad \textbf{B.} x = 22\quad\quad \textbf{C.} x = 11\quad\quad \textbf{D.} x = 0\par\vspace{2pt}
\textbf{\color{primary} Câu 10.} Cho dữ kiện nhận biết sau đây: Bậc của đa thức một biến P(x) = 11x$^4$ - 3x$^2$ + 14x - 1 là. Phương án chính xác là:\par
\hspace*{1.5em}\textbf{A.} Bậc 7\quad\quad \textbf{B.} Bậc 2\quad\quad \textbf{C.} Bậc 4\quad\quad \textbf{D.} Bậc 11\par\vspace{2pt}
\textbf{\color{primary} Câu 11.} Hãy chỉ ra kết quả đúng trong bốn lựa chọn A, B, C, D đối với yêu cầu: Hệ số cao nhất và hệ số tự do của đa thức Q(x) = 13x$^3$ - 5x + 21 lần lượt là:\par
\hspace*{1.5em}\textbf{A.} 13 và -5\quad\quad \textbf{B.} -5 và 21\quad\quad \textbf{C.} 3 và 21\quad\quad \textbf{D.} 13 và 21\par\vspace{2pt}
\textbf{\color{primary} Câu 12.} Câu hỏi ôn tập lý thuyết và kỹ năng cơ bản (Chương VII): Nghiệm của đa thức bậc nhất F(x) = 2x - 28 là:\par
\hspace*{1.5em}\textbf{A.} x = -14\quad\quad \textbf{B.} x = 14\quad\quad \textbf{C.} x = 28\quad\quad \textbf{D.} x = 0\par\vspace{2pt}
\textbf{\color{primary} Câu 13.} Theo nội dung sách giáo khoa Kết nối tri thức với cuộc sống (Chương VII), bậc của đa thức một biến P(x) = 14x$^4$ - 3x$^2$ + 17x - 1 là:\par
\hspace*{1.5em}\textbf{A.} Bậc 2\quad\quad \textbf{B.} Bậc 14\quad\quad \textbf{C.} Bậc 7\quad\quad \textbf{D.} Bậc 4\par\vspace{2pt}
\textbf{\color{primary} Câu 14.} Tìm đáp số đúng của bài toán cơ bản sau: Hệ số cao nhất và hệ số tự do của đa thức Q(x) = 16x$^3$ - 5x + 27 lần lượt là:\par
\hspace*{1.5em}\textbf{A.} -5 và 27\quad\quad \textbf{B.} 3 và 27\quad\quad \textbf{C.} 16 và -5\quad\quad \textbf{D.} 16 và 27\par\vspace{2pt}
\textbf{\color{primary} Câu 15.} Xét đối tượng toán học trong Chương VII: Nghiệm của đa thức bậc nhất F(x) = 2x - 34 là. Giá trị hoặc tính chất tương ứng là:\par
\hspace*{1.5em}\textbf{A.} x = 0\quad\quad \textbf{B.} x = 34\quad\quad \textbf{C.} x = -17\quad\quad \textbf{D.} x = 17\par\vspace{2pt}
\textbf{\color{primary} Câu 16.} Hoàn thành phép xác định cơ bản sau đây: Bậc của đa thức một biến P(x) = 17x$^4$ - 3x$^2$ + 20x - 1 là:\par
\hspace*{1.5em}\textbf{A.} Bậc 7\quad\quad \textbf{B.} Bậc 17\quad\quad \textbf{C.} Bậc 4\quad\quad \textbf{D.} Bậc 2\par\vspace{2pt}
\textbf{\color{primary} Câu 17.} Trong các khẳng định hoặc giá trị dưới đây, đâu là đáp án đúng cho câu hỏi: Hệ số cao nhất và hệ số tự do của đa thức Q(x) = 19x$^3$ - 5x + 33 lần lượt là?\par
\hspace*{1.5em}\textbf{A.} 19 và -5\quad\quad \textbf{B.} 19 và 33\quad\quad \textbf{C.} -5 và 33\quad\quad \textbf{D.} 3 và 33\par\vspace{2pt}
\textbf{\color{primary} Câu 18.} Bài tập khởi động Chương VII: Nghiệm của đa thức bậc nhất F(x) = 2x - 40 là:\par
\hspace*{1.5em}\textbf{A.} x = 0\quad\quad \textbf{B.} x = 40\quad\quad \textbf{C.} x = 20\quad\quad \textbf{D.} x = -20\par\vspace{2pt}
\textbf{\color{primary} Câu 19.} Quan sát dữ kiện đề bài: Bậc của đa thức một biến P(x) = 20x$^4$ - 3x$^2$ + 23x - 1 là. Hãy chọn phương án đúng:\par
\hspace*{1.5em}\textbf{A.} Bậc 20\quad\quad \textbf{B.} Bậc 4\quad\quad \textbf{C.} Bậc 2\quad\quad \textbf{D.} Bậc 7\par\vspace{2pt}
\textbf{\color{primary} Câu 20.} Sử dụng công thức hoặc tính chất đầu bài của Chương VII để giải quyết yêu cầu: Hệ số cao nhất và hệ số tự do của đa thức Q(x) = 22x$^3$ - 5x + 39 lần lượt là:\par
\hspace*{1.5em}\textbf{A.} -5 và 39\quad\quad \textbf{B.} 22 và 39\quad\quad \textbf{C.} 22 và -5\quad\quad \textbf{D.} 3 và 39\par\vspace{2pt}
\subsection*{\color{accent} Phần II. Mức độ Trung bình (Thông hiểu -- Câu 21 đến Câu 40)}
\textbf{\color{primary} Câu 21.} Giá trị của biểu thức đại số A(x) = x$^2$ - 2x + 5 tại x = 3 bằng:\par
\hspace*{1.5em}\textbf{A.} 8\quad\quad \textbf{B.} 5\quad\quad \textbf{C.} 14\quad\quad \textbf{D.} 2\par\vspace{2pt}
\textbf{\color{primary} Câu 22.} Vận dụng quy tắc trong Chương VII (Biểu thức đại số và đa thức một biến) để giải bài toán: Cho hai đa thức M(x) = 3x$^2$ + 3x - 4 và N(x) = -x$^2$ + 4x + 4. Tổng M(x) + N(x) bằng:\par
\hspace*{1.5em}\textbf{A.} 2x$^2$ - x\quad\quad \textbf{B.} 2x$^2$ + 7x\quad\quad \textbf{C.} 4x$^2$ + 7x - 8\quad\quad \textbf{D.} 2x$^2$ + 7x - 8\par\vspace{2pt}
\textbf{\color{primary} Câu 23.} Biến đổi và tính toán theo kiến thức Chương VII: Thực hiện phép nhân đơn thức với đa thức: 2x $\cdot$ (x$^2$ - 5x + 1) ta được. Kết quả thu được là:\par
\hspace*{1.5em}\textbf{A.} 2x$^3$ - 10x$^2$ + 2x\quad\quad \textbf{B.} 2x$^2$ - 10x + 2\quad\quad \textbf{C.} 2x$^3$ - 5x$^2$ + 2x\quad\quad \textbf{D.} 3x$^3$ - 10x$^2$ + 2x\par\vspace{2pt}
\textbf{\color{primary} Câu 24.} Xét bài toán mức độ thông hiểu sau: Giá trị của biểu thức đại số A(x) = x$^2$ - 2x + 8 tại x = 3 bằng. Đáp án đúng là:\par
\hspace*{1.5em}\textbf{A.} 8\quad\quad \textbf{B.} 5\quad\quad \textbf{C.} 11\quad\quad \textbf{D.} 17\par\vspace{2pt}
\textbf{\color{primary} Câu 25.} Để hoàn thành phiếu học tập Chương VII, hãy xác định: Cho hai đa thức M(x) = 3x$^2$ + 6x - 4 và N(x) = -x$^2$ + 7x + 4. Tổng M(x) + N(x) bằng:\par
\hspace*{1.5em}\textbf{A.} 4x$^2$ + 13x - 8\quad\quad \textbf{B.} 2x$^2$ + 13x\quad\quad \textbf{C.} 2x$^2$ - x\quad\quad \textbf{D.} 2x$^2$ + 13x - 8\par\vspace{2pt}
\textbf{\color{primary} Câu 26.} Phân tích dữ kiện bài toán: Thực hiện phép nhân đơn thức với đa thức: 2x $\cdot$ (x$^2$ - 8x + 1) ta được. Giá trị cần tìm bằng:\par
\hspace*{1.5em}\textbf{A.} 3x$^3$ - 16x$^2$ + 2x\quad\quad \textbf{B.} 2x$^3$ - 8x$^2$ + 2x\quad\quad \textbf{C.} 2x$^3$ - 16x$^2$ + 2x\quad\quad \textbf{D.} 2x$^2$ - 16x + 2\par\vspace{2pt}
\textbf{\color{primary} Câu 27.} Áp dụng tính chất đã học ở Chương VII vào bài tập sau: Giá trị của biểu thức đại số A(x) = x$^2$ - 2x + 11 tại x = 3 bằng:\par
\hspace*{1.5em}\textbf{A.} 11\quad\quad \textbf{B.} 8\quad\quad \textbf{C.} 14\quad\quad \textbf{D.} 20\par\vspace{2pt}
\textbf{\color{primary} Câu 28.} Thực hiện các bước suy luận hợp lý để giải quyết yêu cầu: Cho hai đa thức M(x) = 3x$^2$ + 9x - 4 và N(x) = -x$^2$ + 10x + 4. Tổng M(x) + N(x) bằng:\par
\hspace*{1.5em}\textbf{A.} 2x$^2$ + 19x - 8\quad\quad \textbf{B.} 2x$^2$ + 19x\quad\quad \textbf{C.} 4x$^2$ + 19x - 8\quad\quad \textbf{D.} 2x$^2$ - x\par\vspace{2pt}
\textbf{\color{primary} Câu 29.} Cho bài toán tự luyện (Biểu thức đại số và đa thức một biến): Thực hiện phép nhân đơn thức với đa thức: 2x $\cdot$ (x$^2$ - 11x + 1) ta được. Phương án nào sau đây chính xác?\par
\hspace*{1.5em}\textbf{A.} 2x$^3$ - 11x$^2$ + 2x\quad\quad \textbf{B.} 3x$^3$ - 22x$^2$ + 2x\quad\quad \textbf{C.} 2x$^3$ - 22x$^2$ + 2x\quad\quad \textbf{D.} 2x$^2$ - 22x + 2\par\vspace{2pt}
\textbf{\color{primary} Câu 30.} Kết hợp các phép biến đổi trong Chương VII, hãy tìm kết quả của bài toán: Giá trị của biểu thức đại số A(x) = x$^2$ - 2x + 14 tại x = 3 bằng:\par
\hspace*{1.5em}\textbf{A.} 17\quad\quad \textbf{B.} 14\quad\quad \textbf{C.} 23\quad\quad \textbf{D.} 11\par\vspace{2pt}
\textbf{\color{primary} Câu 31.} Một bài kiểm tra định kỳ nội dung Chương VII nêu câu hỏi: Cho hai đa thức M(x) = 3x$^2$ + 12x - 4 và N(x) = -x$^2$ + 13x + 4. Tổng M(x) + N(x) bằng. Hãy chọn đáp số đúng:\par
\hspace*{1.5em}\textbf{A.} 2x$^2$ - x\quad\quad \textbf{B.} 2x$^2$ + 25x - 8\quad\quad \textbf{C.} 2x$^2$ + 25x\quad\quad \textbf{D.} 4x$^2$ + 25x - 8\par\vspace{2pt}
\textbf{\color{primary} Câu 32.} Tìm lời giải đúng cho bài toán thông hiểu: Thực hiện phép nhân đơn thức với đa thức: 2x $\cdot$ (x$^2$ - 14x + 1) ta được:\par
\hspace*{1.5em}\textbf{A.} 3x$^3$ - 28x$^2$ + 2x\quad\quad \textbf{B.} 2x$^2$ - 28x + 2\quad\quad \textbf{C.} 2x$^3$ - 28x$^2$ + 2x\quad\quad \textbf{D.} 2x$^3$ - 14x$^2$ + 2x\par\vspace{2pt}
\textbf{\color{primary} Câu 33.} Sử dụng mối liên hệ giữa các đại lượng trong Chương VII: Giá trị của biểu thức đại số A(x) = x$^2$ - 2x + 17 tại x = 3 bằng:\par
\hspace*{1.5em}\textbf{A.} 14\quad\quad \textbf{B.} 26\quad\quad \textbf{C.} 17\quad\quad \textbf{D.} 20\par\vspace{2pt}
\textbf{\color{primary} Câu 34.} Hãy tính toán và chọn phương án đúng đối với bài tập: Cho hai đa thức M(x) = 3x$^2$ + 15x - 4 và N(x) = -x$^2$ + 16x + 4. Tổng M(x) + N(x) bằng:\par
\hspace*{1.5em}\textbf{A.} 4x$^2$ + 31x - 8\quad\quad \textbf{B.} 2x$^2$ + 31x - 8\quad\quad \textbf{C.} 2x$^2$ - x\quad\quad \textbf{D.} 2x$^2$ + 31x\par\vspace{2pt}
\textbf{\color{primary} Câu 35.} Dựa trên các tính chất trọng tâm của Chương VII, xét bài toán: Thực hiện phép nhân đơn thức với đa thức: 2x $\cdot$ (x$^2$ - 17x + 1) ta được:\par
\hspace*{1.5em}\textbf{A.} 2x$^2$ - 34x + 2\quad\quad \textbf{B.} 2x$^3$ - 17x$^2$ + 2x\quad\quad \textbf{C.} 2x$^3$ - 34x$^2$ + 2x\quad\quad \textbf{D.} 3x$^3$ - 34x$^2$ + 2x\par\vspace{2pt}
\textbf{\color{primary} Câu 36.} Thực hành giải toán (Biểu thức đại số và đa thức một biến): Giá trị của biểu thức đại số A(x) = x$^2$ - 2x + 20 tại x = 3 bằng. Đáp án nào dưới đây thỏa mãn?\par
\hspace*{1.5em}\textbf{A.} 20\quad\quad \textbf{B.} 29\quad\quad \textbf{C.} 23\quad\quad \textbf{D.} 17\par\vspace{2pt}
\textbf{\color{primary} Câu 37.} Kiểm tra kỹ năng biến đổi và suy luận qua bài toán: Cho hai đa thức M(x) = 3x$^2$ + 18x - 4 và N(x) = -x$^2$ + 19x + 4. Tổng M(x) + N(x) bằng:\par
\hspace*{1.5em}\textbf{A.} 2x$^2$ + 37x\quad\quad \textbf{B.} 4x$^2$ + 37x - 8\quad\quad \textbf{C.} 2x$^2$ + 37x - 8\quad\quad \textbf{D.} 2x$^2$ - x\par\vspace{2pt}
\textbf{\color{primary} Câu 38.} Cho giả thiết thuộc phạm vi Chương VII: Thực hiện phép nhân đơn thức với đa thức: 2x $\cdot$ (x$^2$ - 20x + 1) ta được. Kết luận đúng là:\par
\hspace*{1.5em}\textbf{A.} 2x$^2$ - 40x + 2\quad\quad \textbf{B.} 2x$^3$ - 40x$^2$ + 2x\quad\quad \textbf{C.} 3x$^3$ - 40x$^2$ + 2x\quad\quad \textbf{D.} 2x$^3$ - 20x$^2$ + 2x\par\vspace{2pt}
\textbf{\color{primary} Câu 39.} Rút gọn hoặc tính giá trị theo yêu cầu sau của Chương VII: Giá trị của biểu thức đại số A(x) = x$^2$ - 2x + 23 tại x = 3 bằng:\par
\hspace*{1.5em}\textbf{A.} 32\quad\quad \textbf{B.} 20\quad\quad \textbf{C.} 26\quad\quad \textbf{D.} 23\par\vspace{2pt}
\textbf{\color{primary} Câu 40.} Xác định phương án đúng sau khi thực hiện phép tính / suy luận: Cho hai đa thức M(x) = 3x$^2$ + 21x - 4 và N(x) = -x$^2$ + 22x + 4. Tổng M(x) + N(x) bằng:\par
\hspace*{1.5em}\textbf{A.} 4x$^2$ + 43x - 8\quad\quad \textbf{B.} 2x$^2$ + 43x - 8\quad\quad \textbf{C.} 2x$^2$ + 43x\quad\quad \textbf{D.} 2x$^2$ - x\par\vspace{2pt}
\subsection*{\color{accent} Phần III. Mức độ Khó (Vận dụng \& Vận dụng cao -- Câu 41 đến Câu 50)}
\textbf{\color{primary} Câu 41.} Tìm hệ số m để đa thức P(x) = x$^2$ - 4x + m có nghiệm là x = 2.:\par
\hspace*{1.5em}\textbf{A.} m = 12\quad\quad \textbf{B.} m = -4\quad\quad \textbf{C.} m = 4 (1)\quad\quad \textbf{D.} m = 4\par\vspace{2pt}
\textbf{\color{primary} Câu 42.} Bài toán vận dụng (Chương VII – Biểu thức đại số và đa thức một biến): Số dư trong phép chia đa thức A(x) = x$^3$ - 2x$^2$ + 3x + 6 cho nhị thức (x - 1) bằng. Kết quả đúng là:\par
\hspace*{1.5em}\textbf{A.} 2\quad\quad \textbf{B.} 6\quad\quad \textbf{C.} 8\quad\quad \textbf{D.} 4\par\vspace{2pt}
\textbf{\color{primary} Câu 43.} Vận dụng kiến thức nâng cao trong phạm vi Chương VII để giải bài toán: Tìm hệ số m để đa thức P(x) = x$^2$ - 6x + m có nghiệm là x = 2.:\par
\hspace*{1.5em}\textbf{A.} m = 8\quad\quad \textbf{B.} m = 16\quad\quad \textbf{C.} m = -8\quad\quad \textbf{D.} m = 6\par\vspace{2pt}
\textbf{\color{primary} Câu 44.} Phân tích và giải bài toán tổng hợp sau của Chương VII: Số dư trong phép chia đa thức A(x) = x$^3$ - 2x$^2$ + 3x + 8 cho nhị thức (x - 1) bằng:\par
\hspace*{1.5em}\textbf{A.} 6\quad\quad \textbf{B.} 8\quad\quad \textbf{C.} 4\quad\quad \textbf{D.} 10\par\vspace{2pt}
\textbf{\color{primary} Câu 45.} Thử thách tư duy trong phạm vi kiến thức Chương VII: Tìm hệ số m để đa thức P(x) = x$^2$ - 8x + m có nghiệm là x = 2.. Đáp số chính xác là:\par
\hspace*{1.5em}\textbf{A.} m = 20\quad\quad \textbf{B.} m = 12\quad\quad \textbf{C.} m = 8\quad\quad \textbf{D.} m = -12\par\vspace{2pt}
\textbf{\color{primary} Câu 46.} Sử dụng phương pháp giải toán của Chương VII để tìm đáp án cho bài toán: Số dư trong phép chia đa thức A(x) = x$^3$ - 2x$^2$ + 3x + 10 cho nhị thức (x - 1) bằng:\par
\hspace*{1.5em}\textbf{A.} 12\quad\quad \textbf{B.} 8\quad\quad \textbf{C.} 10\quad\quad \textbf{D.} 6\par\vspace{2pt}
\textbf{\color{primary} Câu 47.} Xét bài toán vận dụng thực tiễn / nâng cao thuộc Chương VII: Tìm hệ số m để đa thức P(x) = x$^2$ - 10x + m có nghiệm là x = 2.:\par
\hspace*{1.5em}\textbf{A.} m = 16\quad\quad \textbf{B.} m = -16\quad\quad \textbf{C.} m = 10\quad\quad \textbf{D.} m = 24\par\vspace{2pt}
\textbf{\color{primary} Câu 48.} Tìm kết quả cuối cùng của bài toán phân loại học sinh khá giỏi (Chương VII): Số dư trong phép chia đa thức A(x) = x$^3$ - 2x$^2$ + 3x + 12 cho nhị thức (x - 1) bằng:\par
\hspace*{1.5em}\textbf{A.} 14\quad\quad \textbf{B.} 10\quad\quad \textbf{C.} 12\quad\quad \textbf{D.} 8\par\vspace{2pt}
\textbf{\color{primary} Câu 49.} Kết hợp các tính chất đã học trong Chương VII, hãy giải quyết bài toán: Tìm hệ số m để đa thức P(x) = x$^2$ - 12x + m có nghiệm là x = 2.:\par
\hspace*{1.5em}\textbf{A.} m = -20\quad\quad \textbf{B.} m = 20\quad\quad \textbf{C.} m = 12\quad\quad \textbf{D.} m = 28\par\vspace{2pt}
\textbf{\color{primary} Câu 50.} Câu hỏi vận dụng cao cuối Chương VII: Số dư trong phép chia đa thức A(x) = x$^3$ - 2x$^2$ + 3x + 14 cho nhị thức (x - 1) bằng. Phương án đúng là:\par
\hspace*{1.5em}\textbf{A.} 14\quad\quad \textbf{B.} 16\quad\quad \textbf{C.} 10\quad\quad \textbf{D.} 12\par\vspace{2pt}
\subsection*{\color{primary} Bảng đáp án Chương VII: Biểu thức đại số và đa thức một biến}
\begin{center}
\renewcommand{\arraystretch}{1.25}
\begin{tabularx}{\textwidth}{*{10}{>{\centering\arraybackslash}X}}
\toprule
\textbf{\color{primary} Cột 1} & \textbf{\color{primary} Cột 2} & \textbf{\color{primary} Cột 3} & \textbf{\color{primary} Cột 4} & \textbf{\color{primary} Cột 5} & \textbf{\color{primary} Cột 6} & \textbf{\color{primary} Cột 7} & \textbf{\color{primary} Cột 8} & \textbf{\color{primary} Cột 9} & \textbf{\color{primary} Cột 10} \\
\midrule
\textbf{1.A} & \textbf{2.D} & \textbf{3.C} & \textbf{4.D} & \textbf{5.B} & \textbf{6.B} & \textbf{7.A} & \textbf{8.A} & \textbf{9.C} & \textbf{10.C} \\
\rowcolor{rowtint} \textbf{11.D} & \textbf{12.B} & \textbf{13.D} & \textbf{14.D} & \textbf{15.D} & \textbf{16.C} & \textbf{17.B} & \textbf{18.C} & \textbf{19.B} & \textbf{20.B} \\
\textbf{21.A} & \textbf{22.B} & \textbf{23.A} & \textbf{24.C} & \textbf{25.B} & \textbf{26.C} & \textbf{27.C} & \textbf{28.B} & \textbf{29.C} & \textbf{30.A} \\
\rowcolor{rowtint} \textbf{31.C} & \textbf{32.C} & \textbf{33.D} & \textbf{34.D} & \textbf{35.C} & \textbf{36.C} & \textbf{37.A} & \textbf{38.B} & \textbf{39.C} & \textbf{40.C} \\
\textbf{41.D} & \textbf{42.C} & \textbf{43.A} & \textbf{44.D} & \textbf{45.B} & \textbf{46.A} & \textbf{47.A} & \textbf{48.A} & \textbf{49.B} & \textbf{50.B} \\
\bottomrule
\end{tabularx}
\end{center}\vspace{10pt}
\section*{\color{primary} CHƯƠNG VIII: LÀM QUEN VỚI BIẾN CỐ VÀ XÁC SUẤT CỦA BIẾN CỐ}
\textit{Trọng tâm kiến thức SGK: Biến cố ngẫu nhiên, biến cố chắc chắn, biến cố không thể và xác suất của biến cố trong một số trò chơi đơn giản.}\par\vspace{4pt}
\subsection*{\color{accent} Phần I. Mức độ Dễ (Nhận biết -- Câu 1 đến Câu 20)}
\textbf{\color{primary} Câu 1.} Xác suất của biến cố chắc chắn luôn bằng:\par
\hspace*{1.5em}\textbf{A.} 1\quad\quad \textbf{B.} 0\quad\quad \textbf{C.} 1/2\quad\quad \textbf{D.} 100\par\vspace{2pt}
\textbf{\color{primary} Câu 2.} Thực hiện yêu cầu nhận biết sau: Xác suất của biến cố không thể luôn bằng.\par
\hspace*{1.5em}\textbf{A.} 0,5\quad\quad \textbf{B.} 1/6\quad\quad \textbf{C.} 0\quad\quad \textbf{D.} 1\par\vspace{2pt}
\textbf{\color{primary} Câu 3.} Trong bài học thuộc Chương VIII, xét nội dung: Gieo một con xúc xắc 6 mặt cân đối và đồng chất. Xác suất của biến cố 'Số chấm xuất hiện là số chẵn' bằng. Đáp án đúng là:\par
\hspace*{1.5em}\textbf{A.} 2/3\quad\quad \textbf{B.} 1/2\quad\quad \textbf{C.} 1/6\quad\quad \textbf{D.} 1/3\par\vspace{2pt}
\textbf{\color{primary} Câu 4.} Dựa vào kiến thức chuẩn SGK của Chương VIII (Làm quen với biến cố và xác suất của biến cố), trong một hộp có 8 chiếc bút bi xanh được đánh số từ 1 đến 8. Lấy ngẫu nhiên 1 chiếc bút, biến cố 'Lấy được chiếc bút bi màu xanh' là:\par
\begin{itemize}\setlength{\itemsep}{1pt}\setlength{\parskip}{0pt}
  \item[] \textbf{A.} Biến cố ngẫu nhiên có xác suất 1/2
  \item[] \textbf{B.} Biến cố chắc chắn (có xác suất bằng 1)
  \item[] \textbf{C.} Biến cố không thể (có xác suất bằng 0)
  \item[] \textbf{D.} Biến cố có xác suất 1/8
\end{itemize}\vspace{2pt}
\textbf{\color{primary} Câu 5.} Chọn kết quả đúng nhất khi xét: Một hộp có 8 tấm thẻ ghi các số tự nhiên liên tiếp từ 1 đến 8. Rút ngẫu nhiên 1 tấm thẻ, xác suất để rút được tấm thẻ ghi số 1 là:\par
\hspace*{1.5em}\textbf{A.} 2/8\quad\quad \textbf{B.} 1/8\quad\quad \textbf{C.} 0\quad\quad \textbf{D.} 1\par\vspace{2pt}
\textbf{\color{primary} Câu 6.} Khi làm bài tập cơ bản về làm quen với biến cố và xác suất của biến cố, một học sinh gặp câu hỏi: Xác suất của biến cố chắc chắn luôn bằng. Kết quả nào dưới đây đúng?\par
\hspace*{1.5em}\textbf{A.} 100\quad\quad \textbf{B.} 0\quad\quad \textbf{C.} 1/2\quad\quad \textbf{D.} 1\par\vspace{2pt}
\textbf{\color{primary} Câu 7.} Xác định giá trị hoặc khẳng định đúng cho bài toán sau: Xác suất của biến cố không thể luôn bằng:\par
\hspace*{1.5em}\textbf{A.} 1\quad\quad \textbf{B.} 1/6\quad\quad \textbf{C.} 0\quad\quad \textbf{D.} 0,5\par\vspace{2pt}
\textbf{\color{primary} Câu 8.} Áp dụng trực tiếp định nghĩa và quy tắc trong Chương VIII: Gieo một con xúc xắc 6 mặt cân đối và đồng chất. Xác suất của biến cố 'Số chấm xuất hiện là số chẵn' bằng:\par
\hspace*{1.5em}\textbf{A.} 1/3\quad\quad \textbf{B.} 2/3\quad\quad \textbf{C.} 1/2\quad\quad \textbf{D.} 1/6\par\vspace{2pt}
\textbf{\color{primary} Câu 9.} Kiểm tra nhanh kiến thức nền tảng: Trong một hộp có 13 chiếc bút bi xanh được đánh số từ 1 đến 13. Lấy ngẫu nhiên 1 chiếc bút, biến cố 'Lấy được chiếc bút bi màu xanh' là:\par
\begin{itemize}\setlength{\itemsep}{1pt}\setlength{\parskip}{0pt}
  \item[] \textbf{A.} Biến cố chắc chắn (có xác suất bằng 1)
  \item[] \textbf{B.} Biến cố có xác suất 1/13
  \item[] \textbf{C.} Biến cố không thể (có xác suất bằng 0)
  \item[] \textbf{D.} Biến cố ngẫu nhiên có xác suất 1/2
\end{itemize}\vspace{2pt}
\textbf{\color{primary} Câu 10.} Cho dữ kiện nhận biết sau đây: Một hộp có 13 tấm thẻ ghi các số tự nhiên liên tiếp từ 1 đến 13. Rút ngẫu nhiên 1 tấm thẻ, xác suất để rút được tấm thẻ ghi số 1 là. Phương án chính xác là:\par
\hspace*{1.5em}\textbf{A.} 2/13\quad\quad \textbf{B.} 0\quad\quad \textbf{C.} 1/13\quad\quad \textbf{D.} 1\par\vspace{2pt}
\textbf{\color{primary} Câu 11.} Hãy chỉ ra kết quả đúng trong bốn lựa chọn A, B, C, D đối với yêu cầu: Xác suất của biến cố chắc chắn luôn bằng:\par
\hspace*{1.5em}\textbf{A.} 1\quad\quad \textbf{B.} 100\quad\quad \textbf{C.} 1/2\quad\quad \textbf{D.} 0\par\vspace{2pt}
\textbf{\color{primary} Câu 12.} Câu hỏi ôn tập lý thuyết và kỹ năng cơ bản (Chương VIII): Xác suất của biến cố không thể luôn bằng:\par
\hspace*{1.5em}\textbf{A.} 0,5\quad\quad \textbf{B.} 0\quad\quad \textbf{C.} 1\quad\quad \textbf{D.} 1/6\par\vspace{2pt}
\textbf{\color{primary} Câu 13.} Theo nội dung sách giáo khoa Kết nối tri thức với cuộc sống (Chương VIII), gieo một con xúc xắc 6 mặt cân đối và đồng chất. Xác suất của biến cố 'Số chấm xuất hiện là số chẵn' bằng:\par
\hspace*{1.5em}\textbf{A.} 1/6\quad\quad \textbf{B.} 1/3\quad\quad \textbf{C.} 2/3\quad\quad \textbf{D.} 1/2\par\vspace{2pt}
\textbf{\color{primary} Câu 14.} Tìm đáp số đúng của bài toán cơ bản sau: Trong một hộp có 18 chiếc bút bi xanh được đánh số từ 1 đến 18. Lấy ngẫu nhiên 1 chiếc bút, biến cố 'Lấy được chiếc bút bi màu xanh' là:\par
\begin{itemize}\setlength{\itemsep}{1pt}\setlength{\parskip}{0pt}
  \item[] \textbf{A.} Biến cố chắc chắn (có xác suất bằng 1)
  \item[] \textbf{B.} Biến cố có xác suất 1/18
  \item[] \textbf{C.} Biến cố ngẫu nhiên có xác suất 1/2
  \item[] \textbf{D.} Biến cố không thể (có xác suất bằng 0)
\end{itemize}\vspace{2pt}
\textbf{\color{primary} Câu 15.} Xét đối tượng toán học trong Chương VIII: Một hộp có 18 tấm thẻ ghi các số tự nhiên liên tiếp từ 1 đến 18. Rút ngẫu nhiên 1 tấm thẻ, xác suất để rút được tấm thẻ ghi số 1 là. Giá trị hoặc tính chất tương ứng là:\par
\hspace*{1.5em}\textbf{A.} 1/18\quad\quad \textbf{B.} 1\quad\quad \textbf{C.} 2/18\quad\quad \textbf{D.} 0\par\vspace{2pt}
\textbf{\color{primary} Câu 16.} Hoàn thành phép xác định cơ bản sau đây: Xác suất của biến cố chắc chắn luôn bằng:\par
\hspace*{1.5em}\textbf{A.} 100\quad\quad \textbf{B.} 0\quad\quad \textbf{C.} 1/2\quad\quad \textbf{D.} 1\par\vspace{2pt}
\textbf{\color{primary} Câu 17.} Trong các khẳng định hoặc giá trị dưới đây, đâu là đáp án đúng cho câu hỏi: Xác suất của biến cố không thể luôn bằng?\par
\hspace*{1.5em}\textbf{A.} 0\quad\quad \textbf{B.} 0,5\quad\quad \textbf{C.} 1/6\quad\quad \textbf{D.} 1\par\vspace{2pt}
\textbf{\color{primary} Câu 18.} Bài tập khởi động Chương VIII: Gieo một con xúc xắc 6 mặt cân đối và đồng chất. Xác suất của biến cố 'Số chấm xuất hiện là số chẵn' bằng:\par
\hspace*{1.5em}\textbf{A.} 2/3\quad\quad \textbf{B.} 1/2\quad\quad \textbf{C.} 1/3\quad\quad \textbf{D.} 1/6\par\vspace{2pt}
\textbf{\color{primary} Câu 19.} Quan sát dữ kiện đề bài: Trong một hộp có 23 chiếc bút bi xanh được đánh số từ 1 đến 23. Lấy ngẫu nhiên 1 chiếc bút, biến cố 'Lấy được chiếc bút bi màu xanh' là. Hãy chọn phương án đúng:\par
\begin{itemize}\setlength{\itemsep}{1pt}\setlength{\parskip}{0pt}
  \item[] \textbf{A.} Biến cố không thể (có xác suất bằng 0)
  \item[] \textbf{B.} Biến cố chắc chắn (có xác suất bằng 1)
  \item[] \textbf{C.} Biến cố có xác suất 1/23
  \item[] \textbf{D.} Biến cố ngẫu nhiên có xác suất 1/2
\end{itemize}\vspace{2pt}
\textbf{\color{primary} Câu 20.} Sử dụng công thức hoặc tính chất đầu bài của Chương VIII để giải quyết yêu cầu: Một hộp có 23 tấm thẻ ghi các số tự nhiên liên tiếp từ 1 đến 23. Rút ngẫu nhiên 1 tấm thẻ, xác suất để rút được tấm thẻ ghi số 1 là:\par
\hspace*{1.5em}\textbf{A.} 1/23\quad\quad \textbf{B.} 0\quad\quad \textbf{C.} 2/23\quad\quad \textbf{D.} 1\par\vspace{2pt}
\subsection*{\color{accent} Phần II. Mức độ Trung bình (Thông hiểu -- Câu 21 đến Câu 40)}
\textbf{\color{primary} Câu 21.} Một hộp có 10 chiếc thẻ cùng loại được đánh số từ 1 đến 10. Rút ngẫu nhiên một thẻ. Xác suất của biến cố 'Số ghi trên thẻ là số chẵn' bằng:\par
\hspace*{1.5em}\textbf{A.} 1/10\quad\quad \textbf{B.} 4/10\quad\quad \textbf{C.} 1/3\quad\quad \textbf{D.} 1/2\par\vspace{2pt}
\textbf{\color{primary} Câu 22.} Vận dụng quy tắc trong Chương VIII (Làm quen với biến cố và xác suất của biến cố) để giải bài toán: Một túi đựng 4 viên bi xanh và 6 viên bi đỏ có cùng kích thước và khối lượng. Lấy ngẫu nhiên 1 viên bi. Xác suất để lấy được viên bi xanh là:\par
\hspace*{1.5em}\textbf{A.} 6/10\quad\quad \textbf{B.} 4/10\quad\quad \textbf{C.} 4/6\quad\quad \textbf{D.} 1/4\par\vspace{2pt}
\textbf{\color{primary} Câu 23.} Biến đổi và tính toán theo kiến thức Chương VIII: Gieo một con xúc xắc cân đối. Xác suất của biến cố 'Số chấm xuất hiện là số nguyên tố' (gồm 2, 3, 5) bằng. Kết quả thu được là:\par
\hspace*{1.5em}\textbf{A.} 2/3\quad\quad \textbf{B.} 1/2\quad\quad \textbf{C.} 1/3\quad\quad \textbf{D.} 1/6\par\vspace{2pt}
\textbf{\color{primary} Câu 24.} Xét bài toán mức độ thông hiểu sau: Một hộp có 16 chiếc thẻ cùng loại được đánh số từ 1 đến 16. Rút ngẫu nhiên một thẻ. Xác suất của biến cố 'Số ghi trên thẻ là số chẵn' bằng. Đáp án đúng là:\par
\hspace*{1.5em}\textbf{A.} 7/16\quad\quad \textbf{B.} 1/16\quad\quad \textbf{C.} 1/3\quad\quad \textbf{D.} 1/2\par\vspace{2pt}
\textbf{\color{primary} Câu 25.} Để hoàn thành phiếu học tập Chương VIII, hãy xác định: Một túi đựng 7 viên bi xanh và 9 viên bi đỏ có cùng kích thước và khối lượng. Lấy ngẫu nhiên 1 viên bi. Xác suất để lấy được viên bi xanh là:\par
\hspace*{1.5em}\textbf{A.} 7/9\quad\quad \textbf{B.} 7/16\quad\quad \textbf{C.} 1/7\quad\quad \textbf{D.} 9/16\par\vspace{2pt}
\textbf{\color{primary} Câu 26.} Phân tích dữ kiện bài toán: Gieo một con xúc xắc cân đối. Xác suất của biến cố 'Số chấm xuất hiện là số nguyên tố' (gồm 2, 3, 5) bằng. Giá trị cần tìm bằng:\par
\hspace*{1.5em}\textbf{A.} 1/3\quad\quad \textbf{B.} 2/3\quad\quad \textbf{C.} 1/2\quad\quad \textbf{D.} 1/6\par\vspace{2pt}
\textbf{\color{primary} Câu 27.} Áp dụng tính chất đã học ở Chương VIII vào bài tập sau: Một hộp có 22 chiếc thẻ cùng loại được đánh số từ 1 đến 22. Rút ngẫu nhiên một thẻ. Xác suất của biến cố 'Số ghi trên thẻ là số chẵn' bằng:\par
\hspace*{1.5em}\textbf{A.} 1/22\quad\quad \textbf{B.} 1/2\quad\quad \textbf{C.} 1/3\quad\quad \textbf{D.} 10/22\par\vspace{2pt}
\textbf{\color{primary} Câu 28.} Thực hiện các bước suy luận hợp lý để giải quyết yêu cầu: Một túi đựng 10 viên bi xanh và 12 viên bi đỏ có cùng kích thước và khối lượng. Lấy ngẫu nhiên 1 viên bi. Xác suất để lấy được viên bi xanh là:\par
\hspace*{1.5em}\textbf{A.} 10/12\quad\quad \textbf{B.} 12/22\quad\quad \textbf{C.} 10/22\quad\quad \textbf{D.} 1/10\par\vspace{2pt}
\textbf{\color{primary} Câu 29.} Cho bài toán tự luyện (Làm quen với biến cố và xác suất của biến cố): Gieo một con xúc xắc cân đối. Xác suất của biến cố 'Số chấm xuất hiện là số nguyên tố' (gồm 2, 3, 5) bằng. Phương án nào sau đây chính xác?\par
\hspace*{1.5em}\textbf{A.} 1/2\quad\quad \textbf{B.} 2/3\quad\quad \textbf{C.} 1/6\quad\quad \textbf{D.} 1/3\par\vspace{2pt}
\textbf{\color{primary} Câu 30.} Kết hợp các phép biến đổi trong Chương VIII, hãy tìm kết quả của bài toán: Một hộp có 28 chiếc thẻ cùng loại được đánh số từ 1 đến 28. Rút ngẫu nhiên một thẻ. Xác suất của biến cố 'Số ghi trên thẻ là số chẵn' bằng:\par
\hspace*{1.5em}\textbf{A.} 1/2\quad\quad \textbf{B.} 13/28\quad\quad \textbf{C.} 1/3\quad\quad \textbf{D.} 1/28\par\vspace{2pt}
\textbf{\color{primary} Câu 31.} Một bài kiểm tra định kỳ nội dung Chương VIII nêu câu hỏi: Một túi đựng 13 viên bi xanh và 15 viên bi đỏ có cùng kích thước và khối lượng. Lấy ngẫu nhiên 1 viên bi. Xác suất để lấy được viên bi xanh là. Hãy chọn đáp số đúng:\par
\hspace*{1.5em}\textbf{A.} 15/28\quad\quad \textbf{B.} 13/15\quad\quad \textbf{C.} 1/13\quad\quad \textbf{D.} 13/28\par\vspace{2pt}
\textbf{\color{primary} Câu 32.} Tìm lời giải đúng cho bài toán thông hiểu: Gieo một con xúc xắc cân đối. Xác suất của biến cố 'Số chấm xuất hiện là số nguyên tố' (gồm 2, 3, 5) bằng:\par
\hspace*{1.5em}\textbf{A.} 2/3\quad\quad \textbf{B.} 1/2\quad\quad \textbf{C.} 1/3\quad\quad \textbf{D.} 1/6\par\vspace{2pt}
\textbf{\color{primary} Câu 33.} Sử dụng mối liên hệ giữa các đại lượng trong Chương VIII: Một hộp có 34 chiếc thẻ cùng loại được đánh số từ 1 đến 34. Rút ngẫu nhiên một thẻ. Xác suất của biến cố 'Số ghi trên thẻ là số chẵn' bằng:\par
\hspace*{1.5em}\textbf{A.} 1/34\quad\quad \textbf{B.} 16/34\quad\quad \textbf{C.} 1/3\quad\quad \textbf{D.} 1/2\par\vspace{2pt}
\textbf{\color{primary} Câu 34.} Hãy tính toán và chọn phương án đúng đối với bài tập: Một túi đựng 16 viên bi xanh và 18 viên bi đỏ có cùng kích thước và khối lượng. Lấy ngẫu nhiên 1 viên bi. Xác suất để lấy được viên bi xanh là:\par
\hspace*{1.5em}\textbf{A.} 1/16\quad\quad \textbf{B.} 18/34\quad\quad \textbf{C.} 16/34\quad\quad \textbf{D.} 16/18\par\vspace{2pt}
\textbf{\color{primary} Câu 35.} Dựa trên các tính chất trọng tâm của Chương VIII, xét bài toán: Gieo một con xúc xắc cân đối. Xác suất của biến cố 'Số chấm xuất hiện là số nguyên tố' (gồm 2, 3, 5) bằng:\par
\hspace*{1.5em}\textbf{A.} 1/6\quad\quad \textbf{B.} 1/2\quad\quad \textbf{C.} 2/3\quad\quad \textbf{D.} 1/3\par\vspace{2pt}
\textbf{\color{primary} Câu 36.} Thực hành giải toán (Làm quen với biến cố và xác suất của biến cố): Một hộp có 40 chiếc thẻ cùng loại được đánh số từ 1 đến 40. Rút ngẫu nhiên một thẻ. Xác suất của biến cố 'Số ghi trên thẻ là số chẵn' bằng. Đáp án nào dưới đây thỏa mãn?\par
\hspace*{1.5em}\textbf{A.} 1/3\quad\quad \textbf{B.} 1/2\quad\quad \textbf{C.} 1/40\quad\quad \textbf{D.} 19/40\par\vspace{2pt}
\textbf{\color{primary} Câu 37.} Kiểm tra kỹ năng biến đổi và suy luận qua bài toán: Một túi đựng 19 viên bi xanh và 21 viên bi đỏ có cùng kích thước và khối lượng. Lấy ngẫu nhiên 1 viên bi. Xác suất để lấy được viên bi xanh là:\par
\hspace*{1.5em}\textbf{A.} 19/40\quad\quad \textbf{B.} 1/19\quad\quad \textbf{C.} 21/40\quad\quad \textbf{D.} 19/21\par\vspace{2pt}
\textbf{\color{primary} Câu 38.} Cho giả thiết thuộc phạm vi Chương VIII: Gieo một con xúc xắc cân đối. Xác suất của biến cố 'Số chấm xuất hiện là số nguyên tố' (gồm 2, 3, 5) bằng. Kết luận đúng là:\par
\hspace*{1.5em}\textbf{A.} 1/6\quad\quad \textbf{B.} 1/3\quad\quad \textbf{C.} 2/3\quad\quad \textbf{D.} 1/2\par\vspace{2pt}
\textbf{\color{primary} Câu 39.} Rút gọn hoặc tính giá trị theo yêu cầu sau của Chương VIII: Một hộp có 46 chiếc thẻ cùng loại được đánh số từ 1 đến 46. Rút ngẫu nhiên một thẻ. Xác suất của biến cố 'Số ghi trên thẻ là số chẵn' bằng:\par
\hspace*{1.5em}\textbf{A.} 1/46\quad\quad \textbf{B.} 1/3\quad\quad \textbf{C.} 1/2\quad\quad \textbf{D.} 22/46\par\vspace{2pt}
\textbf{\color{primary} Câu 40.} Xác định phương án đúng sau khi thực hiện phép tính / suy luận: Một túi đựng 22 viên bi xanh và 24 viên bi đỏ có cùng kích thước và khối lượng. Lấy ngẫu nhiên 1 viên bi. Xác suất để lấy được viên bi xanh là:\par
\hspace*{1.5em}\textbf{A.} 22/24\quad\quad \textbf{B.} 24/46\quad\quad \textbf{C.} 1/22\quad\quad \textbf{D.} 22/46\par\vspace{2pt}
\subsection*{\color{accent} Phần III. Mức độ Khó (Vận dụng \& Vận dụng cao -- Câu 41 đến Câu 50)}
\textbf{\color{primary} Câu 41.} Từ các thẻ ghi các số tự nhiên từ 1 đến 12, rút ngẫu nhiên 1 thẻ. Xác suất của biến cố 'Số ghi trên thẻ chia hết cho 3' bằng:\par
\hspace*{1.5em}\textbf{A.} 1/4\quad\quad \textbf{B.} 1/3\quad\quad \textbf{C.} 1/2\quad\quad \textbf{D.} 2/3\par\vspace{2pt}
\textbf{\color{primary} Câu 42.} Bài toán vận dụng (Chương VIII – Làm quen với biến cố và xác suất của biến cố): Chọn ngẫu nhiên một số tự nhiên có hai chữ số (từ 10 đến 99). Xác suất để chọn được một số chính phương (gồm 16, 25, 36, 49, 64, 81) bằng. Kết quả đúng là:\par
\hspace*{1.5em}\textbf{A.} 6/99\quad\quad \textbf{B.} 1/10\quad\quad \textbf{C.} 1/12\quad\quad \textbf{D.} 1/15\par\vspace{2pt}
\textbf{\color{primary} Câu 43.} Vận dụng kiến thức nâng cao trong phạm vi Chương VIII để giải bài toán: Từ các thẻ ghi các số tự nhiên từ 1 đến 18, rút ngẫu nhiên 1 thẻ. Xác suất của biến cố 'Số ghi trên thẻ chia hết cho 3' bằng:\par
\hspace*{1.5em}\textbf{A.} 1/4\quad\quad \textbf{B.} 2/3\quad\quad \textbf{C.} 1/3\quad\quad \textbf{D.} 1/2\par\vspace{2pt}
\textbf{\color{primary} Câu 44.} Phân tích và giải bài toán tổng hợp sau của Chương VIII: Chọn ngẫu nhiên một số tự nhiên có hai chữ số (từ 10 đến 99). Xác suất để chọn được một số chính phương (gồm 16, 25, 36, 49, 64, 81) bằng:\par
\hspace*{1.5em}\textbf{A.} 6/99\quad\quad \textbf{B.} 1/12\quad\quad \textbf{C.} 1/10\quad\quad \textbf{D.} 1/15\par\vspace{2pt}
\textbf{\color{primary} Câu 45.} Thử thách tư duy trong phạm vi kiến thức Chương VIII: Từ các thẻ ghi các số tự nhiên từ 1 đến 12, rút ngẫu nhiên 1 thẻ. Xác suất của biến cố 'Số ghi trên thẻ chia hết cho 3' bằng. Đáp số chính xác là:\par
\hspace*{1.5em}\textbf{A.} 1/2\quad\quad \textbf{B.} 2/3\quad\quad \textbf{C.} 1/3\quad\quad \textbf{D.} 1/4\par\vspace{2pt}
\textbf{\color{primary} Câu 46.} Sử dụng phương pháp giải toán của Chương VIII để tìm đáp án cho bài toán: Chọn ngẫu nhiên một số tự nhiên có hai chữ số (từ 10 đến 99). Xác suất để chọn được một số chính phương (gồm 16, 25, 36, 49, 64, 81) bằng:\par
\hspace*{1.5em}\textbf{A.} 1/15\quad\quad \textbf{B.} 1/12\quad\quad \textbf{C.} 6/99\quad\quad \textbf{D.} 1/10\par\vspace{2pt}
\textbf{\color{primary} Câu 47.} Xét bài toán vận dụng thực tiễn / nâng cao thuộc Chương VIII: Từ các thẻ ghi các số tự nhiên từ 1 đến 18, rút ngẫu nhiên 1 thẻ. Xác suất của biến cố 'Số ghi trên thẻ chia hết cho 3' bằng:\par
\hspace*{1.5em}\textbf{A.} 1/3\quad\quad \textbf{B.} 2/3\quad\quad \textbf{C.} 1/2\quad\quad \textbf{D.} 1/4\par\vspace{2pt}
\textbf{\color{primary} Câu 48.} Tìm kết quả cuối cùng của bài toán phân loại học sinh khá giỏi (Chương VIII): Chọn ngẫu nhiên một số tự nhiên có hai chữ số (từ 10 đến 99). Xác suất để chọn được một số chính phương (gồm 16, 25, 36, 49, 64, 81) bằng:\par
\hspace*{1.5em}\textbf{A.} 1/10\quad\quad \textbf{B.} 1/12\quad\quad \textbf{C.} 6/99\quad\quad \textbf{D.} 1/15\par\vspace{2pt}
\textbf{\color{primary} Câu 49.} Kết hợp các tính chất đã học trong Chương VIII, hãy giải quyết bài toán: Từ các thẻ ghi các số tự nhiên từ 1 đến 12, rút ngẫu nhiên 1 thẻ. Xác suất của biến cố 'Số ghi trên thẻ chia hết cho 3' bằng:\par
\hspace*{1.5em}\textbf{A.} 2/3\quad\quad \textbf{B.} 1/2\quad\quad \textbf{C.} 1/3\quad\quad \textbf{D.} 1/4\par\vspace{2pt}
\textbf{\color{primary} Câu 50.} Câu hỏi vận dụng cao cuối Chương VIII: Chọn ngẫu nhiên một số tự nhiên có hai chữ số (từ 10 đến 99). Xác suất để chọn được một số chính phương (gồm 16, 25, 36, 49, 64, 81) bằng. Phương án đúng là:\par
\hspace*{1.5em}\textbf{A.} 1/12\quad\quad \textbf{B.} 1/10\quad\quad \textbf{C.} 6/99\quad\quad \textbf{D.} 1/15\par\vspace{2pt}
\subsection*{\color{primary} Bảng đáp án Chương VIII: Làm quen với biến cố và xác suất của biến cố}
\begin{center}
\renewcommand{\arraystretch}{1.25}
\begin{tabularx}{\textwidth}{*{10}{>{\centering\arraybackslash}X}}
\toprule
\textbf{\color{primary} Cột 1} & \textbf{\color{primary} Cột 2} & \textbf{\color{primary} Cột 3} & \textbf{\color{primary} Cột 4} & \textbf{\color{primary} Cột 5} & \textbf{\color{primary} Cột 6} & \textbf{\color{primary} Cột 7} & \textbf{\color{primary} Cột 8} & \textbf{\color{primary} Cột 9} & \textbf{\color{primary} Cột 10} \\
\midrule
\textbf{1.A} & \textbf{2.C} & \textbf{3.B} & \textbf{4.B} & \textbf{5.B} & \textbf{6.D} & \textbf{7.C} & \textbf{8.C} & \textbf{9.A} & \textbf{10.C} \\
\rowcolor{rowtint} \textbf{11.A} & \textbf{12.B} & \textbf{13.D} & \textbf{14.A} & \textbf{15.A} & \textbf{16.D} & \textbf{17.A} & \textbf{18.B} & \textbf{19.B} & \textbf{20.A} \\
\textbf{21.D} & \textbf{22.B} & \textbf{23.B} & \textbf{24.D} & \textbf{25.B} & \textbf{26.C} & \textbf{27.B} & \textbf{28.C} & \textbf{29.A} & \textbf{30.A} \\
\rowcolor{rowtint} \textbf{31.D} & \textbf{32.B} & \textbf{33.D} & \textbf{34.C} & \textbf{35.B} & \textbf{36.B} & \textbf{37.A} & \textbf{38.D} & \textbf{39.C} & \textbf{40.D} \\
\textbf{41.B} & \textbf{42.D} & \textbf{43.C} & \textbf{44.D} & \textbf{45.C} & \textbf{46.A} & \textbf{47.A} & \textbf{48.D} & \textbf{49.C} & \textbf{50.D} \\
\bottomrule
\end{tabularx}
\end{center}\vspace{10pt}
\section*{\color{primary} CHƯƠNG IX: QUAN HỆ GIỮA CÁC YẾU TỐ TRONG MỘT TAM GIÁC}
\textit{Trọng tâm kiến thức SGK: Quan hệ giữa góc và cạnh đối diện trong một tam giác, đường vuông góc và đường xiên, bất đẳng thức tam giác, sự đồng quy của ba đường trung tuyến, ba đường phân giác, ba đường trung trực, ba đường cao.}\par\vspace{4pt}
\subsection*{\color{accent} Phần I. Mức độ Dễ (Nhận biết -- Câu 1 đến Câu 20)}
\textbf{\color{primary} Câu 1.} Cho tam giác ABC có cạnh AB = 4 cm, BC = 6 cm, AC = 8 cm. Góc lớn nhất của tam giác ABC là:\par
\begin{itemize}\setlength{\itemsep}{1pt}\setlength{\parskip}{0pt}
  \item[] \textbf{A.} Góc C
  \item[] \textbf{B.} Góc B (đối diện với cạnh lớn nhất AC)
  \item[] \textbf{C.} Góc A
  \item[] \textbf{D.} Ba góc bằng nhau
\end{itemize}\vspace{2pt}
\textbf{\color{primary} Câu 2.} Thực hiện yêu cầu nhận biết sau: Trong một tam giác, giao điểm của ba đường trung tuyến được gọi là.\par
\begin{itemize}\setlength{\itemsep}{1pt}\setlength{\parskip}{0pt}
  \item[] \textbf{A.} Tâm đường tròn nội tiếp
  \item[] \textbf{B.} Trọng tâm của tam giác
  \item[] \textbf{C.} Tâm đường tròn ngoại tiếp
  \item[] \textbf{D.} Trực tâm của tam giác
\end{itemize}\vspace{2pt}
\textbf{\color{primary} Câu 3.} Trong bài học thuộc Chương IX, xét nội dung: Cho tam giác ABC có đường trung tuyến AM = 15 cm và G là trọng tâm của tam giác. Độ dài đoạn thẳng AG bằng. Đáp án đúng là:\par
\hspace*{1.5em}\textbf{A.} 15 cm\quad\quad \textbf{B.} 10 cm\quad\quad \textbf{C.} 7.5 cm\quad\quad \textbf{D.} 5 cm\par\vspace{2pt}
\textbf{\color{primary} Câu 4.} Dựa vào kiến thức chuẩn SGK của Chương IX (Quan hệ giữa các yếu tố trong một tam giác), cho tam giác ABC có cạnh AB = 7 cm, BC = 9 cm, AC = 11 cm. Góc lớn nhất của tam giác ABC là:\par
\begin{itemize}\setlength{\itemsep}{1pt}\setlength{\parskip}{0pt}
  \item[] \textbf{A.} Góc B (đối diện với cạnh lớn nhất AC)
  \item[] \textbf{B.} Góc A
  \item[] \textbf{C.} Góc C
  \item[] \textbf{D.} Ba góc bằng nhau
\end{itemize}\vspace{2pt}
\textbf{\color{primary} Câu 5.} Chọn kết quả đúng nhất khi xét: Trong một tam giác, giao điểm của ba đường trung tuyến được gọi là:\par
\begin{itemize}\setlength{\itemsep}{1pt}\setlength{\parskip}{0pt}
  \item[] \textbf{A.} Tâm đường tròn nội tiếp
  \item[] \textbf{B.} Trọng tâm của tam giác
  \item[] \textbf{C.} Tâm đường tròn ngoại tiếp
  \item[] \textbf{D.} Trực tâm của tam giác
\end{itemize}\vspace{2pt}
\textbf{\color{primary} Câu 6.} Khi làm bài tập cơ bản về quan hệ giữa các yếu tố trong một tam giác, một học sinh gặp câu hỏi: Cho tam giác ABC có đường trung tuyến AM = 24 cm và G là trọng tâm của tam giác. Độ dài đoạn thẳng AG bằng. Kết quả nào dưới đây đúng?\par
\hspace*{1.5em}\textbf{A.} 16 cm\quad\quad \textbf{B.} 24 cm\quad\quad \textbf{C.} 8 cm\quad\quad \textbf{D.} 12.0 cm\par\vspace{2pt}
\textbf{\color{primary} Câu 7.} Xác định giá trị hoặc khẳng định đúng cho bài toán sau: Cho tam giác ABC có cạnh AB = 10 cm, BC = 12 cm, AC = 14 cm. Góc lớn nhất của tam giác ABC là:\par
\begin{itemize}\setlength{\itemsep}{1pt}\setlength{\parskip}{0pt}
  \item[] \textbf{A.} Ba góc bằng nhau
  \item[] \textbf{B.} Góc C
  \item[] \textbf{C.} Góc B (đối diện với cạnh lớn nhất AC)
  \item[] \textbf{D.} Góc A
\end{itemize}\vspace{2pt}
\textbf{\color{primary} Câu 8.} Áp dụng trực tiếp định nghĩa và quy tắc trong Chương IX: Trong một tam giác, giao điểm của ba đường trung tuyến được gọi là:\par
\begin{itemize}\setlength{\itemsep}{1pt}\setlength{\parskip}{0pt}
  \item[] \textbf{A.} Tâm đường tròn nội tiếp
  \item[] \textbf{B.} Tâm đường tròn ngoại tiếp
  \item[] \textbf{C.} Trực tâm của tam giác
  \item[] \textbf{D.} Trọng tâm của tam giác
\end{itemize}\vspace{2pt}
\textbf{\color{primary} Câu 9.} Kiểm tra nhanh kiến thức nền tảng: Cho tam giác ABC có đường trung tuyến AM = 33 cm và G là trọng tâm của tam giác. Độ dài đoạn thẳng AG bằng:\par
\hspace*{1.5em}\textbf{A.} 33 cm\quad\quad \textbf{B.} 16.5 cm\quad\quad \textbf{C.} 11 cm\quad\quad \textbf{D.} 22 cm\par\vspace{2pt}
\textbf{\color{primary} Câu 10.} Cho dữ kiện nhận biết sau đây: Cho tam giác ABC có cạnh AB = 13 cm, BC = 15 cm, AC = 17 cm. Góc lớn nhất của tam giác ABC là. Phương án chính xác là:\par
\begin{itemize}\setlength{\itemsep}{1pt}\setlength{\parskip}{0pt}
  \item[] \textbf{A.} Góc A
  \item[] \textbf{B.} Góc B (đối diện với cạnh lớn nhất AC)
  \item[] \textbf{C.} Góc C
  \item[] \textbf{D.} Ba góc bằng nhau
\end{itemize}\vspace{2pt}
\textbf{\color{primary} Câu 11.} Hãy chỉ ra kết quả đúng trong bốn lựa chọn A, B, C, D đối với yêu cầu: Trong một tam giác, giao điểm của ba đường trung tuyến được gọi là:\par
\begin{itemize}\setlength{\itemsep}{1pt}\setlength{\parskip}{0pt}
  \item[] \textbf{A.} Tâm đường tròn ngoại tiếp
  \item[] \textbf{B.} Trọng tâm của tam giác
  \item[] \textbf{C.} Trực tâm của tam giác
  \item[] \textbf{D.} Tâm đường tròn nội tiếp
\end{itemize}\vspace{2pt}
\textbf{\color{primary} Câu 12.} Câu hỏi ôn tập lý thuyết và kỹ năng cơ bản (Chương IX): Cho tam giác ABC có đường trung tuyến AM = 42 cm và G là trọng tâm của tam giác. Độ dài đoạn thẳng AG bằng:\par
\hspace*{1.5em}\textbf{A.} 28 cm\quad\quad \textbf{B.} 42 cm\quad\quad \textbf{C.} 21.0 cm\quad\quad \textbf{D.} 14 cm\par\vspace{2pt}
\textbf{\color{primary} Câu 13.} Theo nội dung sách giáo khoa Kết nối tri thức với cuộc sống (Chương IX), cho tam giác ABC có cạnh AB = 16 cm, BC = 18 cm, AC = 20 cm. Góc lớn nhất của tam giác ABC là:\par
\begin{itemize}\setlength{\itemsep}{1pt}\setlength{\parskip}{0pt}
  \item[] \textbf{A.} Góc C
  \item[] \textbf{B.} Ba góc bằng nhau
  \item[] \textbf{C.} Góc A
  \item[] \textbf{D.} Góc B (đối diện với cạnh lớn nhất AC)
\end{itemize}\vspace{2pt}
\textbf{\color{primary} Câu 14.} Tìm đáp số đúng của bài toán cơ bản sau: Trong một tam giác, giao điểm của ba đường trung tuyến được gọi là:\par
\begin{itemize}\setlength{\itemsep}{1pt}\setlength{\parskip}{0pt}
  \item[] \textbf{A.} Trực tâm của tam giác
  \item[] \textbf{B.} Tâm đường tròn ngoại tiếp
  \item[] \textbf{C.} Trọng tâm của tam giác
  \item[] \textbf{D.} Tâm đường tròn nội tiếp
\end{itemize}\vspace{2pt}
\textbf{\color{primary} Câu 15.} Xét đối tượng toán học trong Chương IX: Cho tam giác ABC có đường trung tuyến AM = 51 cm và G là trọng tâm của tam giác. Độ dài đoạn thẳng AG bằng. Giá trị hoặc tính chất tương ứng là:\par
\hspace*{1.5em}\textbf{A.} 34 cm\quad\quad \textbf{B.} 17 cm\quad\quad \textbf{C.} 25.5 cm\quad\quad \textbf{D.} 51 cm\par\vspace{2pt}
\textbf{\color{primary} Câu 16.} Hoàn thành phép xác định cơ bản sau đây: Cho tam giác ABC có cạnh AB = 19 cm, BC = 21 cm, AC = 23 cm. Góc lớn nhất của tam giác ABC là:\par
\begin{itemize}\setlength{\itemsep}{1pt}\setlength{\parskip}{0pt}
  \item[] \textbf{A.} Góc A
  \item[] \textbf{B.} Góc C
  \item[] \textbf{C.} Ba góc bằng nhau
  \item[] \textbf{D.} Góc B (đối diện với cạnh lớn nhất AC)
\end{itemize}\vspace{2pt}
\textbf{\color{primary} Câu 17.} Trong các khẳng định hoặc giá trị dưới đây, đâu là đáp án đúng cho câu hỏi: Trong một tam giác, giao điểm của ba đường trung tuyến được gọi là?\par
\begin{itemize}\setlength{\itemsep}{1pt}\setlength{\parskip}{0pt}
  \item[] \textbf{A.} Trọng tâm của tam giác
  \item[] \textbf{B.} Tâm đường tròn ngoại tiếp
  \item[] \textbf{C.} Tâm đường tròn nội tiếp
  \item[] \textbf{D.} Trực tâm của tam giác
\end{itemize}\vspace{2pt}
\textbf{\color{primary} Câu 18.} Bài tập khởi động Chương IX: Cho tam giác ABC có đường trung tuyến AM = 60 cm và G là trọng tâm của tam giác. Độ dài đoạn thẳng AG bằng:\par
\hspace*{1.5em}\textbf{A.} 60 cm\quad\quad \textbf{B.} 30.0 cm\quad\quad \textbf{C.} 40 cm\quad\quad \textbf{D.} 20 cm\par\vspace{2pt}
\textbf{\color{primary} Câu 19.} Quan sát dữ kiện đề bài: Cho tam giác ABC có cạnh AB = 22 cm, BC = 24 cm, AC = 26 cm. Góc lớn nhất của tam giác ABC là. Hãy chọn phương án đúng:\par
\begin{itemize}\setlength{\itemsep}{1pt}\setlength{\parskip}{0pt}
  \item[] \textbf{A.} Góc C
  \item[] \textbf{B.} Góc B (đối diện với cạnh lớn nhất AC)
  \item[] \textbf{C.} Ba góc bằng nhau
  \item[] \textbf{D.} Góc A
\end{itemize}\vspace{2pt}
\textbf{\color{primary} Câu 20.} Sử dụng công thức hoặc tính chất đầu bài của Chương IX để giải quyết yêu cầu: Trong một tam giác, giao điểm của ba đường trung tuyến được gọi là:\par
\begin{itemize}\setlength{\itemsep}{1pt}\setlength{\parskip}{0pt}
  \item[] \textbf{A.} Tâm đường tròn nội tiếp
  \item[] \textbf{B.} Tâm đường tròn ngoại tiếp
  \item[] \textbf{C.} Trọng tâm của tam giác
  \item[] \textbf{D.} Trực tâm của tam giác
\end{itemize}\vspace{2pt}
\subsection*{\color{accent} Phần II. Mức độ Trung bình (Thông hiểu -- Câu 21 đến Câu 40)}
\textbf{\color{primary} Câu 21.} Bộ ba độ dài đoạn thẳng nào dưới đây có thể là độ dài ba cạnh của một tam giác?\par
\hspace*{1.5em}\textbf{A.} 1 cm; 2 cm; 4 cm\quad\quad \textbf{B.} 2 cm; 4 cm; 7 cm\quad\quad \textbf{C.} 5 cm; 6 cm; 8 cm\quad\quad \textbf{D.} 2 cm; 3 cm; 5 cm\par\vspace{2pt}
\textbf{\color{primary} Câu 22.} Vận dụng quy tắc trong Chương IX (Quan hệ giữa các yếu tố trong một tam giác) để giải bài toán: Một tam giác cân có độ dài hai cạnh là 3 cm và 11 cm. Chu vi của tam giác cân đó bằng:\par
\hspace*{1.5em}\textbf{A.} 22 cm\quad\quad \textbf{B.} 17 cm\quad\quad \textbf{C.} 14 cm\quad\quad \textbf{D.} 25 cm\par\vspace{2pt}
\textbf{\color{primary} Câu 23.} Biến đổi và tính toán theo kiến thức Chương IX: Trong một tam giác, điểm cách đều ba cạnh của tam giác là giao điểm của. Kết quả thu được là:\par
\begin{itemize}\setlength{\itemsep}{1pt}\setlength{\parskip}{0pt}
  \item[] \textbf{A.} Ba đường trung trực của tam giác
  \item[] \textbf{B.} Ba đường trung tuyến của tam giác
  \item[] \textbf{C.} Ba đường phân giác của tam giác
  \item[] \textbf{D.} Ba đường cao của tam giác
\end{itemize}\vspace{2pt}
\textbf{\color{primary} Câu 24.} Xét bài toán mức độ thông hiểu sau: Bộ ba độ dài đoạn thẳng nào dưới đây có thể là độ dài ba cạnh của một tam giác. Đáp án đúng là:\par
\hspace*{1.5em}\textbf{A.} 1 cm; 2 cm; 4 cm\quad\quad \textbf{B.} 8 cm; 9 cm; 11 cm\quad\quad \textbf{C.} 2 cm; 4 cm; 10 cm\quad\quad \textbf{D.} 5 cm; 6 cm; 11 cm\par\vspace{2pt}
\textbf{\color{primary} Câu 25.} Để hoàn thành phiếu học tập Chương IX, hãy xác định: Một tam giác cân có độ dài hai cạnh là 3 cm và 17 cm. Chu vi của tam giác cân đó bằng:\par
\hspace*{1.5em}\textbf{A.} 23 cm\quad\quad \textbf{B.} 20 cm\quad\quad \textbf{C.} 37 cm\quad\quad \textbf{D.} 34 cm\par\vspace{2pt}
\textbf{\color{primary} Câu 26.} Phân tích dữ kiện bài toán: Trong một tam giác, điểm cách đều ba cạnh của tam giác là giao điểm của. Giá trị cần tìm bằng:\par
\begin{itemize}\setlength{\itemsep}{1pt}\setlength{\parskip}{0pt}
  \item[] \textbf{A.} Ba đường phân giác của tam giác
  \item[] \textbf{B.} Ba đường trung trực của tam giác
  \item[] \textbf{C.} Ba đường trung tuyến của tam giác
  \item[] \textbf{D.} Ba đường cao của tam giác
\end{itemize}\vspace{2pt}
\textbf{\color{primary} Câu 27.} Áp dụng tính chất đã học ở Chương IX vào bài tập sau: Bộ ba độ dài đoạn thẳng nào dưới đây có thể là độ dài ba cạnh của một tam giác:\par
\hspace*{1.5em}\textbf{A.} 1 cm; 2 cm; 4 cm\quad\quad \textbf{B.} 11 cm; 12 cm; 14 cm\quad\quad \textbf{C.} 2 cm; 4 cm; 13 cm\quad\quad \textbf{D.} 8 cm; 9 cm; 17 cm\par\vspace{2pt}
\textbf{\color{primary} Câu 28.} Thực hiện các bước suy luận hợp lý để giải quyết yêu cầu: Một tam giác cân có độ dài hai cạnh là 3 cm và 23 cm. Chu vi của tam giác cân đó bằng:\par
\hspace*{1.5em}\textbf{A.} 46 cm\quad\quad \textbf{B.} 49 cm\quad\quad \textbf{C.} 29 cm\quad\quad \textbf{D.} 26 cm\par\vspace{2pt}
\textbf{\color{primary} Câu 29.} Cho bài toán tự luyện (Quan hệ giữa các yếu tố trong một tam giác): Trong một tam giác, điểm cách đều ba cạnh của tam giác là giao điểm của. Phương án nào sau đây chính xác?\par
\begin{itemize}\setlength{\itemsep}{1pt}\setlength{\parskip}{0pt}
  \item[] \textbf{A.} Ba đường trung tuyến của tam giác
  \item[] \textbf{B.} Ba đường cao của tam giác
  \item[] \textbf{C.} Ba đường phân giác của tam giác
  \item[] \textbf{D.} Ba đường trung trực của tam giác
\end{itemize}\vspace{2pt}
\textbf{\color{primary} Câu 30.} Kết hợp các phép biến đổi trong Chương IX, hãy tìm kết quả của bài toán: Bộ ba độ dài đoạn thẳng nào dưới đây có thể là độ dài ba cạnh của một tam giác:\par
\hspace*{1.5em}\textbf{A.} 14 cm; 15 cm; 17 cm\quad\quad \textbf{B.} 1 cm; 2 cm; 4 cm\quad\quad \textbf{C.} 2 cm; 4 cm; 16 cm\quad\quad \textbf{D.} 11 cm; 12 cm; 23 cm\par\vspace{2pt}
\textbf{\color{primary} Câu 31.} Một bài kiểm tra định kỳ nội dung Chương IX nêu câu hỏi: Một tam giác cân có độ dài hai cạnh là 3 cm và 29 cm. Chu vi của tam giác cân đó bằng. Hãy chọn đáp số đúng:\par
\hspace*{1.5em}\textbf{A.} 61 cm\quad\quad \textbf{B.} 32 cm\quad\quad \textbf{C.} 58 cm\quad\quad \textbf{D.} 35 cm\par\vspace{2pt}
\textbf{\color{primary} Câu 32.} Tìm lời giải đúng cho bài toán thông hiểu: Trong một tam giác, điểm cách đều ba cạnh của tam giác là giao điểm của:\par
\begin{itemize}\setlength{\itemsep}{1pt}\setlength{\parskip}{0pt}
  \item[] \textbf{A.} Ba đường phân giác của tam giác
  \item[] \textbf{B.} Ba đường cao của tam giác
  \item[] \textbf{C.} Ba đường trung trực của tam giác
  \item[] \textbf{D.} Ba đường trung tuyến của tam giác
\end{itemize}\vspace{2pt}
\textbf{\color{primary} Câu 33.} Sử dụng mối liên hệ giữa các đại lượng trong Chương IX: Bộ ba độ dài đoạn thẳng nào dưới đây có thể là độ dài ba cạnh của một tam giác:\par
\hspace*{1.5em}\textbf{A.} 2 cm; 4 cm; 19 cm\quad\quad \textbf{B.} 14 cm; 15 cm; 29 cm\quad\quad \textbf{C.} 17 cm; 18 cm; 20 cm\quad\quad \textbf{D.} 1 cm; 2 cm; 4 cm\par\vspace{2pt}
\textbf{\color{primary} Câu 34.} Hãy tính toán và chọn phương án đúng đối với bài tập: Một tam giác cân có độ dài hai cạnh là 3 cm và 35 cm. Chu vi của tam giác cân đó bằng:\par
\hspace*{1.5em}\textbf{A.} 41 cm\quad\quad \textbf{B.} 38 cm\quad\quad \textbf{C.} 70 cm\quad\quad \textbf{D.} 73 cm\par\vspace{2pt}
\textbf{\color{primary} Câu 35.} Dựa trên các tính chất trọng tâm của Chương IX, xét bài toán: Trong một tam giác, điểm cách đều ba cạnh của tam giác là giao điểm của:\par
\begin{itemize}\setlength{\itemsep}{1pt}\setlength{\parskip}{0pt}
  \item[] \textbf{A.} Ba đường trung trực của tam giác
  \item[] \textbf{B.} Ba đường phân giác của tam giác
  \item[] \textbf{C.} Ba đường cao của tam giác
  \item[] \textbf{D.} Ba đường trung tuyến của tam giác
\end{itemize}\vspace{2pt}
\textbf{\color{primary} Câu 36.} Thực hành giải toán (Quan hệ giữa các yếu tố trong một tam giác): Bộ ba độ dài đoạn thẳng nào dưới đây có thể là độ dài ba cạnh của một tam giác. Đáp án nào dưới đây thỏa mãn?\par
\hspace*{1.5em}\textbf{A.} 2 cm; 4 cm; 22 cm\quad\quad \textbf{B.} 20 cm; 21 cm; 23 cm\quad\quad \textbf{C.} 1 cm; 2 cm; 4 cm\quad\quad \textbf{D.} 17 cm; 18 cm; 35 cm\par\vspace{2pt}
\textbf{\color{primary} Câu 37.} Kiểm tra kỹ năng biến đổi và suy luận qua bài toán: Một tam giác cân có độ dài hai cạnh là 3 cm và 41 cm. Chu vi của tam giác cân đó bằng:\par
\hspace*{1.5em}\textbf{A.} 82 cm\quad\quad \textbf{B.} 44 cm\quad\quad \textbf{C.} 47 cm\quad\quad \textbf{D.} 85 cm\par\vspace{2pt}
\textbf{\color{primary} Câu 38.} Cho giả thiết thuộc phạm vi Chương IX: Trong một tam giác, điểm cách đều ba cạnh của tam giác là giao điểm của. Kết luận đúng là:\par
\begin{itemize}\setlength{\itemsep}{1pt}\setlength{\parskip}{0pt}
  \item[] \textbf{A.} Ba đường trung tuyến của tam giác
  \item[] \textbf{B.} Ba đường cao của tam giác
  \item[] \textbf{C.} Ba đường trung trực của tam giác
  \item[] \textbf{D.} Ba đường phân giác của tam giác
\end{itemize}\vspace{2pt}
\textbf{\color{primary} Câu 39.} Rút gọn hoặc tính giá trị theo yêu cầu sau của Chương IX: Bộ ba độ dài đoạn thẳng nào dưới đây có thể là độ dài ba cạnh của một tam giác:\par
\hspace*{1.5em}\textbf{A.} 1 cm; 2 cm; 4 cm\quad\quad \textbf{B.} 2 cm; 4 cm; 25 cm\quad\quad \textbf{C.} 23 cm; 24 cm; 26 cm\quad\quad \textbf{D.} 20 cm; 21 cm; 41 cm\par\vspace{2pt}
\textbf{\color{primary} Câu 40.} Xác định phương án đúng sau khi thực hiện phép tính / suy luận: Một tam giác cân có độ dài hai cạnh là 3 cm và 47 cm. Chu vi của tam giác cân đó bằng:\par
\hspace*{1.5em}\textbf{A.} 94 cm\quad\quad \textbf{B.} 53 cm\quad\quad \textbf{C.} 97 cm\quad\quad \textbf{D.} 50 cm\par\vspace{2pt}
\subsection*{\color{accent} Phần III. Mức độ Khó (Vận dụng \& Vận dụng cao -- Câu 41 đến Câu 50)}
\textbf{\color{primary} Câu 41.} Cho tam giác ABC có hai cạnh AB = 3 cm, AC = 7 cm. Có bao nhiêu giá trị nguyên của độ dài cạnh BC (tính theo cm)?\par
\hspace*{1.5em}\textbf{A.} 7 giá trị\quad\quad \textbf{B.} 5 giá trị\quad\quad \textbf{C.} 4 giá trị\quad\quad \textbf{D.} 6 giá trị\par\vspace{2pt}
\textbf{\color{primary} Câu 42.} Bài toán vận dụng (Chương IX – Quan hệ giữa các yếu tố trong một tam giác): Cho tam giác ABC có G là trọng tâm, đường trung tuyến BM. Biết khoảng cách từ G đến M bằng 4 cm. Độ dài cả đường trung tuyến BM bằng. Kết quả đúng là:\par
\hspace*{1.5em}\textbf{A.} 24 cm\quad\quad \textbf{B.} 12 cm\quad\quad \textbf{C.} 8 cm\quad\quad \textbf{D.} 16 cm\par\vspace{2pt}
\textbf{\color{primary} Câu 43.} Vận dụng kiến thức nâng cao trong phạm vi Chương IX để giải bài toán: Cho tam giác ABC có hai cạnh AB = 3 cm, AC = 9 cm. Có bao nhiêu giá trị nguyên của độ dài cạnh BC (tính theo cm):\par
\hspace*{1.5em}\textbf{A.} 6 giá trị\quad\quad \textbf{B.} 5 giá trị\quad\quad \textbf{C.} 4 giá trị\quad\quad \textbf{D.} 7 giá trị\par\vspace{2pt}
\textbf{\color{primary} Câu 44.} Phân tích và giải bài toán tổng hợp sau của Chương IX: Cho tam giác ABC có G là trọng tâm, đường trung tuyến BM. Biết khoảng cách từ G đến M bằng 6 cm. Độ dài cả đường trung tuyến BM bằng:\par
\hspace*{1.5em}\textbf{A.} 18 cm\quad\quad \textbf{B.} 24 cm\quad\quad \textbf{C.} 12 cm\quad\quad \textbf{D.} 36 cm\par\vspace{2pt}
\textbf{\color{primary} Câu 45.} Thử thách tư duy trong phạm vi kiến thức Chương IX: Cho tam giác ABC có hai cạnh AB = 3 cm, AC = 11 cm. Có bao nhiêu giá trị nguyên của độ dài cạnh BC (tính theo cm). Đáp số chính xác là:\par
\hspace*{1.5em}\textbf{A.} 5 giá trị\quad\quad \textbf{B.} 7 giá trị\quad\quad \textbf{C.} 4 giá trị\quad\quad \textbf{D.} 6 giá trị\par\vspace{2pt}
\textbf{\color{primary} Câu 46.} Sử dụng phương pháp giải toán của Chương IX để tìm đáp án cho bài toán: Cho tam giác ABC có G là trọng tâm, đường trung tuyến BM. Biết khoảng cách từ G đến M bằng 8 cm. Độ dài cả đường trung tuyến BM bằng:\par
\hspace*{1.5em}\textbf{A.} 32 cm\quad\quad \textbf{B.} 16 cm\quad\quad \textbf{C.} 48 cm\quad\quad \textbf{D.} 24 cm\par\vspace{2pt}
\textbf{\color{primary} Câu 47.} Xét bài toán vận dụng thực tiễn / nâng cao thuộc Chương IX: Cho tam giác ABC có hai cạnh AB = 3 cm, AC = 13 cm. Có bao nhiêu giá trị nguyên của độ dài cạnh BC (tính theo cm):\par
\hspace*{1.5em}\textbf{A.} 4 giá trị\quad\quad \textbf{B.} 5 giá trị\quad\quad \textbf{C.} 6 giá trị\quad\quad \textbf{D.} 7 giá trị\par\vspace{2pt}
\textbf{\color{primary} Câu 48.} Tìm kết quả cuối cùng của bài toán phân loại học sinh khá giỏi (Chương IX): Cho tam giác ABC có G là trọng tâm, đường trung tuyến BM. Biết khoảng cách từ G đến M bằng 10 cm. Độ dài cả đường trung tuyến BM bằng:\par
\hspace*{1.5em}\textbf{A.} 40 cm\quad\quad \textbf{B.} 60 cm\quad\quad \textbf{C.} 30 cm\quad\quad \textbf{D.} 20 cm\par\vspace{2pt}
\textbf{\color{primary} Câu 49.} Kết hợp các tính chất đã học trong Chương IX, hãy giải quyết bài toán: Cho tam giác ABC có hai cạnh AB = 3 cm, AC = 15 cm. Có bao nhiêu giá trị nguyên của độ dài cạnh BC (tính theo cm):\par
\hspace*{1.5em}\textbf{A.} 5 giá trị\quad\quad \textbf{B.} 7 giá trị\quad\quad \textbf{C.} 4 giá trị\quad\quad \textbf{D.} 6 giá trị\par\vspace{2pt}
\textbf{\color{primary} Câu 50.} Câu hỏi vận dụng cao cuối Chương IX: Cho tam giác ABC có G là trọng tâm, đường trung tuyến BM. Biết khoảng cách từ G đến M bằng 12 cm. Độ dài cả đường trung tuyến BM bằng. Phương án đúng là:\par
\hspace*{1.5em}\textbf{A.} 72 cm\quad\quad \textbf{B.} 48 cm\quad\quad \textbf{C.} 36 cm\quad\quad \textbf{D.} 24 cm\par\vspace{2pt}
\subsection*{\color{primary} Bảng đáp án Chương IX: Quan hệ giữa các yếu tố trong một tam giác}
\begin{center}
\renewcommand{\arraystretch}{1.25}
\begin{tabularx}{\textwidth}{*{10}{>{\centering\arraybackslash}X}}
\toprule
\textbf{\color{primary} Cột 1} & \textbf{\color{primary} Cột 2} & \textbf{\color{primary} Cột 3} & \textbf{\color{primary} Cột 4} & \textbf{\color{primary} Cột 5} & \textbf{\color{primary} Cột 6} & \textbf{\color{primary} Cột 7} & \textbf{\color{primary} Cột 8} & \textbf{\color{primary} Cột 9} & \textbf{\color{primary} Cột 10} \\
\midrule
\textbf{1.B} & \textbf{2.B} & \textbf{3.B} & \textbf{4.A} & \textbf{5.B} & \textbf{6.A} & \textbf{7.C} & \textbf{8.D} & \textbf{9.D} & \textbf{10.B} \\
\rowcolor{rowtint} \textbf{11.B} & \textbf{12.A} & \textbf{13.D} & \textbf{14.C} & \textbf{15.A} & \textbf{16.D} & \textbf{17.A} & \textbf{18.C} & \textbf{19.B} & \textbf{20.C} \\
\textbf{21.C} & \textbf{22.D} & \textbf{23.C} & \textbf{24.B} & \textbf{25.C} & \textbf{26.A} & \textbf{27.B} & \textbf{28.B} & \textbf{29.C} & \textbf{30.A} \\
\rowcolor{rowtint} \textbf{31.A} & \textbf{32.A} & \textbf{33.C} & \textbf{34.D} & \textbf{35.B} & \textbf{36.B} & \textbf{37.D} & \textbf{38.D} & \textbf{39.C} & \textbf{40.C} \\
\textbf{41.B} & \textbf{42.B} & \textbf{43.B} & \textbf{44.A} & \textbf{45.A} & \textbf{46.D} & \textbf{47.B} & \textbf{48.C} & \textbf{49.A} & \textbf{50.C} \\
\bottomrule
\end{tabularx}
\end{center}\vspace{10pt}
\section*{\color{primary} CHƯƠNG X: MỘT SỐ HÌNH KHỐI TRONG THỰC TIỄN}
\textit{Trọng tâm kiến thức SGK: Hình hộp chữ nhật, hình lập phương, hình lăng trụ đứng tam giác và hình lăng trụ đứng tứ giác (diện tích xung quanh và thể tích, không dùng hình vẽ).}\par\vspace{4pt}
\subsection*{\color{accent} Phần I. Mức độ Dễ (Nhận biết -- Câu 1 đến Câu 20)}
\textbf{\color{primary} Câu 1.} Một hình lập phương có độ dài cạnh bằng 3 cm. Thể tích của hình lập phương đó là:\par
\hspace*{1.5em}\textbf{A.} 27 cm$^3$\quad\quad \textbf{B.} 27 cm$^3$ (1)\quad\quad \textbf{C.} 36 cm$^3$\quad\quad \textbf{D.} 54 cm$^3$\par\vspace{2pt}
\textbf{\color{primary} Câu 2.} Thực hiện yêu cầu nhận biết sau: Một hình hộp chữ nhật có ba kích thước là 3 cm, 4 cm và 5 cm. Thể tích hình hộp chữ nhật đó bằng.\par
\hspace*{1.5em}\textbf{A.} 60 cm$^3$\quad\quad \textbf{B.} 70 cm$^3$\quad\quad \textbf{C.} 12 cm$^3$\quad\quad \textbf{D.} 24 cm$^3$\par\vspace{2pt}
\textbf{\color{primary} Câu 3.} Trong bài học thuộc Chương X, xét nội dung: Hình lăng trụ đứng tam giác có tất cả bao nhiêu đỉnh, bao nhiêu cạnh và bao nhiêu mặt. Đáp án đúng là:\par
\hspace*{1.5em}\textbf{A.} 6 đỉnh, 6 cạnh, 5 mặt\quad\quad \textbf{B.} 4 đỉnh, 6 cạnh, 4 mặt\quad\quad \textbf{C.} 6 đỉnh, 9 cạnh, 5 mặt\quad\quad \textbf{D.} 8 đỉnh, 12 cạnh, 6 mặt\par\vspace{2pt}
\textbf{\color{primary} Câu 4.} Dựa vào kiến thức chuẩn SGK của Chương X (Một số hình khối trong thực tiễn), một hình lập phương có độ dài cạnh bằng 6 cm. Thể tích của hình lập phương đó là:\par
\hspace*{1.5em}\textbf{A.} 216 cm$^3$ (1)\quad\quad \textbf{B.} 144 cm$^3$\quad\quad \textbf{C.} 216 cm$^3$\quad\quad \textbf{D.} 72 cm$^3$\par\vspace{2pt}
\textbf{\color{primary} Câu 5.} Chọn kết quả đúng nhất khi xét: Một hình hộp chữ nhật có ba kích thước là 6 cm, 7 cm và 8 cm. Thể tích hình hộp chữ nhật đó bằng:\par
\hspace*{1.5em}\textbf{A.} 21 cm$^3$\quad\quad \textbf{B.} 84 cm$^3$\quad\quad \textbf{C.} 208 cm$^3$\quad\quad \textbf{D.} 336 cm$^3$\par\vspace{2pt}
\textbf{\color{primary} Câu 6.} Khi làm bài tập cơ bản về một số hình khối trong thực tiễn, một học sinh gặp câu hỏi: Hình lăng trụ đứng tam giác có tất cả bao nhiêu đỉnh, bao nhiêu cạnh và bao nhiêu mặt. Kết quả nào dưới đây đúng?\par
\hspace*{1.5em}\textbf{A.} 4 đỉnh, 6 cạnh, 4 mặt\quad\quad \textbf{B.} 6 đỉnh, 9 cạnh, 5 mặt\quad\quad \textbf{C.} 6 đỉnh, 6 cạnh, 5 mặt\quad\quad \textbf{D.} 8 đỉnh, 12 cạnh, 6 mặt\par\vspace{2pt}
\textbf{\color{primary} Câu 7.} Xác định giá trị hoặc khẳng định đúng cho bài toán sau: Một hình lập phương có độ dài cạnh bằng 9 cm. Thể tích của hình lập phương đó là:\par
\hspace*{1.5em}\textbf{A.} 324 cm$^3$\quad\quad \textbf{B.} 486 cm$^3$\quad\quad \textbf{C.} 108 cm$^3$\quad\quad \textbf{D.} 729 cm$^3$\par\vspace{2pt}
\textbf{\color{primary} Câu 8.} Áp dụng trực tiếp định nghĩa và quy tắc trong Chương X: Một hình hộp chữ nhật có ba kích thước là 9 cm, 10 cm và 11 cm. Thể tích hình hộp chữ nhật đó bằng:\par
\hspace*{1.5em}\textbf{A.} 418 cm$^3$\quad\quad \textbf{B.} 990 cm$^3$\quad\quad \textbf{C.} 30 cm$^3$\quad\quad \textbf{D.} 180 cm$^3$\par\vspace{2pt}
\textbf{\color{primary} Câu 9.} Kiểm tra nhanh kiến thức nền tảng: Hình lăng trụ đứng tam giác có tất cả bao nhiêu đỉnh, bao nhiêu cạnh và bao nhiêu mặt?\par
\hspace*{1.5em}\textbf{A.} 6 đỉnh, 6 cạnh, 5 mặt\quad\quad \textbf{B.} 8 đỉnh, 12 cạnh, 6 mặt\quad\quad \textbf{C.} 6 đỉnh, 9 cạnh, 5 mặt\quad\quad \textbf{D.} 4 đỉnh, 6 cạnh, 4 mặt\par\vspace{2pt}
\textbf{\color{primary} Câu 10.} Cho dữ kiện nhận biết sau đây: Một hình lập phương có độ dài cạnh bằng 12 cm. Thể tích của hình lập phương đó là. Phương án chính xác là:\par
\hspace*{1.5em}\textbf{A.} 144 cm$^3$\quad\quad \textbf{B.} 576 cm$^3$\quad\quad \textbf{C.} 1728 cm$^3$\quad\quad \textbf{D.} 864 cm$^3$\par\vspace{2pt}
\textbf{\color{primary} Câu 11.} Hãy chỉ ra kết quả đúng trong bốn lựa chọn A, B, C, D đối với yêu cầu: Một hình hộp chữ nhật có ba kích thước là 12 cm, 13 cm và 14 cm. Thể tích hình hộp chữ nhật đó bằng:\par
\hspace*{1.5em}\textbf{A.} 700 cm$^3$\quad\quad \textbf{B.} 2184 cm$^3$\quad\quad \textbf{C.} 312 cm$^3$\quad\quad \textbf{D.} 39 cm$^3$\par\vspace{2pt}
\textbf{\color{primary} Câu 12.} Câu hỏi ôn tập lý thuyết và kỹ năng cơ bản (Chương X): Hình lăng trụ đứng tam giác có tất cả bao nhiêu đỉnh, bao nhiêu cạnh và bao nhiêu mặt:\par
\hspace*{1.5em}\textbf{A.} 8 đỉnh, 12 cạnh, 6 mặt\quad\quad \textbf{B.} 6 đỉnh, 9 cạnh, 5 mặt\quad\quad \textbf{C.} 6 đỉnh, 6 cạnh, 5 mặt\quad\quad \textbf{D.} 4 đỉnh, 6 cạnh, 4 mặt\par\vspace{2pt}
\textbf{\color{primary} Câu 13.} Theo nội dung sách giáo khoa Kết nối tri thức với cuộc sống (Chương X), một hình lập phương có độ dài cạnh bằng 15 cm. Thể tích của hình lập phương đó là:\par
\hspace*{1.5em}\textbf{A.} 3375 cm$^3$\quad\quad \textbf{B.} 180 cm$^3$\quad\quad \textbf{C.} 900 cm$^3$\quad\quad \textbf{D.} 1350 cm$^3$\par\vspace{2pt}
\textbf{\color{primary} Câu 14.} Tìm đáp số đúng của bài toán cơ bản sau: Một hình hộp chữ nhật có ba kích thước là 15 cm, 16 cm và 17 cm. Thể tích hình hộp chữ nhật đó bằng:\par
\hspace*{1.5em}\textbf{A.} 1054 cm$^3$\quad\quad \textbf{B.} 480 cm$^3$\quad\quad \textbf{C.} 4080 cm$^3$\quad\quad \textbf{D.} 48 cm$^3$\par\vspace{2pt}
\textbf{\color{primary} Câu 15.} Xét đối tượng toán học trong Chương X: Hình lăng trụ đứng tam giác có tất cả bao nhiêu đỉnh, bao nhiêu cạnh và bao nhiêu mặt. Giá trị hoặc tính chất tương ứng là:\par
\hspace*{1.5em}\textbf{A.} 6 đỉnh, 6 cạnh, 5 mặt\quad\quad \textbf{B.} 6 đỉnh, 9 cạnh, 5 mặt\quad\quad \textbf{C.} 8 đỉnh, 12 cạnh, 6 mặt\quad\quad \textbf{D.} 4 đỉnh, 6 cạnh, 4 mặt\par\vspace{2pt}
\textbf{\color{primary} Câu 16.} Hoàn thành phép xác định cơ bản sau đây: Một hình lập phương có độ dài cạnh bằng 18 cm. Thể tích của hình lập phương đó là:\par
\hspace*{1.5em}\textbf{A.} 216 cm$^3$\quad\quad \textbf{B.} 5832 cm$^3$\quad\quad \textbf{C.} 1296 cm$^3$\quad\quad \textbf{D.} 1944 cm$^3$\par\vspace{2pt}
\textbf{\color{primary} Câu 17.} Trong các khẳng định hoặc giá trị dưới đây, đâu là đáp án đúng cho câu hỏi: Một hình hộp chữ nhật có ba kích thước là 18 cm, 19 cm và 20 cm. Thể tích hình hộp chữ nhật đó bằng?\par
\hspace*{1.5em}\textbf{A.} 684 cm$^3$\quad\quad \textbf{B.} 6840 cm$^3$\quad\quad \textbf{C.} 1480 cm$^3$\quad\quad \textbf{D.} 57 cm$^3$\par\vspace{2pt}
\textbf{\color{primary} Câu 18.} Bài tập khởi động Chương X: Hình lăng trụ đứng tam giác có tất cả bao nhiêu đỉnh, bao nhiêu cạnh và bao nhiêu mặt:\par
\hspace*{1.5em}\textbf{A.} 4 đỉnh, 6 cạnh, 4 mặt\quad\quad \textbf{B.} 8 đỉnh, 12 cạnh, 6 mặt\quad\quad \textbf{C.} 6 đỉnh, 6 cạnh, 5 mặt\quad\quad \textbf{D.} 6 đỉnh, 9 cạnh, 5 mặt\par\vspace{2pt}
\textbf{\color{primary} Câu 19.} Quan sát dữ kiện đề bài: Một hình lập phương có độ dài cạnh bằng 21 cm. Thể tích của hình lập phương đó là. Hãy chọn phương án đúng:\par
\hspace*{1.5em}\textbf{A.} 252 cm$^3$\quad\quad \textbf{B.} 9261 cm$^3$\quad\quad \textbf{C.} 1764 cm$^3$\quad\quad \textbf{D.} 2646 cm$^3$\par\vspace{2pt}
\textbf{\color{primary} Câu 20.} Sử dụng công thức hoặc tính chất đầu bài của Chương X để giải quyết yêu cầu: Một hình hộp chữ nhật có ba kích thước là 21 cm, 22 cm và 23 cm. Thể tích hình hộp chữ nhật đó bằng:\par
\hspace*{1.5em}\textbf{A.} 1978 cm$^3$\quad\quad \textbf{B.} 66 cm$^3$\quad\quad \textbf{C.} 924 cm$^3$\quad\quad \textbf{D.} 10626 cm$^3$\par\vspace{2pt}
\subsection*{\color{accent} Phần II. Mức độ Trung bình (Thông hiểu -- Câu 21 đến Câu 40)}
\textbf{\color{primary} Câu 21.} Một hình lập phương có diện tích xung quanh bằng 16 cm$^2$. Độ dài cạnh của hình lập phương đó là:\par
\hspace*{1.5em}\textbf{A.} 3 cm\quad\quad \textbf{B.} 2 cm (1)\quad\quad \textbf{C.} 4 cm\quad\quad \textbf{D.} 2 cm\par\vspace{2pt}
\textbf{\color{primary} Câu 22.} Vận dụng quy tắc trong Chương X (Một số hình khối trong thực tiễn) để giải bài toán: Một hình lăng trụ đứng tam giác có chu vi đáy bằng 14 cm và chiều cao bằng 6 cm. Diện tích xung quanh của hình lăng trụ đứng đó là:\par
\hspace*{1.5em}\textbf{A.} 42 cm$^2$\quad\quad \textbf{B.} 168 cm$^2$\quad\quad \textbf{C.} 20 cm$^2$\quad\quad \textbf{D.} 84 cm$^2$\par\vspace{2pt}
\textbf{\color{primary} Câu 23.} Biến đổi và tính toán theo kiến thức Chương X: Một hình lăng trụ đứng có diện tích đáy bằng 21 dm$^2$ và chiều cao bằng 6 dm. Thể tích của hình lăng trụ đứng đó bằng. Kết quả thu được là:\par
\hspace*{1.5em}\textbf{A.} 252 dm$^3$\quad\quad \textbf{B.} 126 dm$^3$\quad\quad \textbf{C.} 27 dm$^3$\quad\quad \textbf{D.} 42 dm$^3$\par\vspace{2pt}
\textbf{\color{primary} Câu 24.} Xét bài toán mức độ thông hiểu sau: Một hình lập phương có diện tích xung quanh bằng 100 cm$^2$. Độ dài cạnh của hình lập phương đó là. Đáp án đúng là:\par
\hspace*{1.5em}\textbf{A.} 5 cm\quad\quad \textbf{B.} 6 cm\quad\quad \textbf{C.} 10 cm\quad\quad \textbf{D.} 25 cm\par\vspace{2pt}
\textbf{\color{primary} Câu 25.} Để hoàn thành phiếu học tập Chương X, hãy xác định: Một hình lăng trụ đứng tam giác có chu vi đáy bằng 20 cm và chiều cao bằng 9 cm. Diện tích xung quanh của hình lăng trụ đứng đó là:\par
\hspace*{1.5em}\textbf{A.} 180 cm$^2$\quad\quad \textbf{B.} 29 cm$^2$\quad\quad \textbf{C.} 360 cm$^2$\quad\quad \textbf{D.} 90 cm$^2$\par\vspace{2pt}
\textbf{\color{primary} Câu 26.} Phân tích dữ kiện bài toán: Một hình lăng trụ đứng có diện tích đáy bằng 30 dm$^2$ và chiều cao bằng 9 dm. Thể tích của hình lăng trụ đứng đó bằng. Giá trị cần tìm bằng:\par
\hspace*{1.5em}\textbf{A.} 90 dm$^3$\quad\quad \textbf{B.} 540 dm$^3$\quad\quad \textbf{C.} 270 dm$^3$\quad\quad \textbf{D.} 39 dm$^3$\par\vspace{2pt}
\textbf{\color{primary} Câu 27.} Áp dụng tính chất đã học ở Chương X vào bài tập sau: Một hình lập phương có diện tích xung quanh bằng 256 cm$^2$. Độ dài cạnh của hình lập phương đó là:\par
\hspace*{1.5em}\textbf{A.} 16 cm\quad\quad \textbf{B.} 9 cm\quad\quad \textbf{C.} 8 cm\quad\quad \textbf{D.} 64 cm\par\vspace{2pt}
\textbf{\color{primary} Câu 28.} Thực hiện các bước suy luận hợp lý để giải quyết yêu cầu: Một hình lăng trụ đứng tam giác có chu vi đáy bằng 26 cm và chiều cao bằng 12 cm. Diện tích xung quanh của hình lăng trụ đứng đó là:\par
\hspace*{1.5em}\textbf{A.} 156 cm$^2$\quad\quad \textbf{B.} 38 cm$^2$\quad\quad \textbf{C.} 624 cm$^2$\quad\quad \textbf{D.} 312 cm$^2$\par\vspace{2pt}
\textbf{\color{primary} Câu 29.} Cho bài toán tự luyện (Một số hình khối trong thực tiễn): Một hình lăng trụ đứng có diện tích đáy bằng 39 dm$^2$ và chiều cao bằng 12 dm. Thể tích của hình lăng trụ đứng đó bằng. Phương án nào sau đây chính xác?\par
\hspace*{1.5em}\textbf{A.} 51 dm$^3$\quad\quad \textbf{B.} 156 dm$^3$\quad\quad \textbf{C.} 468 dm$^3$\quad\quad \textbf{D.} 936 dm$^3$\par\vspace{2pt}
\textbf{\color{primary} Câu 30.} Kết hợp các phép biến đổi trong Chương X, hãy tìm kết quả của bài toán: Một hình lập phương có diện tích xung quanh bằng 484 cm$^2$. Độ dài cạnh của hình lập phương đó là:\par
\hspace*{1.5em}\textbf{A.} 11 cm\quad\quad \textbf{B.} 12 cm\quad\quad \textbf{C.} 121 cm\quad\quad \textbf{D.} 22 cm\par\vspace{2pt}
\textbf{\color{primary} Câu 31.} Một bài kiểm tra định kỳ nội dung Chương X nêu câu hỏi: Một hình lăng trụ đứng tam giác có chu vi đáy bằng 32 cm và chiều cao bằng 15 cm. Diện tích xung quanh của hình lăng trụ đứng đó là. Hãy chọn đáp số đúng:\par
\hspace*{1.5em}\textbf{A.} 960 cm$^2$\quad\quad \textbf{B.} 240 cm$^2$\quad\quad \textbf{C.} 47 cm$^2$\quad\quad \textbf{D.} 480 cm$^2$\par\vspace{2pt}
\textbf{\color{primary} Câu 32.} Tìm lời giải đúng cho bài toán thông hiểu: Một hình lăng trụ đứng có diện tích đáy bằng 48 dm$^2$ và chiều cao bằng 15 dm. Thể tích của hình lăng trụ đứng đó bằng:\par
\hspace*{1.5em}\textbf{A.} 240 dm$^3$\quad\quad \textbf{B.} 720 dm$^3$\quad\quad \textbf{C.} 63 dm$^3$\quad\quad \textbf{D.} 1440 dm$^3$\par\vspace{2pt}
\textbf{\color{primary} Câu 33.} Sử dụng mối liên hệ giữa các đại lượng trong Chương X: Một hình lập phương có diện tích xung quanh bằng 784 cm$^2$. Độ dài cạnh của hình lập phương đó là:\par
\hspace*{1.5em}\textbf{A.} 196 cm\quad\quad \textbf{B.} 28 cm\quad\quad \textbf{C.} 15 cm\quad\quad \textbf{D.} 14 cm\par\vspace{2pt}
\textbf{\color{primary} Câu 34.} Hãy tính toán và chọn phương án đúng đối với bài tập: Một hình lăng trụ đứng tam giác có chu vi đáy bằng 38 cm và chiều cao bằng 18 cm. Diện tích xung quanh của hình lăng trụ đứng đó là:\par
\hspace*{1.5em}\textbf{A.} 56 cm$^2$\quad\quad \textbf{B.} 342 cm$^2$\quad\quad \textbf{C.} 684 cm$^2$\quad\quad \textbf{D.} 1368 cm$^2$\par\vspace{2pt}
\textbf{\color{primary} Câu 35.} Dựa trên các tính chất trọng tâm của Chương X, xét bài toán: Một hình lăng trụ đứng có diện tích đáy bằng 57 dm$^2$ và chiều cao bằng 18 dm. Thể tích của hình lăng trụ đứng đó bằng:\par
\hspace*{1.5em}\textbf{A.} 75 dm$^3$\quad\quad \textbf{B.} 342 dm$^3$\quad\quad \textbf{C.} 1026 dm$^3$\quad\quad \textbf{D.} 2052 dm$^3$\par\vspace{2pt}
\textbf{\color{primary} Câu 36.} Thực hành giải toán (Một số hình khối trong thực tiễn): Một hình lập phương có diện tích xung quanh bằng 1156 cm$^2$. Độ dài cạnh của hình lập phương đó là. Đáp án nào dưới đây thỏa mãn?\par
\hspace*{1.5em}\textbf{A.} 289 cm\quad\quad \textbf{B.} 17 cm\quad\quad \textbf{C.} 18 cm\quad\quad \textbf{D.} 34 cm\par\vspace{2pt}
\textbf{\color{primary} Câu 37.} Kiểm tra kỹ năng biến đổi và suy luận qua bài toán: Một hình lăng trụ đứng tam giác có chu vi đáy bằng 44 cm và chiều cao bằng 21 cm. Diện tích xung quanh của hình lăng trụ đứng đó là:\par
\hspace*{1.5em}\textbf{A.} 1848 cm$^2$\quad\quad \textbf{B.} 462 cm$^2$\quad\quad \textbf{C.} 65 cm$^2$\quad\quad \textbf{D.} 924 cm$^2$\par\vspace{2pt}
\textbf{\color{primary} Câu 38.} Cho giả thiết thuộc phạm vi Chương X: Một hình lăng trụ đứng có diện tích đáy bằng 66 dm$^2$ và chiều cao bằng 21 dm. Thể tích của hình lăng trụ đứng đó bằng. Kết luận đúng là:\par
\hspace*{1.5em}\textbf{A.} 87 dm$^3$\quad\quad \textbf{B.} 462 dm$^3$\quad\quad \textbf{C.} 1386 dm$^3$\quad\quad \textbf{D.} 2772 dm$^3$\par\vspace{2pt}
\textbf{\color{primary} Câu 39.} Rút gọn hoặc tính giá trị theo yêu cầu sau của Chương X: Một hình lập phương có diện tích xung quanh bằng 1600 cm$^2$. Độ dài cạnh của hình lập phương đó là:\par
\hspace*{1.5em}\textbf{A.} 20 cm\quad\quad \textbf{B.} 40 cm\quad\quad \textbf{C.} 21 cm\quad\quad \textbf{D.} 400 cm\par\vspace{2pt}
\textbf{\color{primary} Câu 40.} Xác định phương án đúng sau khi thực hiện phép tính / suy luận: Một hình lăng trụ đứng tam giác có chu vi đáy bằng 50 cm và chiều cao bằng 24 cm. Diện tích xung quanh của hình lăng trụ đứng đó là:\par
\hspace*{1.5em}\textbf{A.} 1200 cm$^2$\quad\quad \textbf{B.} 74 cm$^2$\quad\quad \textbf{C.} 2400 cm$^2$\quad\quad \textbf{D.} 600 cm$^2$\par\vspace{2pt}
\subsection*{\color{accent} Phần III. Mức độ Khó (Vận dụng \& Vận dụng cao -- Câu 41 đến Câu 50)}
\textbf{\color{primary} Câu 41.} Một bể cá dạng hình hộp chữ nhật có chiều dài 8 dm, chiều rộng 5 dm và chiều cao 6 dm. Hiện bể đang chứa nước đến 3/4 chiều cao của bể. Thể tích nước trong bể là bao nhiêu lít (1 dm$^3$ = 1 lít)?\par
\hspace*{1.5em}\textbf{A.} 240 lít\quad\quad \textbf{B.} 180 lít\quad\quad \textbf{C.} 120 lít\quad\quad \textbf{D.} 60 lít\par\vspace{2pt}
\textbf{\color{primary} Câu 42.} Bài toán vận dụng (Chương X – Một số hình khối trong thực tiễn): Một căn phòng hình hộp chữ nhật dài 7 m, rộng 4 m, cao 4 m. Người ta muốn quét sơn trần nhà và bốn bức tường phía trong phòng, biết tổng diện tích các cửa là 10 m$^2$. Diện tích cần quét sơn là. Kết quả đúng là:\par
\hspace*{1.5em}\textbf{A.} 134 m$^2$\quad\quad \textbf{B.} 102 m$^2$\quad\quad \textbf{C.} 78 m$^2$\quad\quad \textbf{D.} 106 m$^2$\par\vspace{2pt}
\textbf{\color{primary} Câu 43.} Vận dụng kiến thức nâng cao trong phạm vi Chương X để giải bài toán: Một bể cá dạng hình hộp chữ nhật có chiều dài 10 dm, chiều rộng 5 dm và chiều cao 6 dm. Hiện bể đang chứa nước đến 3/4 chiều cao của bể. Thể tích nước trong bể là bao nhiêu lít (1 dm$^3$ = 1 lít):\par
\hspace*{1.5em}\textbf{A.} 75 lít\quad\quad \textbf{B.} 150 lít\quad\quad \textbf{C.} 225 lít\quad\quad \textbf{D.} 300 lít\par\vspace{2pt}
\textbf{\color{primary} Câu 44.} Phân tích và giải bài toán tổng hợp sau của Chương X: Một căn phòng hình hộp chữ nhật dài 6 m, rộng 4 m, cao 4 m. Người ta muốn quét sơn trần nhà và bốn bức tường phía trong phòng, biết tổng diện tích các cửa là 10 m$^2$. Diện tích cần quét sơn là:\par
\hspace*{1.5em}\textbf{A.} 70 m$^2$\quad\quad \textbf{B.} 94 m$^2$\quad\quad \textbf{C.} 118 m$^2$\quad\quad \textbf{D.} 86 m$^2$\par\vspace{2pt}
\textbf{\color{primary} Câu 45.} Thử thách tư duy trong phạm vi kiến thức Chương X: Một bể cá dạng hình hộp chữ nhật có chiều dài 12 dm, chiều rộng 5 dm và chiều cao 6 dm. Hiện bể đang chứa nước đến 3/4 chiều cao của bể. Thể tích nước trong bể là bao nhiêu lít (1 dm$^3$ = 1 lít). Đáp số chính xác là:\par
\hspace*{1.5em}\textbf{A.} 360 lít\quad\quad \textbf{B.} 90 lít\quad\quad \textbf{C.} 270 lít\quad\quad \textbf{D.} 180 lít\par\vspace{2pt}
\textbf{\color{primary} Câu 46.} Sử dụng phương pháp giải toán của Chương X để tìm đáp án cho bài toán: Một căn phòng hình hộp chữ nhật dài 8 m, rộng 4 m, cao 4 m. Người ta muốn quét sơn trần nhà và bốn bức tường phía trong phòng, biết tổng diện tích các cửa là 10 m$^2$. Diện tích cần quét sơn là:\par
\hspace*{1.5em}\textbf{A.} 86 m$^2$\quad\quad \textbf{B.} 118 m$^2$ (1)\quad\quad \textbf{C.} 118 m$^2$\quad\quad \textbf{D.} 150 m$^2$\par\vspace{2pt}
\textbf{\color{primary} Câu 47.} Xét bài toán vận dụng thực tiễn / nâng cao thuộc Chương X: Một bể cá dạng hình hộp chữ nhật có chiều dài 14 dm, chiều rộng 5 dm và chiều cao 6 dm. Hiện bể đang chứa nước đến 3/4 chiều cao của bể. Thể tích nước trong bể là bao nhiêu lít (1 dm$^3$ = 1 lít):\par
\hspace*{1.5em}\textbf{A.} 420 lít\quad\quad \textbf{B.} 105 lít\quad\quad \textbf{C.} 315 lít\quad\quad \textbf{D.} 210 lít\par\vspace{2pt}
\textbf{\color{primary} Câu 48.} Tìm kết quả cuối cùng của bài toán phân loại học sinh khá giỏi (Chương X): Một căn phòng hình hộp chữ nhật dài 7 m, rộng 4 m, cao 4 m. Người ta muốn quét sơn trần nhà và bốn bức tường phía trong phòng, biết tổng diện tích các cửa là 10 m$^2$. Diện tích cần quét sơn là:\par
\hspace*{1.5em}\textbf{A.} 134 m$^2$\quad\quad \textbf{B.} 102 m$^2$\quad\quad \textbf{C.} 106 m$^2$\quad\quad \textbf{D.} 78 m$^2$\par\vspace{2pt}
\textbf{\color{primary} Câu 49.} Kết hợp các tính chất đã học trong Chương X, hãy giải quyết bài toán: Một bể cá dạng hình hộp chữ nhật có chiều dài 16 dm, chiều rộng 5 dm và chiều cao 6 dm. Hiện bể đang chứa nước đến 3/4 chiều cao của bể. Thể tích nước trong bể là bao nhiêu lít (1 dm$^3$ = 1 lít):\par
\hspace*{1.5em}\textbf{A.} 120 lít\quad\quad \textbf{B.} 360 lít\quad\quad \textbf{C.} 480 lít\quad\quad \textbf{D.} 240 lít\par\vspace{2pt}
\textbf{\color{primary} Câu 50.} Câu hỏi vận dụng cao cuối Chương X: Một căn phòng hình hộp chữ nhật dài 6 m, rộng 4 m, cao 4 m. Người ta muốn quét sơn trần nhà và bốn bức tường phía trong phòng, biết tổng diện tích các cửa là 10 m$^2$. Diện tích cần quét sơn là. Phương án đúng là:\par
\hspace*{1.5em}\textbf{A.} 94 m$^2$\quad\quad \textbf{B.} 70 m$^2$\quad\quad \textbf{C.} 86 m$^2$\quad\quad \textbf{D.} 118 m$^2$\par\vspace{2pt}
\subsection*{\color{primary} Bảng đáp án Chương X: Một số hình khối trong thực tiễn}
\begin{center}
\renewcommand{\arraystretch}{1.25}
\begin{tabularx}{\textwidth}{*{10}{>{\centering\arraybackslash}X}}
\toprule
\textbf{\color{primary} Cột 1} & \textbf{\color{primary} Cột 2} & \textbf{\color{primary} Cột 3} & \textbf{\color{primary} Cột 4} & \textbf{\color{primary} Cột 5} & \textbf{\color{primary} Cột 6} & \textbf{\color{primary} Cột 7} & \textbf{\color{primary} Cột 8} & \textbf{\color{primary} Cột 9} & \textbf{\color{primary} Cột 10} \\
\midrule
\textbf{1.A} & \textbf{2.A} & \textbf{3.C} & \textbf{4.C} & \textbf{5.D} & \textbf{6.B} & \textbf{7.D} & \textbf{8.B} & \textbf{9.C} & \textbf{10.C} \\
\rowcolor{rowtint} \textbf{11.B} & \textbf{12.B} & \textbf{13.A} & \textbf{14.C} & \textbf{15.B} & \textbf{16.B} & \textbf{17.B} & \textbf{18.D} & \textbf{19.B} & \textbf{20.D} \\
\textbf{21.D} & \textbf{22.D} & \textbf{23.B} & \textbf{24.A} & \textbf{25.A} & \textbf{26.C} & \textbf{27.C} & \textbf{28.D} & \textbf{29.C} & \textbf{30.A} \\
\rowcolor{rowtint} \textbf{31.D} & \textbf{32.B} & \textbf{33.D} & \textbf{34.C} & \textbf{35.C} & \textbf{36.B} & \textbf{37.D} & \textbf{38.C} & \textbf{39.A} & \textbf{40.A} \\
\textbf{41.B} & \textbf{42.D} & \textbf{43.C} & \textbf{44.B} & \textbf{45.C} & \textbf{46.C} & \textbf{47.C} & \textbf{48.C} & \textbf{49.B} & \textbf{50.A} \\
\bottomrule
\end{tabularx}
\end{center}\vspace{10pt}
\end{document}
